% BEGIN INCLUDED SGA4_English_current_buildable_source_20260906/source/05bis/drafts/SGA4_VBIS_LINES_59_205_English.texfrag
\section*{引言}

\begin{enumerate}
\renewcommand{\labelenumi}{\theenumi.}
\item 设\pageoriginale{} $X$ 为拓扑空间, $\Uu=(U_i)_{i\in I}$ 为 $X$ 的覆盖,
假定它或为开覆盖, 或为局部有限的闭覆盖.
若 $\mathfrak{F}$ 为 $X$ 上的阿贝尔层, 则由 $\Uu$ 定义的 Leray 谱序列
\begin{equation*}
\Hv{p}(\Uu,\mathcal{H}^q(\mathfrak{F}))\Rightarrow H^{\ast}
(X,\mathfrak{F})\tag{1.1}
\end{equation*}\etiquette[1.1]{Vbis.eq1.1}
[\cite{sga4.exp.Vbis.bib.5} II (5.2.4) 与 (5.4.1)] 可描述如下.

覆盖 $\Uu$ 定义一个对 $\mathfrak{F}$ 函子性的“\v{C}ech”分解
$\mathcal{C}^{\ast}(\Uu,\mathfrak{F})$ (同上 (\hyperref[Vbis.subsubsec5.2.1]{5.2.1})).
另一方面, 对每个 $\mathfrak{F}$, 存在对 $\mathfrak{F}$ 函子性的“典范松弛”分解
$C^{\ast}(\mathfrak{F})$ (同上 (\hyperref[Vbis.subsec4.3]{4.3})).
采用此记号, 当 $\Uu$ 为开覆盖时, 谱序列 (\hyperref[Vbis.subsec1.1]{1.1}) 来自双复形
\begin{equation*}
\Gamma(X,\mathcal{C}^{\ast}(\Uu,C^{\ast}(\mathfrak{F}))){\quoi}.
\tag{1.2}
\end{equation*}\etiquette[1.2]{Vbis.eq1.2}
当 $\Uu$ 为局部有限的闭覆盖时, 则考虑双复形
\begin{equation*}
\Gamma(X,C^{\ast}(\mathcal{C}^{\ast}(\Uu,\mathfrak{F}))){\quoi}.
\tag{1.3}
\end{equation*}\etiquette[1.3]{Vbis.eq1.3}

\item 试为这些双复形寻找统一的描述.
令 $X_0$ 为各 $U_i$ 的拓扑不交并, 并对 $n\geq0$,
令 $X_n$ 为 $X_0$ 在 $X$ 上与自身的 $(n+1)$ 重纤维积:
\begin{equation*}
X_n=\mathop{\coprod}_{i_0,\ldots,i_n\in I}U_{i_0}\cap\cdots\cap U_{i_n}
=\mathop{\coprod}_{\sigma\in\Hom([n],I)}
\mathop{\bigcap}_{i=0}^{n}U_{\sigma(i)}{\quoi}.\tag{2.1}
\end{equation*}\etiquette[2.1]{Vbis.eq2.1}
各 $X_n$ 形成拓扑空间的单纯系统, 若 $j_n$ 表示从 $X_n$ 到 $X$ 的投影, 则有
\begin{equation*}
\mathcal{C}^n(\Uu,\mathfrak{F})=j_{n\ast}j_n^{\ast}(\mathfrak{F}){\quoi}.
\tag{2.2}
\end{equation*}\etiquette[2.2]{Vbis.eq2.2}

注意, 典范松弛分解的形成与\pageoriginale{}限制到开子集以及闭浸入的正像交换. 因而:
\begin{enumerate}
\renewcommand{\theenumi}{\alph{enumi}}
\renewcommand{\labelenumi}{\theenumi)}
\item 若 $\Uu$ 为开覆盖,
\begin{equation*}
\mathcal{C}^q(\Uu,C^p(\mathfrak{F}))=j_{q\ast}j_q^{\ast}
C^p(\mathfrak{F})=j_{q\ast}C^p(j_q^{\ast}\mathfrak{F});\tag{2.3}
\end{equation*}\etiquette[2.3]{Vbis.eq2.3}
\item 若 $\Uu$ 为局部有限的闭覆盖,
\begin{equation*}
C^p(\mathcal{C}^q(\Uu,\mathfrak{F}))=C^p(j_{q\ast}j_q^{\ast}
\mathfrak{F})=j_{q\ast}C^p(j_q^{\ast}\mathfrak{F}){\quoi}.\tag{2.4}
\end{equation*}\etiquette[2.4]{Vbis.eq2.4}
\end{enumerate}
因此, 为统一描述 \eqhyperref[Vbis.eq1.2]{1.2} 与 \eqhyperref[Vbis.eq1.3]{1.3},
只须对每个 $q$ 取 $X_q$ 上 $j_q^{\ast}(\mathfrak{F})$ 的典范松弛分解,
再对该分解应用函子 $j_{q\ast}$.

\item 上述描述对于 $X$ 上任意拓扑空间单纯系统仍有意义:
\hypertarget{sga4.exp.Vbis.equation.3.1}{}\begin{gather*}
\Delta^0\longrightarrow\Top/X\tag{3.1}\\
[n]\longmapsto X_n.\notag
\end{gather*}
这样的系统不一定形如 \eqhyperref[Vbis.eq2.1]{2.1}.
然而由 $\mathfrak{F}$ 余增广的双复形
\begin{equation*}
\mathfrak{F}\longrightarrow
(j_{q\ast}C^p(j_q^{\ast}(\mathfrak{F})))_{p,q}\tag{3.2}
\end{equation*}\etiquette[3.2]{Vbis.eq3.2}
一般并不定义 $\mathfrak{F}$ 的分解.

本文致力于寻找使 \eqhyperref[Vbis.eq3.2]{3.2} 定义 $\mathfrak{F}$ 的分解的充分条件.
在此情形中, 谱序列 (\hyperref[Vbis.subsec1.1]{1.1}) 推广为谱序列
\begin{equation*}
\Hv{p}(H^q(X_p,j_p^{\ast}(\mathfrak{F})))\Rightarrow
H^{p+q}(X,\mathfrak{F}),\tag{3.3}
\end{equation*}\etiquette[3.3]{Vbis.eq3.3}
称为“下降谱序列”.

在\pageoriginale“常系数”情形中, Segal (参见 \cite{sga4.exp.Vbis.bib.8})
曾用其他方法, 对例如 $K$-理论等其他“上同调理论”得到类似谱序列.
Segal 在 CW-复形范畴中工作: 他使用“几何实现”函子,
它将拓扑空间的半单纯复形对应到一个新的拓扑空间;
这个新空间可与单纯拓扑斯的相伴拓扑斯比较
[参见 \hyperref[Vbis.prop1.2.12]{(1.2.12)}].

\item 在 \hyperref[Vbis.sec5]{第 5 节}, 我们用奇点消解说明所得判别准则的应用:
对 $\CC$ 上的每个解析空间 $X$, 构造 $X$ 上非奇异解析空间的单纯系统
$$
[n]\longmapsto X_n,
$$
使 \eqhyperref[Vbis.eq3.2]{3.2} 定义 $\mathfrak{F}$ 的分解.
以常值层 $\CC$ 作为 $\mathfrak{F}$, 特别得到谱序列
\begin{equation*}
H^q(X_p,\CC)\Longrightarrow H^{p+q}(X,\CC)\tag{4.1}
\end{equation*}\etiquette[4.1]{Vbis.eq4.1}
它用非奇异解析空间的复上同调表示 $X$ 的复上同调.
此外, 若 $X$ 射影, 可假定所有 $X_p$ 均射影;
这是得到 $X$ 的某种“Hodge 理论”的关键因素 (参见 \cite{sga4.exp.Vbis.bib.2}).

\item 将层 $\mathfrak{F}$ 换为下有界层复形时, 上述构造逐字推广.
它们导出导出范畴中的“局部化”技术.

已知, 对 \eqhyperref[Vbis.eq2.1]{2.1} 给出的 $X_0$, 箭头
$$
j_0^{\ast}:D^b(X)\longrightarrow D^b(X_0)
$$
一般并不忠实. 我们将证明, 涉及各 $X_n$ 的、比 $j_0^{\ast}(K^{\bigdot})$\pageoriginale{}
(其中 $K^{\bigdot}\in D^+(X)$) 更精细的资料, 有时能使我们重构复形 $K^{\bigdot}$.

所得命题将在第 XVII 篇附录中使用, 将函子 $Rf_!$
(其中 $f$ 为概形之间有限型的分离态射) 的定义扩张到不假定 $f$ 可紧化的情形.

为作此应用, 不能只考虑拓扑空间; 这里以概形的平展景代之.
此外, 为完成证明, 必须同时考虑空间的单纯系统及多重单纯系统.
这解释了、说明了, 或为我们起步时所取的极高的一般性程度提供了辩解.
\end{enumerate}


% END INCLUDED SGA4_English_current_buildable_source_20260906/source/05bis/drafts/SGA4_VBIS_LINES_59_205_English.texfrag


% BEGIN INCLUDED SGA4_English_current_buildable_source_20260906/source/05bis/drafts/SGA4_VBIS_LINES_206_516_English.texfrag
\section{预备知识}\etiquette[1]{Vbis.sec1}

\subsection{记号}\etiquette[1.1]{Vbis.subsec1.1}

\subsubsection{}\etiquette[1.1.1]{Vbis.subsubsec1.1.1}
下文始终设 $\Uu$ 为使 $\ZZ \in \Uu$ 的宇宙:
\emph{所考虑的一切拓扑斯均为 $\Uu\text{-拓扑斯}$.}

设 $T$ 与 $T'$ 为两个拓扑斯: 态射 $\varphi : T \to T'$ 由一对函子
$\varphi_{\ast} : T \to T'$ 和 $\varphi^{\ast} : T'\to T$ 组成, 配备伴随关系
$$\Hom_{T'}(.,\varphi_{\ast}.) \isomto{\Hom}_{T}(\varphi^{\ast} . ,
.),$$
使得 $\varphi^{\ast}$ 左正合 (即保持有限投射极限).

设 $(T, \mathcal{O}_{T})$ 与 $(T', \mathcal{O}_{T'})$ 为两个环化拓扑斯:
从 $(T, \mathcal{O}_{T})$ 到 $(T', \mathcal{O}_{T'})$ 的态射为二元组 $(\varphi, \theta)$,
其中 $\varphi : T \to T'$ 为拓扑斯态射,
$\theta : \mathcal{O}_{T'} \to \varphi_{\ast} (\mathcal{O}_T)$ 为环态射.

\subsubsection{}\etiquette[1.1.2]{Vbis.subsubsec1.1.2}
我们将不断使用 [\SGA 1 VI] 中介绍的纤维范畴语言; 读者也可参阅 \cite{sga4.exp.Vbis.bib.4}.
这里只固定一些记号: 若 $E \to B$ 为纤维 (分别为余纤维) 函子,
则对 $B$ 中的态射 $m : i \to j$, 用 $m^{\ast} : E_{j} \to E_i$
(分别用 $m_{\ast} : E_i \to E_j$) 表示相应的“逆像”
(分别为\pageoriginale“正像”) 函子;
每个这样的函子均定义到唯一的函子性同构.
若 $\varphi : E \to E'$ 为 $B$-函子, 则对 $B$ 的每个对象 $i$,
用 $\varphi_i : E_i \to E'_i$ 表示 $\varphi$ 在 $E_i$ 上的限制.

\subsubsection{}\etiquette[1.1.3]{Vbis.subsubsec1.1.3}

最后, 用 $\Delta$ 表示如下范畴:
$\Delta$ 的对象为有序集 $[n] = \{ 0, 1, \ldots, n \}$,
$\Delta$ 的态射为所有弱递增映射.
$\Delta^{+}$ (分别为 $\Delta^{-}$) 表示其对象与 $\Delta$ 相同,
而箭头为 $\Delta$ 的单态射 (分别为满态射) 的范畴.
我们将按需要自由使用与 $\Delta$ 有关的经典记号 [参见 \cite{sga4.exp.Vbis.bib.3} II 2].

\subsection{D-拓扑斯}\etiquette[1.2]{Vbis.subsec1.2}

\begin{enonce*}{定义 1.2.1}\etiquette[1.2.1]{Vbis.def1.2.1}
设 $D$ 为范畴, $E \to D$ 为纤维且余纤维函子.
若下述条件成立, 就称 $E$ 在 $D$ 上为拓扑斯双纤维化的,
或称 $E$ 为 $D$-拓扑斯:
\begin{enumerate}
\renewcommand{\theenumi}{\alph{enumi}}
\renewcommand{\labelenumi}{(\theenumi)}
\item 对 $D$ 的每个对象 $i$, 纤维范畴 $E_i$ 为拓扑斯.

\item 对 $D$ 的每个态射 $m : i \to j$, 存在拓扑斯态射 $f : E_j \to E_i$,
使 $f_{\ast} = m^{\ast}$ 且 $f^{\ast} = m_{\ast}$.
\end{enumerate}
\end{enonce*}

\begin{enonce*}[remark]{注}
考虑到 (a), 条件 (b) 也可等价表述为函子 $m_{\ast}$ 左正合
\nde{这里须注意: 在纤维范畴中, $m_\star$ 为 $m^\star$ 的左伴随,
因而其右正合性是重言的; 参见 [\SGA 1 VI.10].
这解释了公式 $f_{\ast} = m^{\ast}$ 和 $f^{\ast}=m_{\ast}$.}.
\end{enonce*}

当 $D = \Delta$ (分别为 $\Delta \times \Delta$) 时,
用\emph{单纯拓扑斯} (分别用\emph{双重单纯拓扑斯})
表示 $\Delta$-拓扑斯 (分别为 $\Delta \times \Delta$-拓扑斯).

实际中, 我们将通过下述考虑遇到 $D$-拓扑斯.

\begin{enonce*}{定义 1.2.2}\etiquette[1.2.2]{Vbis.def1.2.2}
设\pageoriginale{} $\mathcal{E}$ 为 $D$ 上的纤维且余纤维范畴.
若 $\mathcal{E}^{\circ}$ 为 $D^{\circ}$-拓扑斯,
就称 $\mathcal{E}$ \emph{在 $D$ 上为拓扑斯对偶双纤维化的}.
\end{enonce*}

\hypertarget{sga4.exp.Vbis.remark.1.2.3}{}\begin{enonce*}[remark]{注 1.2.3}
明确地说, $\mathcal{E}$ 在 $D$ 上为拓扑斯对偶双纤维化的,
当且仅当下述条件成立:
\begin{enumerate}
\renewcommand{\theenumi}{\alph{enumi}}
\renewcommand{\labelenumi}{(\theenumi)}
\item 对 $D$ 的每个对象 $i$, 纤维范畴 $\mathcal{E}_i$ 的反范畴
\sisi{}{${\mathcal{E}_i}^\circ$} 为拓扑斯.

\item 对 $D$ 的每个态射 $m : i \to j$, 函子 $m^{\ast}$ 左正合
\nde{参见前注. 换言之, 函子族
$((m_{\ast})^\circ,(m^{\ast})^\circ)$ 定义一个相容的拓扑斯态射系统
$\mathcal{E}_i^\circ\rightarrow \mathcal{E}_j^\circ$.}.
\end{enumerate}
\end{enonce*}

\setcounter{subsection}{2}
\setcounter{subsubsection}{3}
\subsubsection{}\etiquette[1.2.4]{Vbis.subsubsec1.2.4}
读者可在 \hyperref[Vbis.sec4]{第 4 节} 找到拓扑斯对偶双纤维化范畴的例子.

\subsubsection{}\etiquette[1.2.5]{Vbis.subsubsec1.2.5}
设 $\mathcal{E} \to B$ 为 $B$ 上拓扑斯对偶双纤维化的范畴,
$X : D^{\circ} \to B$ 为函子.
请读者验证 $(\fprod{D^{\circ}}{\mathcal{E}}{B})^{\circ}$ 是 $D$-拓扑斯,
我们将其记为 $\bar{X}$.

\subsubsection{}\etiquette[1.2.6]{Vbis.subsubsec1.2.6}
通常称函子 $X : D^{\circ} \to B$ 为 $B$ \emph{的} $D$-\emph{对象},
用 $X_i$ 表示 $D$ 的对象 $i$ 在此函子下的像;
$B$ 的 $D$-对象形成范畴, 记为 $D^{\circ}B$.
若 $S$ 为 $B$ 的对象, 则 $B/S$ 的 $D$-对象称为\emph{向 $S$ 增广的 $B$ 的 $D$-对象}.
给出向 $S$ 增广的 $B$ 的 $D$-对象, 显然等价于给出函子性态射 $X \to C^{D}_S$,
其中 $X$ 为 $B$ 的 $D$-对象, $C^D_S$ 为 $B$ 中取值为 $S$ 的常值 $D$-对象.

当 $D = \Delta$ 时, 称为\emph{单纯对象} (分别为\emph{增广单纯对象}).

我们还将使用\emph{双重单纯对象} (取 $D = \Delta \times \Delta$).

\subsubsection{}\etiquette[1.2.7]{Vbis.subsubsec1.2.7}
现假定\pageoriginale{}范畴 $B$ 有有限纤维积.

设 $f : R \to S$ 为 $B$ 中的箭头: 从 $\Delta^{\circ} \times (B / S)^{\circ}$ 到 $\Ens$ 的双函子
$$
([n], X) \rightsquigarrow \Hom_{\Ens} ([n], \Hom_S (X, R))
$$
定义函子
$$
[n] \rightsquigarrow X_{[n]} = X_n
$$
其中取 $X_n$ 为函子
$$
Z \rightsquigarrow \Hom_{\Ens} ([n], \Hom_S (Z, R))
$$
的一个表示对象
$(X_n \simeq \displaystyle\mathop{\overbrace{R \times_S R
    \times \ldots \times_S R}}^{n+1 \text{ 次 }})$.

用 $[R |_{f} S]$ 或 $[R | S]$ 表示这样构造的向 $S$ 增广的半单纯对象.

最后, 若 $X$ 与 $X'$ 为 $B$ (或 $B/S$) 的两个半单纯对象,
并且 $$u : X \to X'$$ 为函子性态射,
则出于技术原因, 我们引入在范畴 $\Delta^{\circ}B$ 中计算的对象 $$[X |_{u} X']$$.
此对象可解释为 $B$ 的双重单纯对象
\nde{明确地说, 按指标的通常次序, 有
$$[[X |_u X']]_{(n,m)}=\displaystyle\mathop{\overbrace{X_m \times_{X'_m} X_m
    \times \ldots \times_{X'_m} X_m}}^{n+1 \text{ 次 }}$$
其函子性由各因子的函子性导出.}, 此时记为 $[[X |_u X']]$:
事实上, 有典范范畴同构
$$
\Delta^{\circ} (\Delta^{\circ} B) \xrightarrow{\sim} (\Delta
\times \Delta)^{\circ} B{\quoi}.
$$

现在回到 $D$-拓扑斯的一般概念.

\subsubsection{}\etiquette[1.2.8]{Vbis.subsubsec1.2.8}
设 $E$ 为 $D$-拓扑斯: 用 $\Ysheaf{\Gamma}(E)$ 表示范畴 $\Hom_D (D, E)$.
设 $f : D'\to D$ 为函子: 范畴 $\fprod{D'}{E}{D}$ 是 $D'$-拓扑斯,
且与 $f$ 复合得到函子
$$
f^{\ast} : \Ysheaf{\Gamma} (E) \longrightarrow \Ysheaf{\Gamma} (\fprod{D'}{E}{D}){\quoi}.
$$

在\pageoriginale{} $D'$ 仅由 $D$ 的一个对象 $i$ 组成
(以 $i$ 的恒等态射为唯一态射)、且 $f$ 为记作 $e_i$ 的典范包含的情形中,
可取 $\fprod{D'}{E}{D}$ 为纤维范畴 $E_i$,
此时 $e^{\ast}_i$ 等同于“\emph{在 $i$ 处求值}”函子.

\begin{enonce*}{命题 1.2.9}\etiquette[1.2.9]{Vbis.prop1.2.9}
若 $D'$ 为 $\Uu$-小范畴, 则函子 $f^{\ast}$ 有右伴随和左伴随
(分别记为 $f_{\ast}$ 和 $f_{!}$).
\end{enonce*}

这来自 Kan 引理 [III 1.1] 的一个小推广, 我们将反复使用它.

\begin{enonce*}{引理 1.2.10}\etiquette[1.2.10]{Vbis.lem1.2.10}
设 $I$、$J$ 和 $A$ 为同一范畴 $B$ 上的三个范畴.
假定 $I$ 为 $\Uu$-小, 且 $A$ 在 $B$ 上纤维化且余纤维化.
给定 $B$-函子 $f : I \to J$, 用 $f^{\ast}$ 表示与 $f$ 复合所定义的函子
$$
\Hom_B (J, A) \longrightarrow \Hom_B (I, A)
$$
若 $A$ 在 $B$ 上的每个纤维中, $\Uu$-归纳极限 (分别为投射极限) 存在,
则 $f^{\ast}$ 有左伴随 (分别为右伴随).
\end{enonce*}

\begin{enonce*}[remark]{证明}
只证明左伴随的存在性, 引理括号内的对应部分由对偶性得出.
使用下述事实, 其验证留待读者完成.
\end{enonce*}

\begin{enonce*}{引理 1.2.10.1}\etiquette[1.2.10.1]{Vbis.lem1.2.10.1}
设 $A$ 为范畴 $B$ 上的双纤维范畴, $b$ 为 $B$ 的对象,
$(F_{\lambda})_{\lambda \in \Lambda}$ 为从范畴 $\Lambda$ 到纤维 $A_b$ 的函子.
对 $A_{b'}$ 中的每个 $G$, 以及 $B$ 中的每个 $u : b'\to b$
(分别为 $u : b \to b'$), 只要左端有定义, 就存在双射
$$
\Hom_u(G, \varprojlim F_{\lambda}) \xrightarrow{\sim}
\varprojlim \Hom_u (G, F_{\lambda})
$$
(分别为 $\Hom_u (\varinjlim F_{\lambda}, G)
\xrightarrow{\sim}\varprojlim \Hom_u (F_{\lambda}, G)$).
\end{enonce*}

\sisi{$+$}{} 明确这点后,\pageoriginale{}令 $I \amalg_f J$ 为如下定义的 $B$ 上范畴:
$$
\ob (I \textstyle{\amalg}_f J) = \ob (I) \amalg \ob (J)
$$
\begin{equation*}
\Hom(x, y)=
\begin{cases}
\Hom_I(x, y) & \text{ 若 } x, y \in \ob (I)\\
\Hom_J(x, y) & \text{ 若 } x, y \in \ob (J)\\
\Hom_J(f(x), y) & \text{ 若 } x \in \ob(I) \text{ 且 } y \in \ob(J)\\
∅ & \text{ 若 } x \in \ob(J) \text{ 且 } y \in
\ob(I){\quoi}.\\
\end{cases}
\end{equation*}

\sisi{$+$}{} 范畴 $\Hom_B(I \amalg_f J, A)$ 等价于如下三元组组成的范畴:
从 $I$ 到 $A$ 的 B-函子 $F$, 从 $J$ 到 $A$ 的 $B$-函子 $G$,
以及从 $F$ 到 $G \circ f = f(G)$ 的 $B$-函子态射 $\varphi$.

\sisi{$+$}{} 对 $J$ 的每个对象 $j$, 用 $I /j$ 表示 $I$ 中“由 $f$ 放在 $j$ 上”的对象所成范畴:
$I /j$ 的对象为二元组 $(i, \alpha)$, 其中 $i$ 为 $I$ 的对象,
$\alpha : i \to j$ 为 $I \amalg_f J/j$ 中的箭头;
$I/j$ 的态射为 $I\amalg_{f}J/j$ 中的态射.
若 $p_I$ 和 $p_J$ 为 $I$ 和 $J$ 到 $B$ 的投影, 且 $(i, \alpha)$ 为 $I/j$ 的对象,
用 $\alpha_{\ast} : A_{p_{I}(i)} \to A_{p_{J}(j)}$ 表示函子 $p_{J}(\alpha)_{\ast}$.

\sisi{$+$}{} 给出从 $I \amalg_f J$ 到 $A$ 的 $B$-函子,
等价于给出 $F \in \Hom_B (I, A)$、$G \in \Hom_B (J, A)$
以及对 $j$ 函子性的态射
$$
\psi_j : \varinjlim_{(i, \alpha)\in I/j} \alpha_{\ast}(F(i))
\longrightarrow G(j){\quoi}.
$$
(左端对 $j$ 的函子性来自 \hyperref[Vbis.lem1.2.10.1]{1.2.10.1}.)
因此 $f^{\ast}$ 的左伴随由下式给出
$$
f_{!}(F) (j) = \varinjlim_{(i, \alpha)\in I/j}
\alpha_{\ast}(F(i)){\quoi}.
$$

\begin{enonce*}{推论 1.2.11}\etiquette[1.2.11]{Vbis.cor1.2.11}
设\pageoriginale{} $E$ 为 $D$-拓扑斯, $i$ 为 $D$ 的对象;
函子 $e^{\ast}_i$ 有左伴随 (分别为右伴随), 由下式给出
$$
e_{i!}(a)(j) = \coprod_{\alpha \in \Hom(i, j)} \alpha_{\ast}(a)
\left( \emph{\resp} e_{i\ast} (a) (j) = \prod_{\alpha \in \Hom(j,
  i)} \alpha^{\ast}(a)\right)
$$
其中 a 为 $E$ 中 $i$ 上的对象.
\end{enonce*}

\begin{enonce*}{命题 1.2.12}\etiquette[1.2.12]{Vbis.prop1.2.12}
设 $D$ 为 $\Uu$-小范畴, $E$ 为 $D$-拓扑斯.
则范畴 $\Ysheaf{\Gamma}(E)$ 为 $\Uu$-拓扑斯.
\end{enonce*}

\sisi{$+$}{} 利用 \hyperref[Vbis.lem1.2.10.1]{1.2.10.1},
可“逐纤维”验证范畴 $\Ysheaf{\Gamma}(E)$ 具有下述性质:
\begin{enumerate}
\renewcommand{\theenumi}{\alph{enumi}}
\renewcommand{\labelenumi}{\theenumi)}
\item 有限投射极限可表.
\item 以 $\Uu$ 的元素为指标的余积可表. 它们不交且普遍.
\item 等价关系普遍有效.
\end{enumerate}

\sisi{$+$}{} 还须证明 $\Ysheaf{\Gamma}(E)$ 有以 $\Uu$ 的某个元素为指标的生成系.
但若对 $D$ 的每个对象 $i$, $(G_{i\lambda})_{\lambda \in \Lambda_{i}}$
为 $E_i$ 的生成系 (其中 $\Lambda_i$ 为 $\Uu$-小集),
则族 $(e_{i!}(G_{i\lambda}))_{i, \lambda}$ 为 $\Ysheaf{\Gamma}(E)$ 的生成系.

\setcounter{subsection}{2}
\setcounter{subsubsection}{12}

% END INCLUDED SGA4_English_current_buildable_source_20260906/source/05bis/drafts/SGA4_VBIS_LINES_206_516_English.texfrag


% BEGIN INCLUDED SGA4_English_current_buildable_source_20260906/source/05bis/drafts/SGA4_VBIS_LINES_517_829_English.texfrag
\subsubsection{}\etiquette[1.2.13]{Vbis.subsubsec1.2.13}
现在引入 $D$-拓扑斯之间的态射概念.
首先说明, 若 $F$ 和 $F'$ 为同一范畴 $B$ 上的两个范畴,
$T:F\to F'$ 为 $B$-函子, 则 $T$ 的 \emph{$B$-左伴随}
意指 $T$ 的一个左伴随函子 $S:F'\to F$,
使典范态射 $1\to T\circ S$ 与 $S\circ T\to1$ 均为函子的 $B$-态射.
在此条件下, 立即可验 $T$ 笛卡尔且 $S$ 余笛卡尔.

\begin{enonce*}{定义 1.2.14}\etiquette[1.2.14]{Vbis.def1.2.14}
设\pageoriginale{} $E$ 和 $E'$ 为两个 $D$-拓扑斯.
从 $E$ 到 $E'$ 的态射为一对 $D$-函子 $\Phi_{\ast}:E\to E'$ 与 $\Phi^{\ast}:E'\to E$,
配备 $D$-伴随关系
$$
\Hom(\Phi^{\ast}(-),-)\xrightarrow[\xi]{\sim}
\Hom(-,\Phi_{\ast}(-)),
$$
使对 $D$ 的每个对象 $i$, 二元组 $(\Phi_{\ast i},\Phi_i^{\ast})$
配备由 $\xi$ 诱导的伴随关系后, 为从 $E_i$ 到 $E'_i$ 的拓扑斯态射.
\end{enonce*}

\begin{enonce*}{命题 1.2.15}\etiquette[1.2.15]{Vbis.prop1.2.15}
设 $E$ 与 $E'$ 为两个 $D$-拓扑斯, $(\Phi_{\ast},\Phi^{\ast}):E\to E'$ 为态射.
若 $D$ 为 $\Uu$-小范畴, 则二元组
$(\Gamma(\Phi_{\ast}),\Gamma(\Phi^{\ast})):\Ysheaf{\Gamma}(E)\to\Ysheaf{\Gamma}(E')$
为拓扑斯态射.
\end{enonce*}

这立即来自前面的定义.

下述引理提供构造 $D$-拓扑斯态射的一种方法, 其证明留待读者完成.

\begin{enonce*}{引理 1.2.16}\etiquette[1.2.16]{Vbis.lem1.2.16}
设 $E$ 与 $E'$ 为同一范畴 $D$ 上的两个双纤维范畴,
$\Phi$ 为从 $E$ 到 $E'$ 的\emph{笛卡尔} $D$-函子,
使对 $D$ 的每个对象 $i$, $\Phi_i:E_i\to E'_i$ 有左伴随.
则对每个 $i$ 选取 $\Phi_i$ 的左伴随, 便典范确定一个为 $\Phi$ 的 $D$-左伴随的
$D$-函子 $\Psi:E'\to E$.
\end{enonce*}

\begin{enonce*}[remark]{评注 1.2.17}\etiquette[1.2.17]{Vbis.sch1.2.17}
在 (\hyperref[Vbis.lem1.2.16]{1.2.16}) 的条件下, 假定 $E$ 和 $E'$ 为两个 $D$-拓扑斯,
且对 $D$ 的每个对象 $i$, $\Phi_i$ 的每个左伴随都左正合.
则若 $\Psi:E'\to E$ 为 $\Phi$ 的 $D$-左伴随,
二元组 $(\Phi,\Psi):E\to E'$ 即为 $D$-拓扑斯态射.
\end{enonce*}

现设 $\mathcal E\to B$ 为 $B$ 上拓扑斯对偶双纤维化的范畴,
$X$ 与 $X'$ 为 $B$ 的两个 $D$-对象, $\alpha:X\to X'$ 为函子性态射.
为 $\mathcal E$ 与 $\mathcal E^{\circ}$ 选定规范化分裂后,
便典范确定一个 $D$-拓扑斯态射 $(\alpha_{\ast},\alpha^{\ast}):\bar X\to\bar X'$,
使对 $D$ 的每个对象\pageoriginale{} $i$,
$(\alpha_{\ast i},\alpha_i^{\ast}):\bar X_i\to\bar X'_i$
等于 $(\alpha_{i\ast}^{\circ},\alpha_i^{\ast\circ})$,
其中 $\alpha_i:X_i\to X'_i$ 为 $\alpha$ 给出的 $B$ 中的箭头.
对 $\mathcal E$ 和 $\mathcal E^{\circ}$ 的两种不同分裂选择,
所得态射之间存在唯一的 $D$-同构.
(关于这些事实的验证, 读者可参阅 \cite{sga4.exp.Vbis.bib.4}, (1.17).)

\subsection{环化 $D$-拓扑斯}\etiquette[1.3]{Vbis.subsec1.3}

\begin{enonce*}{定义 1.3.1}\etiquette[1.3.1]{Vbis.def1.3.1}
环化 $D$-拓扑斯为二元组 $(E,A)$, 其中 $E$ 为 $D$-拓扑斯,
$A$ 为 $\Ysheaf{\Gamma}(E)$ 的环.
\end{enonce*}

随即可验, 对 $D$ 的每个对象 $i$, $A_i$ 为拓扑斯 $E_i$ 的环,
且对每个态射 $m:i\to j$, 典范箭头 $A_i\to m^{\ast}(A_j)$ 为环同态.

\begin{enonce*}{定义 1.3.2}\etiquette[1.3.2]{Vbis.def1.3.2}
设 $(E,A)$ 与 $(E',A')$ 为两个环化 $D$-拓扑斯.
从 $(E,A)$ 到 $(E',A')$ 的态射为二元组 $(\Phi,\theta)$,
其中 $\Phi:E\to E'$ 为 $D$-拓扑斯态射,
$\theta:A'\to\Gamma(\Phi_{\ast})(A)$ 为环同态.
\end{enonce*}

\begin{enonce*}[remark]{注 1.3.3}\etiquette[1.3.3]{Vbis.rem1.3.3}
当 $D$ 为 $\Uu$-小范畴时, 环化 $D$-拓扑斯态射 $\Phi:(E,A)\to(E',A')$
诱导一个环化 $\Uu$-拓扑斯态射
$(\Gamma(\Phi),\theta):(\Ysheaf{\Gamma}(E),A)\to (\Ysheaf{\Gamma}(E'),A')$
(参见 \hyperref[Vbis.prop1.2.15]{1.2.15}).
\end{enonce*}

\setcounter{subsubsection}{3}
\subsubsection{}\etiquette[1.3.4]{Vbis.subsubsec1.3.4}
设 $\mathcal E\to B$ 为 $B$ 上拓扑斯对偶双纤维化的范畴,
$\mathcal O$ 为 $\Ysheaf{\Gamma}(\mathcal E^{\circ})=\Hom_{B^{\circ}}(B^{\circ},\mathcal E^{\circ})$ 的环.
若 $X:D^{\circ}\to B$ 为 $B$ 的 $D$-对象, 则 $D$-拓扑斯 $\bar X$
(参见 \hyperref[Vbis.subsubsec1.2.5]{(1.2.5)})
自然地由 $(\mathcal O\mathbin{.}X^{\circ}):D\to\mathcal E$ 环化;
所得环化 $D$-拓扑斯记为 $(\bar X,\mathcal O)$.
若 $\alpha:X\to X'$ 为函子性态射, 则态射
$(\alpha_{\ast},\alpha^{\ast}):\bar X\to\bar X'$
(参见 \hyperref[Vbis.sch1.2.17]{1.2.17}) 典范诱导环化 $D$-拓扑斯态射
$(\bar X,\mathcal O)\to(\bar X',\mathcal O)$, 仍记为 $\alpha$.

\subsubsection{}\etiquette[1.3.5]{Vbis.subsubsec1.3.5}
环化 $D$-拓扑斯 $(E,A)$ 典范定义一个在 $D$ 上为阿贝尔范畴双纤维化的范畴
$\Ysheaf{\Mod}(E,A)$, 其在 $D$ 的对象 $i$ 处\pageoriginale{}的纤维
为环化拓扑斯 $(E_i,A_i)$ 上的模范畴 $\Mod(E_i,A_i)$.
采用此记号, $\Ysheaf{\Gamma}(E)$ 的 $A$-模范畴,
记为 $\Mod(\Ysheaf{\Gamma}(E),A)$,
等同于范畴 $\Hom_D(D,\Ysheaf{\Mod}(E,A))$.

\subsubsection{}\etiquette[1.3.6]{Vbis.subsubsec1.3.6}
设 $\varphi=(\Phi,\theta):(E,A)\to(E',A')$ 为环化 $D$-拓扑斯态射.
它在相应的模范畴之间定义两个函子
$\varphi_{\ast}:\underline{\Mod}(E,A)\to\underline{\Mod}(E',A')$ 和
$\varphi^{\ast}:\underline{\Mod}(E',A')\to\Mod(E,A)$:

\begin{itemize}
\item 设 $M$ 为 $\Ysheaf{\Mod}(E,A)$ 中 $D$ 的对象 $i$ 上的对象.
则 $\Phi_{\ast}(M)$ 为 $\Phi_{\ast}(A_i)$ 上的模,
而态射 $\theta_i:A'_i\to\Phi_{\ast}(A_i)$ 使之成为 $A'_i$-模,
记为 $\varphi_{\ast}(M)$. 函子 $\varphi_{\ast}$ 称为\emph{态射 $\varphi$ 下的正像函子}.

\item 设 $M'$ 为 $\Ysheaf{\Mod}(E',A')$ 中 $D$ 的对象 $i$ 上的对象.
则 $\Phi^{\ast}(M')$ 为 $\Phi^{\ast}(A'_i)$ 上的模,
且 $\varphi^{\ast}(M')=\Phi^{\ast}(M')\otimes_{\Phi^{\ast}(A'_i)}A_i$
典范配备 $A_i$-模结构. 借助 (\hyperref[Vbis.lem1.2.16]{1.2.16}),
这定义函子 $\varphi^{\ast}$, 它是 $\varphi_{\ast}$ 的左伴随,
称为\emph{态射 $\varphi$ 下的逆像函子}.
\end{itemize}

\paragraph{}\etiquette[1.3.6.1]{Vbis.par1.3.6.1}
若函子 $\Gamma(\varphi^{\ast})$ 正合, 就称 $\varphi$ \emph{平坦}.

\begin{enonce*}{命题 1.3.7}\etiquette[1.3.7]{Vbis.prop1.3.7}
设 $f:D'\to D$ 为函子, $(E,A)$ 为环化 $D$-拓扑斯.
则典范函子
$$
f^{\ast}:\Mod(\Ysheaf{\Gamma}(E),A)\longrightarrow
\Mod(\Ysheaf{\Gamma}(\fprod{E}{D'}{D}),A\circ f)
$$
在 $D'$ 为 $\Uu$-小范畴时, 既有右伴随也有左伴随. 特别地, 它正合.
\end{enonce*}

这立即来自 (\hyperref[Vbis.lem1.2.10]{1.2.10}) 以及等同
$\Hom_D(D,\Ysheaf{\Mod}(E,A))\simeq \Mod(\Ysheaf{\Gamma}(E),A)$.

按照\pageoriginale{}一般记号, 将 $f^{\ast}$ 的左伴随 (分别为右伴随)
记为 $f_!$ (分别为 $f_{\ast}$).

\setcounter{subsubsection}{7}
\subsubsection{}\etiquette[1.3.8]{Vbis.subsubsec1.3.8}
下面的讨论为形如
$$
\sisi{\RR^{+}}{R}\Gamma(\varphi_{\ast}):
D^{+}(\Ysheaf{\Gamma}(E),A)\to
D^{+}(\Ysheaf{\Gamma}(E'),A'),
$$
的导出函子提供方便的计算方法,
其中 $\varphi:(E,A)\to(E',A')$ 为环化 $D$-拓扑斯态射
(参见 (\hyperref[Vbis.subsubsec1.3.6]{1.3.6})).
[若 $(S,\mathcal O_S)$ 为环化拓扑斯, 用 $D^{+}(S,\mathcal O_S)$
表示 $S$ 上 $\mathcal O_S$-模范畴的下有界导出范畴
\nde{若 $f:(S,\mathcal O_S)\rightarrow(T,\mathcal O_T)$ 为环化拓扑斯的平坦态射,
则如原版本一样, 将导出函子
$$f^*:D^+(T,\mathcal O_T)\rightarrow D^+(S,\mathcal O_S),$$
简记为 $f^*$ (而非 $L^+(f^*)$).}.]

\subsubsection{}\etiquette[1.3.9]{Vbis.subsubsec1.3.9}
此外, 此形式计算也可用于其他背景, 例如“$\ell$-进上同调的下降” (参见 \SGA 5).

\begin{enonce*}{命题 1.3.10}\etiquette[1.3.10]{Vbis.prop1.3.10}
设 $D$ 为 $\Uu$-小范畴. 若 $(E,A)$ 为环化 $D$-拓扑斯,
用 $I_{(E,A)}$ 表示 $\Mod(\Ysheaf{\Gamma}(E),A)$ 中
同构于形如 $\prod_{i\in\ob(D)}e_{i\ast}(Q_i)$ 的对象所成的类,
其中对 $D$ 的每个对象 $i$, $Q_i$ 完全无上同调
\emph{(参见 V \hyperref[V.def4.1]{4.1})}\footnote{按第 V 篇的术语称为“松弛”.};
$I_{(E,A)}$ 具有下述性质:
\begin{enumerate}
\renewcommand{\theenumi}{\roman{enumi}}
\renewcommand{\labelenumi}{(\theenumi)}
\item 对 $I_{(E,A)}$ 的每个对象 $F$ 及 $D$ 的每个对象 $i$,
$e_i^{\ast}(F)$ 完全无上同调.

\item $\Mod(\Ysheaf{\Gamma}(E),A)$ 的每个对象均嵌入 $I_{(E,A)}$ 的某个对象.

\item $I_{(E,A)}$ 对有限直和稳定.

\item 对每个环化 $D$-拓扑斯态射 $\varphi:(E,A)\to(E',A')$,
函子 $\Gamma(\varphi_{\ast})$ 将 $C^{+}(\Ysheaf{\Gamma}(E),A)$ 中
由 $I_{(E,A)}$ 的对象组成的每个无上同调复形,
映为由 $I_{(E',A')}$ 的对象组成的无上同调复形.
\end{enumerate}
\end{enonce*}

\noindent{\underline{\emph{证明}}}

\begin{enumerate}
\renewcommand{\theenumi}{\roman{enumi}}
\renewcommand{\labelenumi}{(\theenumi)}
\item 根据 $e_i^{\ast}$ 的显式表达式
(参见 \hyperref[Vbis.cor1.2.11]{(1.2.1.1)}), 只须证明下述引理:

\begin{enonce*}{引理 1.3.10.1}\etiquette[1.3.10.1]{Vbis.lem1.3.10.1}
设\pageoriginale{} $(S,\mathcal O_S)$ 为环化拓扑斯,
$(F_t)_{t\in T}$ 为以 $\Uu$-小集 $T$ 为指标的一族完全无上同调的 $\mathcal O_S$-模.
则 $\prod_{t\in T}F_t$ 完全无上同调.
\end{enonce*}

设 $X$ 为 $S$ 的对象. 有谱序列
\begin{equation*}
H^p\left(X,{\mathop{\prod}_{t\in T}}^{(q)}F_t\right)\longrightarrow
\mathop{\prod}_{t\in T}H^{p+q}(X,F_t)\tag{1.3.10.2}
\end{equation*}\etiquette[1.3.10.2]{Vbis.eq1.3.10.2}
其中 $\prod_{t\in T}^{(q)}$ 表示“以 $T$ 为指标的积”函子的第 $q$ 个导出函子.
由于 $\prod_{t\in T}^{(q)}F_t$ 是预层 $U\mapsto\prod_{t\in T}H^q(U,F_t)$ 的伴随层,
\eqhyperref[Vbis.eq1.3.10.2]{1.3.10.2} 退化并给出同构
\begin{equation*}
H^n\left(X,\prod_{t\in T}F_t\right)\xrightarrow{\sim}
\prod_{t\in T}H^n(X,F_t)\tag{1.3.10.3}
\end{equation*}\etiquette[1.3.10.3]{Vbis.eq1.3.10.3}
从而最终得到对每个 $n>0$, 有 $H^n(X,\prod_{t\in T}F_t)=0$.

\item 设 $F$ 为 $\Mod(\Ysheaf{\Gamma}(E),A)$ 的对象, 令 $F_i=e_i^{\ast}(F)$.
对每个 $i$, 选取单态射 $F_i\to Q_i$,
其中 $Q_i$ 在 $\Mod(E_i,A_i)$ 中完全无上同调.
则 $e_{i\ast}(F_i)\to e_{i\ast}(Q_i)$ 为单态射
(因为 $e_{i\ast}$ 有左伴随), 且典范箭头
\begin{equation*}
F\to\prod_{i\in\Ob(D)}e_{i\ast}(F_i)\to
\prod_{i\in\Ob(D)}e_{i\ast}(Q_i)\tag{1.3.10.4}
\end{equation*}\etiquette[1.3.10.4]{Vbis.eq1.3.10.4}
仍为单态射, 这立即可验.

\item 证明留待读者完成.

\item 读者还可验证 $\Gamma(\varphi_{\ast})$ 将 $I_{(E,A)}$ 的对象映为
$I_{(E',A')}$ 的对象. 这一点确定后, 由 (i), 只须验证下述引理:
\end{enumerate}

\begin{enonce*}{引理 1.3.10.5}\etiquette[1.3.10.5]{Vbis.lem1.3.10.5}
若\pageoriginale{} $0\to F\to G\to H\to0$ 为 $\Mod(\Ysheaf{\Gamma}(E),A)$ 中的短正合列,
其中 $F,G\in I_{(E,A)}$, 则序列
$$
0\to\Gamma(\varphi_{\ast})(F)\to
\Gamma(\varphi_{\ast})(G)\to
\Gamma(\varphi_{\ast})(H)\to0
$$
在 $\Mod(\Ysheaf{\Gamma}(E'),A')$ 中正合.
\end{enonce*}

只须注意, 对 $D$ 的每个对象 $i$, $e_i^{\ast}(H)$ 无上同调,
并且 $\Gamma(\varphi_{\ast})$ 逐纤维计算.
这就完成 (\hyperref[Vbis.prop1.3.10]{1.3.10}) 的证明.

\begin{enonce*}{推论 1.3.11}\etiquette[1.3.11]{Vbis.cor1.3.11}
设 $\varphi:(E,A)\to(E',A')$ 为环化 $D$-拓扑斯态射.
函子 $\sisi{\RR^+}{R}\Gamma(\varphi_{\ast})$ 可借助由 $I_{(E,A)}$ 的对象组成的分解计算.
\end{enonce*}

设 $L\subset K^{+}(\Ysheaf{\Gamma}(E),A)$ 为由 $I_{(E,A)}$ 的对象组成的闭复形的满子范畴.
由 ((\hyperref[Vbis.prop1.3.10]{1.3.10}), (ii)), $L$ 为三角子范畴,
[\cite{sga4.exp.Vbis.bib.7}, 第 1 章] 的定理 (\hyperref[Vbis.subsec5.1]{5.1}) 可应用于它.

\begin{enonce*}{推论 1.3.12}\etiquette[1.3.12]{Vbis.cor1.3.12}
设 $\varphi:(E,A)\to(E',A')$ 为环化 $D$-拓扑斯态射.
对 $D$ 的每个对象 $i$, 用 $\varphi_i:(E_i,A_i)\to(E'_i,A'_i)$
表示在 $i$ 上诱导的环化拓扑斯态射. 则图
\[
\xymatrix@C=2.75cm@R=1.75cm{
D^{+}(\Gamma(E),A) \ar[r]^{\sisi{\RR^{+}(e_i^{\ast})}{e_i^{\ast}}}
\ar[d]_{\sisi{\RR^{+}}{R}\Gamma(\varphi_{\ast})}&
D^{+}(E_i,A_i) \ar[d]^{\sisi{\RR^{+}}{R}\varphi_{i\ast}}\\
D^{+}(\Gamma(E'),A')
\ar[r]_{\sisi{\RR^{+}(e_i^{\ast})}{e_i^{\ast}}}&D^{+}(E'_i,A'_i)
}
\]
本质交换.
\end{enonce*}

存在态射
$$
\sisi{\RR^{+}(e_i^{\ast})}{e_i^{\ast}}\circ
\sisi{\RR^{+}}{R}\Gamma(\varphi_{\ast})\to
\sisi{\RR^{+}}{R}\varphi_{i\ast}\circ
\sisi{\RR^{+}(e_i^{\ast})}{e_i^{\ast}}
$$
由 (\hyperref[Vbis.prop1.3.10]{1.3.10}) 可验证它是同构.

\begin{enonce*}[remark]{注 1.3.13}\etiquette[1.3.13]{Vbis.rem1.3.13}
用\pageoriginale{} $\mathcal A_{(E,A)}$ 表示 $\Mod(\Ysheaf{\Gamma}(E),A)$ 中
使对 $D$ 的每个对象 $i$, $e_i^{\ast}(F)$ 完全无上同调的对象 $F$ 所成的类.
由 (\hyperref[Vbis.prop1.3.10]{1.3.10}) 的证明得出,
$\sisi{\RR^{+}}{R}\Gamma(\varphi_{\ast})$ 可借助由 $\mathcal A_{(E,A)}$ 的对象组成的分解计算;
特别地, 引言中使用“典范松弛分解”\footnote{此概念的推广参见第 XVII 篇.} 就是如此.
由 [(\hyperref[Vbis.prop1.3.10]{1.3.10})(i)], 有 $I_{(E,A)}\subset\mathcal A_{(E,A)}$,
但在 \hyperref[Vbis.sec2]{第 2 节} 我们将明确使用 $I_{(E,A)}$.
\end{enonce*}

% END INCLUDED SGA4_English_current_buildable_source_20260906/source/05bis/drafts/SGA4_VBIS_LINES_517_829_English.texfrag


% BEGIN INCLUDED SGA4_English_current_buildable_source_20260906/source/05bis/drafts/SGA4_VBIS_LINES_830_1073_English.texfrag
\section{上同调下降方法}\etiquette[2]{Vbis.sec2}

本节中, $D$ 为 $\Uu$-小范畴.

\subsection{一般讨论与记号}\etiquette[2.1]{Vbis.subsec2.1}

\subsubsection{}\etiquette[2.1.1]{Vbis.subsubsec2.1.1}
设 $(S,\mathcal{O}_S)$ 为环化拓扑斯.
范畴 $S\times D$ 配备典范投影 $S\times D\to D$ 后, 为 $D$-拓扑斯;
此外, 取值为 $\mathcal{O}_S$ 的常值截面定义一个环化 $D$-拓扑斯
$(S\times D,\mathcal{O}_S)$,
称为\emph{取值为 $(S,\mathcal{O}_S)$ 的常值环化 $D$-拓扑斯}.
采用此记号, 范畴 $\Mod(\Ysheaf{\Gamma}(S\times D),\mathcal{O}_S)$
等同于从 $D$ 到 $\Mod(S,\mathcal{O}_S)$ 的协变函子范畴.

将 $S$ 上的每个模 $F$ 对应到取值为 $F$ 的常值函子, 定义正合函子
$$
\epsilon^{\ast}:\Mod(S,\mathcal{O}_S)\to
\Mod(\Ysheaf{\Gamma}(S\times D),\mathcal{O}_S)
$$
函子 $\epsilon^{\ast}$ 有右伴随 $\epsilon_\ast$,
它将每个函子 $H:D\to\Mod(\mathcal{O}_S)$ 对应到其投射极限;
伴随态射 $F\to\epsilon_\ast\epsilon^{\ast}(F)$
是从 $F$ 到取值为 $F$ 的常值函子的投射极限的态射.
因此 $\epsilon^{\ast}$ 全忠实当且仅当 $D$ 连通.

\begin{enonce*}{定义 2.1.2}\etiquette[2.1.2]{Vbis.def2.1.2}
设\pageoriginale{} $(E,A)$ 为环化 $D$-拓扑斯.
$(E,A)$ 的增广是从 $(E,A)$ 到常值环化 $D$-拓扑斯的一个 (环化 $D$-拓扑斯) 态射.
配备增广的环化 $D$-拓扑斯称为增广环化 $D$-拓扑斯.
\end{enonce*}

\setcounter{subsubsection}{2}
\subsubsection{}\etiquette[2.1.3]{Vbis.subsubsec2.1.3}
设 $\mathcal{E}\to B$ 为拓扑斯对偶双纤维化的范畴,
$\mathcal{O}$ 为
$\Ysheaf{\Gamma}(\mathcal{E}^{\circ})=\Hom_{B^{\circ}}(B^{\circ},\mathcal{E}^{\circ})$ 中的环.
设 $S$ 为 $B$ 的对象. 函子 $D^{\circ}\to B/S$,
即函子性态射 $X\xrightarrow{\theta}C^D_S$,
其中 $X$ 为函子 $D^{\circ}\to B$ (参见 (\hyperref[Vbis.subsubsec1.2.6]{1.2.6})),
诱导一个增广
$$
\displaylines{\hfill
\theta:(\bar{X},\mathcal{O})\longrightarrow
(\mathcal{E}^{\circ}_S\times D,\mathcal{O}_S)
\hfill\llap{(\cf \hyperref[Vbis.subsubsec1.3.4]{(1.3.4)}){\quoi}.}}
$$
我们将用与产生它的函子性态射相同的字母表示它.


\subsection{上同调下降}\etiquette[2.2]{Vbis.subsec2.2}

\subsubsection{}\etiquette[2.2.1]{Vbis.subsubsec2.2.1}
设 $\theta:(E,A)\to(S\times D,\mathcal{O}_S)$ 为增广环化 $D$-拓扑斯.
采用 (\hyperref[Vbis.subsubsec1.3.6]{1.3.6})
及 (\hyperref[Vbis.subsubsec2.1.1]{2.1.1}) 的记号, 令
$$
\bar{\theta}^{\ast}=\Gamma(\theta^{\ast})\circ\epsilon^{\ast}:
\Mod(S,\mathcal{O}_S)\to\Mod(\Ysheaf{\Gamma}(E),A)
$$
以及
$\bar{\theta}_\ast=\epsilon_\ast\circ\Gamma(\theta_\ast):\Mod(\Ysheaf{\Gamma}(E),A)\to\Mod(S,\mathcal{O}_S)$.

$\bar{\theta}^{\ast}$ 的像属于 $\Mod(\Ysheaf{\Gamma}(E),A)$ 的满子范畴,
该子范畴由 $\Ysheaf{\Gamma}(D,\Ysheaf{\Mod}(E,A))$ 的余笛卡尔截面组成.
我们将此范畴记为 $\Ysheaf{\Gamma}^{\text{cocart}}(D,\Ysheaf{\Mod}(E,A))$
\nde{换言之, 一个增广环化 $D$-拓扑斯, 就是给出一个相容的 $S$-拓扑斯态射系统
$$(m^\ast,m_\ast):E_j\longrightarrow E_i,$$
其中 $m\in\Hom_D(i,j)$.
余笛卡尔截面范畴就是下降资料范畴; 其对象为配备同构余循环
$$F_j\isomto m^\ast F_i.$$
的族 $F_i\in\text{Ob}(E_i)$.
将 $S$ 上的每个对象 $F$ 对应到取值为 $F$ 的常值函子的函子,
是拓扑斯态射
$$\epsilon:\Gamma(S\times D)\longrightarrow S,$$
的逆像, 它诱导 \hyperref[Vbis.subsubsec2.1.3]{2.1.3} 的环化拓扑斯态射.
因此二元组 $(\bar\theta_\ast,\bar\theta^\ast)$
就是由拓扑斯态射的复合
$\bar\theta=\epsilon\circ\Gamma(\theta):\Gamma(E)\to S$
诱导的 (参见 \hyperref[Vbis.prop1.2.15]{1.2.15}).}

\begin{enonce*}{定义 2.2.2}\etiquette[2.2.2]{Vbis.def2.2.2}
若
$$
\sisi{}{\bar\theta^{\ast}=}
\Gamma(\theta^{\ast})\circ\epsilon^{\ast}:
\Mod(S,\mathcal{O}_S)\to
\Ysheaf{\Gamma}^{\text{\rm cocart}}(D,\Ysheaf{\Mod}(E,A))
$$
为范畴等价, 就称态射 $\theta$ 为\emph{有效下降增广}.
\end{enonce*}

\setcounter{subsubsection}{2}
\subsubsection{}\etiquette[2.2.3]{Vbis.subsubsec2.2.3}
采用\pageoriginale{} (\hyperref[Vbis.subsubsec2.2.1]{2.2.1}) 的记号,
假定 $\theta$ 平坦, 于是 $\bar\theta^{\ast}$ 正合, 并直接过渡到导出范畴;
将由此得到的函子\sisi{记为 $\LL^{+}(\bar\theta^{\ast})$}{仍记为 $\bar\theta^{\ast}$}.
采用此记号:

\begin{enonce*}{引理 2.2.3.1}\etiquette[2.2.3.1]{Vbis.lem2.2.3.1}
存在两个函子性态射
$$
\alpha:\id_{D^{+}(S,\mathcal{O}_S)}\to
\sisi{\RR^{+}}{R}\bar\theta_\ast\circ
\sisi{\LL^{+}(\bar\theta^{\ast})}{\bar\theta^{\ast}}
\quad\mbox{和}\quad
\beta:\sisi{\LL^{+}(\bar\theta^{\ast})}{\bar\theta^{\ast}}\circ
\sisi{\RR^{+}}{R}\bar\theta_\ast\to
\id_{D^{+}(\Ysheaf{\Gamma}(E),A)}
$$
使这两个函子互为伴随.
\end{enonce*}

(参见 \cite{sga4.exp.Vbis.bib.9} (3.3)).

\begin{enonce*}{定义 2.2.4}\etiquette[2.2.4]{Vbis.def2.2.4}
若
\begin{enumerate}
\renewcommand{\theenumi}{\arabic{enumi}}
\renewcommand{\labelenumi}{\theenumi$^{\circ}$)}
\item $\theta$ 平坦;
\item
\sisi{$\LL^{+}(\bar\theta^{\ast})$}
{$\bar\theta^{\ast}:D^{+}(S,\mathcal{O}_S)\to
D^{+}(\Ysheaf{\Gamma}(E),A)$} 全忠实,
\end{enumerate}
就称态射 $\theta$ 为 \emph{1-上同调下降增广}.
\end{enonce*}

\begin{enonce*}[remark]{注 2.2.5}\etiquette[2.2.5]{Vbis.rem2.2.5}
前一定义的条件 $2^{\circ}$) 也可表述为:
(\hyperref[Vbis.lem2.2.3.1]{2.2.3.1}) 中的态射
$\alpha\sisi{}{:\id_{D^{+}(S,\mathcal{O}_S)}\to R\bar\theta_\ast\circ\bar\theta^{\ast}}$ 为同构.
\sisi{}{这正是实际中所证明的, 也是有用的性质.}
\end{enonce*}

\begin{enonce*}{定义 2.2.6}\etiquette[2.2.6]{Vbis.def2.2.6}
若 $\theta$ 既是有效下降增广又是 1-上同调下降增广,
就称态射 $\theta$ 为 \emph{2-上同调下降增广} (或\emph{有效上同调下降增广}).
\end{enonce*}

下述结果说明了上述术语的合理性:

\begin{enonce*}{命题 2.2.7}\etiquette[2.2.7]{Vbis.prop2.2.7}
设 $\theta$ 为有效上同调下降增广.
则 $\sisi{\LL^{+}(\bar\theta^{\ast})}{\bar\theta^{\ast}}$ 的本质像,
为 $D^{+}(\Ysheaf{\Gamma}(E),A)$ 中由满足下述条件的复形 $F^{\bigdot}$ 组成的满子范畴:
对每个 $i$, $H^{i}(F^{\bigdot})$ 为 $\Ysheaf{\Mod}(E,A)$ 的余笛卡尔截面.
\end{enonce*}

为简便起见, 将满足上述条件的复形 $F^{\bigdot}$ 称为\emph{上同调下降资料}.

\sisi{}{\begin{proof}}
显然, 若 $K^{\bigdot}\in D^{+}(S,\mathcal{O}_S)$,
\sisi{}{则}
$\sisi{\LL^+(\bar\theta^{\ast})}{\bar\theta^{\ast}(K^{\bigdot})}$ 为上同调下降资料.
因而只须证明, 对每个上同调下降资料 $F^{\bigdot}$, 典范态射
$$
\beta(F^{\bigdot}):
\sisi{\LL^{+}(\bar\theta^{\ast})}{\bar\theta^{\ast}}\circ
\sisi{\RR^{+}(\bar\theta_\ast)}{R\bar\theta_\ast}(F^{\bigdot})
\to F^{\bigdot}
$$
在\pageoriginale{} $D^{+}(\Ysheaf{\Gamma}(E),A)$ 中为同构.

\begin{enumerate}
\renewcommand{\theenumi}{\alph{enumi}}
\renewcommand{\labelenumi}{\theenumi)}
\item \emph{$F^{\bigdot}$ 有界的情形}: 对使 $H^{i}(F^{\bigdot})\neq0$ 的整数
$i$ 所在区间的长度 $\ell$ 归纳:
\begin{itemize}
\item 当 $\ell\leq1$ 时, 因 $\theta$ 为有效下降, 断言成立.

\item 设 $\ell>1$, 并令 $n$ 为使 $H^n(F^{\bigdot})\neq0$ 的最大整数.
有正合三角\footnote{对每个复形 $K^{\bigdot}$,
$K^{\bigdot}[-n]$ 表示复形 $T^{-n}(K^{\bigdot})$, 其中 $T$ 为平移函子.}
\[
\xymatrix@R=1.5cm{
          & H^{n}(F^{\bigdot})[-n] \ar[ld]&\\
F'^{\bigdot} \ar[rr] & & F^{\bigdot} \ar[lu]
}
\]
使归纳假设适用于 $F'^{\bigdot}$.
由于 $\beta(H^{n}(F^{\bigdot})[-n])$ 与 $\beta(F'^{\bigdot})$ 为同构,
$\beta(F^{\bigdot})$ 也为同构.
\end{itemize}

\item \emph{一般情形}: 用 $\sisi{\sigma}{\tau}_{\leq n}(F^{\bigdot})$ 表示复形
$$
\ldots\to F^{n-2}\to F^{n-1}\to\Ker d^n\to0\to0\ldots{\quoi},
$$
(参见 \cite{sga4.exp.Vbis.bib.7} (7.1)), 从而对每个 $n$, 有交换图
\[
\xymatrix@R=1.5cm@C=2cm{
\bar{\theta}^{\ast}\circ R\bar{\theta}_{\ast}
(\tau_{\leq n}(F^{\bigdot}))
\ar[r]^{\sim}\ar[d]&
\tau_{\leq n}(F^{\bigdot})\ar[d]\\
\bar{\theta}^{\ast}\circ R\bar{\theta}_{\ast}
(F^{\bigdot})\ar[r]&F^{\bigdot}
}
\]
只须看出态射
$$
\sisi{\RR^{+}(\bar\theta_\ast)}{R\bar\theta_\ast}
(\sisi{\sigma}{\tau}_{\leq n}(F^{\bigdot}))
\to
\sisi{\RR^{+}(\bar\theta_\ast)}{R\bar\theta_\ast}(F^{\bigdot})
$$
对 $i\leq n$ 在 $H^{i}$ 上诱导\pageoriginale{}同构.
但有三角
\[
\xymatrix{
&F''^{\bigdot}\ar[ld]&\\
\tau_{\leq n}(F^{\bigdot})\ar[rr]&&
F^{\bigdot}\ar[lu]
}
\]
其中对 $i\leq n$, 有 $H^i({F''}^{\bigdot})=0$.
由此得对 $i\leq n$, 有
$H^{i}(\sisi{\RR^{+}(\bar\theta_\ast)}{R\bar\theta_\ast}(F''^{\bigdot}))=0$,
继而对 $i\leq n$,
$H^{i}(\sisi{\RR^{+}(\bar\theta_\ast)}{R\bar\theta_\ast}(\sisi{\sigma}{\tau}_{\leq n}(F^{\bigdot})))\to
H^{i}(\sisi{\RR^{+}(\bar\theta_\ast)}{R\bar\theta_\ast}(F^{\bigdot}))$
为同构, 这就完成了证明.
\end{enumerate}

\sisi{}{\end{proof}}

现在给出一个显式方法来计算
$$
\sisi{\RR^{+}}{R}\bar\theta_\ast
\sisi{}{=R\epsilon_\ast\circ R\Gamma(\theta)_\ast},
$$
这对于证明上同调下降的某些判别准则及应用此概念很有用.

\sisi{\subsection{计算
$R^+\epsilon_\ast$ 的一种方法}\etiquette[2.3]{Vbis.subsec2.3}}
{\subsection{计算
$R\epsilon_\ast$ 的一种方法}\etiquette[2.3]{Vbis.subsec2.3}}

先给出两组连锁命题.

% END INCLUDED SGA4_English_current_buildable_source_20260906/source/05bis/drafts/SGA4_VBIS_LINES_830_1073_English.texfrag


% BEGIN INCLUDED SGA4_English_current_buildable_source_20260906/source/05bis/drafts/SGA4_VBIS_LINES_1074_1378_English.texfrag
\subsubsection{}\etiquette[2.3.1]{Vbis.subsubsec2.3.1}
设 $D$ 为 $\Uu$-小范畴, $\Ab$ 为属于 $\Uu$ 的阿贝尔群的范畴.
用 $r_D$ 表示由下式定义的函子 $D\to(\Hom(D,\Ab))^{\circ}$:
$$
r_D(i)(j)=\ZZ^{\Hom(i,j)}
$$
此式对 $D$ 的每对对象 $(i,j)$ 成立.
对由 $D$ 的对象和 $\Hom(D,\Ab)$ 的对象组成的每一对 $(i,X)$,
可构造一个典范且函子性的同构
$$
\Hom(r_D(i),X)\xrightarrow{\sim}X(i).
$$

\paragraph{}\etiquette[2.3.1.1]{Vbis.par2.3.1.1}
设 $A$ 为加性范畴 (参见 \sisi{Tohokû}{Tôhoku}), $F$ 为函子 $D\to A$.
由下式定义函子 $\bar F:(\Hom(D,\Ab))^{\circ}\to\Hom(A^{\circ},\Ens)$:
$$
\bar F(X)(a)=\Hom_{\Hom(D,\Ab)}
(X,\Hom_A(a,F(.))),
$$
从而\pageoriginale{}下图本质交换, 其中 $h_A$ 表示典范函子:
\[
\xymatrix@C=3cm@R=1.5cm{
D \ar[r]^{r_D}\ar[d]_{F}&(\Hom(D,\Ab))^{\circ}\ar[d]^{\bar F}\\
A \ar[r]_{h_A}&\Hom(A^{\circ},\Ens).
}
\]

请读者验证, $\bar F$ 将 $\Hom(D,\Ab)$ 中的有限直和映为积,
且若用 $\ZZ$ 表示取值为 $\ZZ$ 的常值函子 $D\to\Ab$,
则 $\bar F(\ZZ)=\varprojlim F$.

最后, 对应 $F\rightsquigarrow\bar F$ 定义协变函子
$$
\Hom(D,A)\to\Hom((\Hom(D,\Ab))^{\circ},
\Hom(A^{\circ},\Ens)){\quoi}.
$$

\paragraph{}\etiquette[2.3.1.2]{Vbis.par2.3.1.2}
用 $\Add(D)$ 表示 $(\Hom(D,\Ab))^{\circ}$ 的满子范畴,
其由形如 $r_D(i)$ (其中 $i$ 为 $D$ 的对象) 的对象及其有限直和定义.
范畴 $\Add(D)$ 加性, 且函子 $r_D:D\to\Add(D)$ 具有下述泛性质:

\paragraph{}\etiquette[2.3.1.3]{Vbis.par2.3.1.3}
\emph{对每个加性范畴 $A$, 函子 $G\longmapsto G\circ r_D$
诱导从 $\Hom(\Add(D),A)$ 中由加性函子组成的满子范畴到范畴 $\Hom(D,A)$ 的等价.}

范畴 $\Add(D)$ 由 (\hyperref[Vbis.par2.3.1.3]{2.3.1.3}) 确定到等价,
称为\emph{由 $D$ 生成的加性范畴}.

%% compilation problem!

%\nde{In more concrete terms, used implicitly below, the objects of
%$\Add(D)$ are the finite sequences
%$$[i_1]\oplus\cdots\oplus[i_n]:=(i_1,\ldots,i_n)$$
%of objects of $D$, with morphisms in $\Add(D)$ defined by the rules
%\begin{itemize}
%  \item the empty sequence is the zero object;
%  \item for all $i,j\in\text{ob}(D)$,
%      $\Hom([i],[j])=\ZZ^{\Hom_D(i,j)}$;
%  \item for all finite families $(i),(j)$ of objects of $D$,
%      $\Hom(\oplus[i],\oplus[j])=\prod\Hom([i],[j])$,
%\end{itemize}
%while the composition rules are deduced from those of $D$ by
%bilinearity. With this presentation, $r_D$ is the tautological
%functor $i\rightsquigarrow[i]$; and, if $F$ is a functor from $D$ to
%an additive category $A$, it extends uniquely to an additive functor
%$$\widetilde F:\Add(D)\longrightarrow A$$
%defined by $\widetilde F(\oplus[i])=\oplus F_i$, with the evident
%functoriality rules.}

\subsubsection{}\etiquette[2.3.2]{Vbis.subsubsec2.3.2}
设 $A$ 为阿贝尔范畴, $K^{\bigdot\bigdot}$ 为双复形,
将其视为 $C^{+}(C^+(A))$ 的对象, 第一个指标对应于第一个\pageoriginale{}符号 $C^+$.
用 $(K^{\bigdot\bigdot})_s$ 表示\emph{相伴单复形}
(参见 \sisi{Tohokû}{Tôhoku} (\hyperref[Vbis.subsec2.4]{2.4}), 并沿用那里的记号).
这定义函子
\begin{equation*}
(\quad)_s:C^+(C^+(A))\to C^+(A)\tag{2.3.2.1}
\end{equation*}\etiquette[2.3.2.1]{Vbis.eq2.3.2.1}
下述引理留待读者验证:

\begin{enonce*}{引理 2.3.2.2}\etiquette[2.3.2.2]{Vbis.lem2.3.2.2}
函子 $(\quad)_s$ 定义保持拟同构的三角函子
$$
K^+(C^+(A))\to K^+(A)
$$
因而定义三角函子
$$
D^+(C^+(A))\to D^+(A){\quoi},
$$
仍记为 $(\quad)_s$.
\end{enonce*}

\begin{enonce*}[remark]{注}
对函子 $(\quad)_s:C(C^b(A))\to C(A)$ 也有同一结论,
因为由 (\EGA III (11.3.3)), 所考虑的谱序列双正则.
\end{enonce*}

\subsubsection{}\etiquette[2.3.3]{Vbis.subsubsec2.3.3}
现设 $(S,\mathcal{O}_S)$ 为环化拓扑斯, $D$ 为 $\Uu$-小范畴,
$Z^{\bigdot}\in C^{+}(\Add(D))$ 为上链复形, 且当 $n<0$ 时 $Z^n=0$.

\sisi{$\Mod(\Ysheaf{\Gamma}(S\times D),\mathcal{O}_S)$ 的每个对象 $F$
(由 (\hyperref[Vbis.subsubsec2.1.1]{2.1.1}), 它等同于函子 $D\to\Mod(S,\mathcal{O}_S)$)
通过 $\Add(D)$ 的泛性质定义 $C^{+}(S,\mathcal{O}_S)$ 的对象
$\varepsilon_{Z^{\bigdot}\ast}(F)^{\bigdot}$,
由 (\hyperref[Vbis.par2.3.1.3]{2.3.1.3}), 它对 $F$ 函子性地变动}
{$\Mod(\Ysheaf{\Gamma}(S\times D),\mathcal{O}_S)$ 的每个对象 $F$
(由 (\hyperref[Vbis.subsubsec2.1.1]{2.1.1}), 它等同于函子 $D\to\Mod(S,\mathcal{O}_S)$)
通过 $\Add(D)$ 的泛性质 (\hyperref[Vbis.par2.3.1.3]{2.3.1.3}) 定义加性函子
$$\widetilde F:\Add(D)\longrightarrow\Mod(S,\mathcal{O}_S).$$
像复形
$$\varepsilon_{Z^{\bigdot}\ast}(F):=\widetilde F(Z^{\bigdot})$$
为 $C^{+}(S,\mathcal{O}_S)$ 的对象, 对 $F$ 函子性地变动.
这样便定义了一个显然正合的加性函子
$$
\varepsilon_{Z^{\bigdot}\ast}:
\Mod(\Ysheaf{\Gamma}(S\times D),\mathcal{O}_S)
\longrightarrow C^{+}(S,\mathcal{O}_S),
$$
}.

这定义保持拟同构的三角函子
\nde{$K^+$ 表示下有界复形的同伦范畴.}:
\begin{equation*}
K^{+}(\varepsilon_{Z^{\bigdot}\ast}):
K^+(\Ysheaf{\Gamma}(S\times D),\mathcal{O}_S)
\to K^+(C^+(S,\mathcal{O}_S)){\quoi}.\tag{2.3.3.1}
\end{equation*}\etiquette[2.3.3.1]{Vbis.eq2.3.3.1}
此外, 由 (\hyperref[Vbis.lem2.3.2.2]{2.3.2.2}) 得出复合函子
$$
(\quad)_s\circ K^{+}(\varepsilon_{Z^{\bigdot}\ast}):
K^+(\Ysheaf{\Gamma}(S\times D),\mathcal{O}_S)
\to K^{+}(S,\mathcal{O}_S)
$$
将无上同调对象映为无上同调对象. 因而它定义三角\pageoriginale{}函子
\begin{equation*}
\sisi{\RR^+}{R}\varepsilon_{Z^{\bigdot}\ast}:
D^+(\Ysheaf{\Gamma}(S\times D),\mathcal{O}_S)
\to D^{+}(S,\mathcal{O}_S)\tag{2.3.3.2}
\end{equation*}\etiquette[2.3.3.2]{Vbis.eq2.3.3.2}
满足关系
\nde{换言之, $\varepsilon_{Z^{\bigdot}\ast}$ 直接导出为
$F^{\bigdot}\rightsquigarrow \varepsilon_{Z^{\bigdot}\ast}(F^{\bigdot})$.}
$Q\circ(\quad)_s\circ K^+(\varepsilon_{Z^{\bigdot}\ast})
=\sisi{\RR^{+}}{R}\varepsilon_{Z^{\bigdot}\ast}\circ Q$
\footnote{对每个阿贝尔范畴 $A$, 字母 $Q$ 表示典范函子 $K^{+}(A)\to D^{+}(A)$.}

\subsubsection{}\etiquette[2.3.4]{Vbis.subsubsec2.3.4}
保留上述资料\sisi{.}{, 并令 $A=C^+(S,\mathcal{O}_S)$.
%%%%%%%
%%%%%%
首先, $\widetilde F$ 与 $\bar F$ 相容, 即下列 Yoneda 型图交换:
$$
\xymatrix{
&A\ar[r]^h&\Hom(A^\circ,\Ens)\\
D\ar[rd]_{{r_D}}\ar[r]\ar[ru]^{F}&
\Add(D)\ar[d]_{\widetilde{r_D}}\ar[u]_{\widetilde F}\\
&\Hom(D,\Ab)\ar[ruu]_{\bar F}
}
$$
其中 $h$ 和 $r_D$ 分别为下式定义的 Yoneda 型箭头:
$$
h(a)=(\alpha\longmapsto\Hom_A(\alpha,a))
\quad\text{和}\quad
r_D(d)\stackrel{\hyperref[Vbis.subsubsec2.3.1]{2.3.1}}{=}
(\delta\longmapsto\Hom_D(d,\delta)),
$$
即 $h\circ\widetilde F=\bar F\circ\widetilde{r_D}$. 由此得
$$
h(\varepsilon_{Z^{\bigdot}\ast}(F))
=\bar F(\widetilde{r_D}(Z^{\bigdot})).
$$

由投射极限的定义, $h(\epsilon_\ast F)=\varprojlim F$,
且如已指出的, 重言地有 $\bar F(\ZZ)=\varprojlim F$; 故
$$\bar F(\ZZ)=h(\epsilon_\ast(F)).$$
%%%%%%
%%%%%%%
}


将 $Z^{\bigdot}$ 视为 $\Hom(D,\Ab)$ 中的链复形
\sisi{.}{, 即借助函子 $\widetilde r_D$ 将它视为阿贝尔范畴
$\Hom(D,\Ab)^{\circ}$ 的复形.
因此可以谈论正合列和分解. 明确地说, $\Add(D)$ 中的序列
$$Z\longrightarrow Z'\longrightarrow Z''$$
正合, 当且仅当对 $D$ 的每个对象 $j$, 序列
$$
\Hom_{\Add(D)}(Z'',[j])\longrightarrow
\Hom_{\Add(D)}(Z',[j])\longrightarrow
\Hom_{\Add(D)}(Z,[j])
$$
在 $\Ab$ 中正合.}
设 $t:Z^0\to\ZZ$ 为从 $Z^0$ 到取值为 $\ZZ$ 的常值函子的态射,
使复合态射
\begin{equation*}
Z^1\to Z^0\xrightarrow{t}\ZZ\tag{2.3.4.1}
\end{equation*}\etiquette[2.3.4.1]{Vbis.eq2.3.4.1}
为零.

由函子性 (参见 (\hyperref[Vbis.par2.3.1.1]{2.3.1.1})), 得到态射
$\varepsilon_\ast(F)\to\varepsilon_{Z^{\bigdot}\ast}(F)^0$, 使复合态射
\begin{equation*}
\varepsilon_\ast(F)\to\varepsilon_{Z^{\bigdot}\ast}(F)^0\to
\varepsilon_{Z^{\bigdot}\ast}(F)^1\tag{2.3.4.2}
\end{equation*}\etiquette[2.3.4.2]{Vbis.eq2.3.4.2}
为零.

于是存在三角函子的典范态射
\begin{equation*}
Q\circ K^+(\varepsilon_\ast)\xrightarrow{j}
Q\circ(\quad)_s\circ K^+(\varepsilon_{Z^{\bigdot}\ast})
=\sisi{\RR^{+}}{R}\varepsilon_{Z^{\bigdot}\ast}\circ Q
\tag{2.3.4.3}
\end{equation*}\etiquette[2.3.4.3]{Vbis.eq2.3.4.3}

\begin{enonce*}{命题 2.3.5}\etiquette[2.3.5]{Vbis.prop2.3.5}
保留 (\hyperref[Vbis.subsubsec2.3.3]{2.3.3}) 及
{(\hyperref[Vbis.subsubsec2.3.4]{2.3.4})} 的条件与记号;
再假定增广复形
$$
\ldots\to Z^n\to Z^{n-1}\to\ldots\to Z^1\to Z^0
\xrightarrow{t}\ZZ
$$
定义 $\ZZ$ 的分解.

则对由 $I_{(S\times D,\mathcal{O}_S)}$ 的对象组成的每个复形 $F^{\bigdot}$
(参见 \hyperref[Vbis.prop1.3.10]{1.3.10}),
$$
j(F^{\bigdot}):
Q\circ K^+(\varepsilon_\ast)(F^{\bigdot})\to
(\sisi{\RR^+}{R}\varepsilon_{Z^{\bigdot}\ast}\circ Q)(F^{\bigdot})
$$
为\pageoriginale{}同构.
\end{enonce*}

\sisi{$+$}{} 设 $i$ 为 $D$ 的对象, $Q_i$ 为 $(S,\mathcal{O}_S)$ 上的模.
证明 $j(e_{i\ast}(Q_i))$ 为同构.
只须看出序列
\begin{equation*}
0\to Q_i\to\varepsilon_{Z^{\bigdot}\ast}(e_{i\ast}(Q_i))^0\to
\varepsilon_{Z^{\bigdot}\ast}(e_{i\ast}(Q_i))^1\to\ldots
\tag{2.3.5.1}
\end{equation*}\etiquette[2.3.5.1]{Vbis.eq2.3.5.1}
正合.

为此, 只须证明对 $\Mod(S,\mathcal{O}_S)$ 的每个对象 $X$, 序列
\begin{equation*}
0\to\Hom(X,Q_i)\to
\Hom(X,\varepsilon_{Z^{\bigdot}\ast}(e_{i\ast}(Q_i))^0)\to
\Hom(X,\varepsilon_{Z^{\bigdot}\ast}(e_{i\ast}(Q_i))^1)\to\ldots
\tag{2.3.5.2}
\end{equation*}\etiquette[2.3.5.2]{Vbis.eq2.3.5.2}
正合. 直接计算表明, \eqhyperref[Vbis.eq2.3.5.2]{2.3.5.2} 归结为
\hypertarget{sga4.exp.Vbis.equation.2.3.5.3}{}\begin{gather*}
0\to\Hom(\ZZ,\Hom(X,Q_i))\to
\Hom(\Hom_{\Add(D)}(Z^0,[i]),\Hom(X,Q_i))\to\tag{2.3.5.3}\\
\Hom(\Hom_{\Add(D)}(Z^1,[i]),\Hom(X,Q_i))\to\ldots{\quoi}.
\end{gather*}
\sisi{}{由假设, 序列
$$
0\leftarrow\ZZ\leftarrow\Hom_{\Add(D)}(Z^0,[i])
\leftarrow\Hom_{\Add(D)}(Z^1,[i])\leftarrow\ldots
$$
是自由模的正合列, 因而零伦.
由此得 \hyperlink{sga4.exp.Vbis.equation.2.3.5.3}{(2.3.5.3)} 正合.}

\sisi{$+$}{} 若对 $D$ 的每个对象 $i$, $Q_i$ 为完全无上同调的模,
则对各复形 \eqhyperref[Vbis.eq2.3.5.1]{2.3.5.1} 应用函子 $\prod_{i\in\Ob(D)}$,
由 (\hyperref[Vbis.lem1.3.10.1]{1.3.10.1}) 可见所得复形仍无上同调.
故 $j(\prod_{i\in\Ob(D)}e_{i\ast}(Q_i))$ 为同构.

\sisi{$+$}{} 请读者推出, 若对每个 $n$, 有 $F^n\in I_{(S\times D,\mathcal{O}_S)}$,
则 $j(F^{\bigdot})$ 为同构 (例如可用谱序列).

\begin{enonce*}{推论 2.3.6}\etiquette[2.3.6]{Vbis.cor2.3.6}
在\pageoriginale{} (\hyperref[Vbis.prop2.3.5]{2.3.5}) 的条件下,
态射 $j$ 将 $\sisi{\RR^+}{R}\varepsilon_{Z^{\bigdot}\ast}$
等同于 $\varepsilon_\ast$ 的导出函子.
\end{enonce*}

由 (\hyperref[Vbis.prop1.3.10]{1.3.10}) 及 (\hyperref[Vbis.prop2.3.5]{2.3.5}),
$K^{+}(\Ysheaf{\Gamma}(S\times D),\mathcal{O}_S)$ 中
由 $I_{(S\times D,\mathcal{O}_S)}$ 的对象组成的复形所定义的子范畴,
对于函子 $K^{+}(\varepsilon_\ast)$ 满足
[\cite{sga4.exp.Vbis.bib.7}, I] 的定理 (\hyperref[Vbis.subsec5.1]{5.1}) 的条件.
推论立即由此观察得出.

\begin{enonce*}{推论 2.3.7}\etiquette[2.3.7]{Vbis.cor2.3.7}
在 (\hyperref[Vbis.prop2.3.5]{2.3.5}) 的条件下,
设 $(E,A)$ 为环化 $D$-拓扑斯,
$\theta:(E,A)\to(S\times D,\mathcal{O}_S)$ 为增广. 存在典范同构
$$
\sisi{\RR^+}{R}\bar\theta_\ast\xrightarrow{\sim}
\sisi{\RR^+}{R}\varepsilon_{Z^{\bigdot}\ast}\circ
\sisi{\RR^+}{R}\Gamma(\theta_\ast){\quoi}.
$$
\end{enonce*}

这来自 [\cite{sga4.exp.Vbis.bib.7}, I, 5.4].

\begin{enonce*}[remark]{注 2.3.8}\etiquette[2.3.8]{Vbis.rem2.3.8}
我们先验地知道 $\sisi{\RR^+}{R}\varepsilon_\ast$ 存在,
且存在同构 $\sisi{\RR^+}{R}\bar\theta_\ast\xrightarrow{\sim}
\sisi{\RR^+}{R}\varepsilon_\ast\circ\sisi{\RR^+}{R}\Gamma(\theta_\ast)$,
因为 $\Gamma(\theta_\ast)$ 和 $\varepsilon_\ast$ 由环化拓扑斯态射诱导
(参见 (\hyperref[Vbis.subsubsec1.3.6]{1.3.6}) 及
(\hyperref[Vbis.subsubsec2.1.1]{2.1.1})).
尽管如此, 上述计算将很有用, 并可用于其他背景,
因为实际中 (\hyperref[Vbis.prop2.3.5]{2.3.5}) 的条件总会满足.
\end{enonce*}

现在给出可应用 (\hyperref[Vbis.prop2.3.5]{2.3.5}) 的基本例子.

\setcounter{subsubsection}{8}

% END INCLUDED SGA4_English_current_buildable_source_20260906/source/05bis/drafts/SGA4_VBIS_LINES_1074_1378_English.texfrag


% BEGIN INCLUDED SGA4_English_current_buildable_source_20260906/source/05bis/drafts/SGA4_VBIS_LINES_1379_1421_English.texfrag
\subsubsection{}\etiquette[2.3.9]{Vbis.subsubsec2.3.9}
在 $\Add(\Delta)$ 中, 令 $Z^n=[n]$, 并由公式
\begin{equation*}
d^n=\sum_{i=0}^{n+1}(-1)^i\partial^i.\tag{2.3.9.1}
\end{equation*}\etiquette[2.3.9.1]{Vbis.eq2.3.9.1}
定义 $d^n:Z^n\to Z^{n+1}$.
这里 $\partial^i:[n]\to[n+1]$ 为第 $i$ 个面
[参见 \cite{sga4.exp.Vbis.bib.3} II 2].
取显然的自然增广作为 $t:Z^0\to\ZZ$
\nde{明确地说, 若 $F$ 为微分为 $\mathrm d_F$ 的单纯复形,
则 $\varepsilon_{Z^{\bigdot}*}(F)$ 为双复形 $F^j(i)$ 的相伴单复形,
其水平微分为单纯微分 $\mathrm d^i$, 垂直微分为 $\mathrm d_F^j$.}.

\hyperref[Vbis.subsubsec2.3.4]{2.3.4} 的条件显然满足.
\hyperref[Vbis.prop2.3.5]{2.3.5} 的条件来自 \cite{sga4.exp.Vbis.bib.5} I (3.7.4).

\subsubsection{}\etiquette[2.3.10]{Vbis.subsubsec2.3.10}
在\pageoriginale{} $\Add(\Delta\times\Delta)$ 中,
考虑下列双复形的相伴单复形
\hypertarget{sga4.exp.Vbis.equation.2.3.10.1}{}\[\vcenter{
\xymatrix@C=2cm@R=1.75cm{
\big([n],[m+1]\big)  \ar[r]^{s_1^{n,m+1}} & \big([n+1],[m+1]\big)\\
\big([n],[m]\big) \ar[r]^{s_1^{n,m}} \ar[u]^{s_2^{n,m}} &
\big([n+1],[m]\big)\ar[u]_{s_2^{n+1,m}}
}}\leqno{(2.3.10.1)}
\]
其中
\begin{equation*}
\tag{2.3.10.2}
\begin{cases}
S^{n,m}_1=\displaystyle\mathop{\sum}_{i=0}^{n+1}(-1)^i
(\partial^i,\id_{[m]})\,,\\[15pt]
S^{n,m}_2=\displaystyle\mathop{\sum}_{j=0}^{m+1}(-1)^j
(\id_{[n]},\partial^j){\quoi}.
\end{cases}
\end{equation*}\etiquette[2.3.10.2]{Vbis.eq2.3.10.2}

读者可利用 \hyperref[Vbis.subsubsec2.3.9]{2.3.9} 验证,
这样定义的复形配备显然的自然增广后, 满足 \hyperref[Vbis.prop2.3.5]{2.3.5} 的条件.

\subsubsection{}\etiquette[2.3.11]{Vbis.subsubsec2.3.11}
多重单纯情形的处理类似
(取 $\underbrace{\Delta\times\Delta\times\cdots\times\Delta}_{p\text{ 次}}$).


% END INCLUDED SGA4_English_current_buildable_source_20260906/source/05bis/drafts/SGA4_VBIS_LINES_1379_1421_English.texfrag


% BEGIN INCLUDED SGA4_English_current_buildable_source_20260906/source/05bis/drafts/SGA4_VBIS_LINES_1422_1690_English.texfrag
\subsection{相对上同调下降}\etiquette[2.4]{Vbis.subsec2.4}\pageoriginale

\setcounter{subsubsection}{-1}
\subsubsection{}\etiquette[2.4.0]{Vbis.subsubsec2.4.0}
设 $\Ee\to B$ 为拓扑斯对偶双纤维化的范畴,
$\mathcal O$ 为 $\Ysheaf{\Gamma}(\Ee^{\circ})$ 的环.

这些资料典范定义一个在 $B^{\circ}$ 上纤维化且余纤维化的阿贝尔范畴系统,
记为 $\Ysheaf{\Mod}(\Ee^{\circ},\mathcal O)$
(参见 (\hyperref[Vbis.rem1.3.3]{1.3.3}) 及
(\hyperref[Vbis.subsubsec1.3.4]{1.3.4})).
下文始终相对于纤维且余纤维函子
$\Ysheaf{\Mod}(\mathcal E^{\circ},\mathcal O)\to B^{\circ}$
取“正像”与“逆像”函子.

\begin{enonce*}{定义 2.4.0.1}\etiquette[2.4.0.1]{Vbis.def2.4.0.1}
若 $f_{\ast}:\Mod(\mathcal E^{\circ}_S,\mathcal O_S)\to
\Mod(\mathcal E^{\circ}_T,\mathcal O_T)$ 为正合函子,
就称 $B$ 中的态射 $f:T\to S$ (相对于 $(\mathcal E,\mathcal O)$) \emph{平坦}
\nde{此记号可能引起混淆: 或许应写为 $(f^\circ)_\star$,
它在形式上等于 $(f^\star)^\circ$,
其中 $f^\star$ 为与 $\Ee\rightarrow B$ 相联系的纤维函子.
例如, 对拓扑空间上的层构成的、在拓扑空间范畴上纤维化的范畴
(参见 \hyperref[Vbis.subsec4.1]{4.1}), 有
$(f^\circ)_\star(S,F)=(T,f^{-1}F)=(f^\star)^\circ(S,F)$,
其中 $f:T\rightarrow S$ 为拓扑空间的连续映射, $F$ 为 $S$ 上的层,
$f^{-1}$ 为通常的逆像. 此时便认出通常的平坦性条件.}.
\end{enonce*}

\subsubsection{}\etiquette[2.4.1]{Vbis.subsubsec2.4.1}
设 $G$ 为 $\Ysheaf{\Mod}(\mathcal E^{\circ},\mathcal O)$ 的纤维且余纤维子范畴,
使对 $B$ 的每个对象 $S$, $G_S$ 为
$(\Ysheaf{\Mod}(\mathcal E^{\circ},\mathcal O))_S$ 的\emph{厚}子范畴
(参见 \sisi{Tohokû}{Tôhoku} (1.11)).
读者将在 \hyperref[Vbis.sec4]{第 4 节} 找到此类情形的例子.

若 $S$ 为 $B$ 的对象, 用 $D^+(S)$ 表示 $\Mod(\mathcal E^{\circ}_S,\mathcal O_S)$ 的导出范畴.
按照 [\cite{sga4.exp.Vbis.bib.7} (I \S\ 4)], 引入范畴 $D^+_{G_S}(S)$,
它等同于 $D^+(S)$ 中由满足对每个 $n$ 均有 $H^n(X^{\bigdot})\in G_S$
的对象 $X^{\bigdot}$ 组成的满子范畴.

\emph{始终假定下述条件成立.}

\paragraph{}\etiquette[2.4.1.1]{Vbis.par2.4.1.1}
\emph{对 $B$ 中的每个态射 $h:T\to S$,
函子 $\sisi{\RR^+}{Rh^{\ast}}:D^+(T)\to D^+(S)$ 具有如下性质:
只要 $F\in G_T$, 就有 $\sisi{\RR^+}{Rh^{\ast}}(F)\in D^+_{G_S}(S)$}
\nde{如上所述, 与其写 $Rh^\ast$, 更正确的写法为
$R(h^\circ)^\ast=R(h_\ast)^\circ$,
这使记号与通常的几何直觉更一致.
同样也可以相对于 $\Mod(\mathcal E,\mathcal O)\rightarrow B$
而非其反范畴来取纤维和余纤维函子.
但此时纤维只能是拓扑斯的对偶.}.

由 [\cite{sga4.exp.Vbis.bib.7} (I. (7.3). (ii))], 此条件蕴含
$\sisi{\RR^+}{Rh^{\ast}}(D^+_{G_T}(T))\subset D^+_{G_S}(S)$.

\subsubsection{}\etiquette[2.4.2]{Vbis.susubsec2.4.2}
设\pageoriginale{} $D$ 为 $\Uu$-小范畴.
\emph{下文始终假定给定 $\Add(D)$ 的一个复形 $Z^{\bigdot}$,
满足} (\hyperref[Vbis.prop2.3.5]{2.3.5}) 的条件
(实际应用中将有 $D=\Delta$ 或 $D=\Delta\times\Delta$).

若 $X:D^{\circ}\to B$ 为 $B$ 的 $D$-对象,
用 $\Mod_G(\Ysheaf{\Gamma}(\bar X),\mathcal O)$
表示 $\Mod(\Ysheaf{\Gamma}(\bar X),\mathcal O)$ 中由满足下述条件的对象 $F$ 组成的满子范畴:
对 $D$ 的每个对象 $i$, 有 $e_i^{\ast}(F)\in G_i$ $(=G_{X_i})$.
范畴 $\Mod_G(\Ysheaf{\Gamma}(\bar X),\mathcal O)$
为 $\Mod(\Ysheaf{\Gamma}(\bar X),\mathcal O)$ 的厚子范畴.
用 $K_G^+(\Ysheaf{\Gamma}(\bar X),\Oo)$
(分别用 $D_G^+(\Ysheaf{\Gamma}(\bar X),\Oo)$)
表示 $K^+(\Ysheaf{\Gamma}(\bar X),\Oo)$
(分别为 $D^+(\Ysheaf{\Gamma}(\bar X),\Oo)$)
中由满足对每个 $n$ 均有
$H^n(X^{\bigdot})\in\Mod_G(\Ysheaf{\Gamma}(\bar X),\Oo)$ 的对象 $X^{\bigdot}$ 组成的满子范畴.

\subsubsection{}\etiquette[2.4.3]{Vbis.subsubsec2.4.3}
设 $X$ 和 $X'$ 为 $B$ 的两个 $D$-对象, $\alpha:X\to X'$ 为函子性态射.
回顾 (参见 (\hyperref[Vbis.subsubsec1.3.4]{1.3.4})),
相应的环化 $D$-拓扑斯态射仍记为 $\alpha:(\bar X,\mathcal O)\to(\bar X',\mathcal O)$.

\begin{enonce*}{定义 2.4.3.1}\etiquette[2.4.3.1]{Vbis.def2.4.3.1}
若相应的环化 $D$-拓扑斯态射在 (\hyperref[Vbis.par1.3.6.1]{1.3.6.1}) 的意义下平坦,
就称态射 $\alpha:X\to X'$ \emph{平坦}.
\end{enonce*}

于是有显然的引理:

\begin{enonce*}{引理 2.4.3.2}\etiquette[2.4.3.2]{Vbis.lem2.4.3.2}
态射 $\alpha:X\to X'$ 平坦当且仅当对 $D$ 的每个对象 $i$,
$\alpha_i:X_i\to X'_i$ 为平坦态射
\emph{(参见 (\hyperref[Vbis.def2.4.0.1]{2.4.0.1}))}.
\end{enonce*}

\subsubsection{}\etiquette[2.4.4]{Vbis.subsubsec2.4.4}
态射 $\alpha:X\to X'$ 定义函子
\begin{equation*}
K^+(\Gamma(\alpha_*)):K^+(\Ysheaf{\Gamma}(\bar X),\Oo)
\longrightarrow K^+(\Ysheaf{\Gamma}(\bar X'),\Oo){\quoi}.
\tag{2.4.4.1}
\end{equation*}\etiquette[2.4.4.1]{Vbis.eq2.4.4.1}

\begin{enonce*}{引理 2.4.4.2}\etiquette[2.4.4.2]{Vbis.lem2.4.4.2}
函子
$$
K^+(\Gamma(\alpha_*))|_{K_G^+(\Ysheaf{\Gamma}(\bar X),\Oo)}:
K_G^+(\Ysheaf{\Gamma}(\bar X),\Oo)\longrightarrow
K^+(\Ysheaf{\Gamma}(\bar X'),\Oo)
$$
有右导出函子, 记为 $\sisi{\RR^+}{R}_G\Gamma(\alpha_{\ast})$,
且典范态射
$$
\sisi{\RR^+}{R}_G\Gamma(\alpha_{\ast})\longrightarrow
\sisi{\RR^+}{R}\Gamma(\alpha_{\ast})|_{D_G^+
(\Ysheaf{\Gamma}(\bar X),\Oo)}
$$
为\pageoriginale{}同构.

此外,
$$
\sisi{\RR^+}{R}_G\Gamma(\alpha_{\ast})
(D_G^+(\Ysheaf{\Gamma}(\bar X),\Oo))
\subset D_G^+(\Ysheaf{\Gamma}(\bar X'),\Oo).
$$
\end{enonce*}

第一部分来自 (\hyperref[Vbis.prop1.3.10]{1.3.10}) 及
[\cite{sga4.exp.Vbis.bib.7}, (I. (5.2))].
为验证最后一个断言, 由 [\cite{sga4.exp.Vbis.bib.7}, (I. (7.3). (ii))],
只须证明只要 $F\in\Mod_G(\Ysheaf{\Gamma}(\bar X),\Oo)$, 就有
$\sisi{\RR^+}{R}\Gamma(\alpha_{\ast})(F)\in D_G^+(\Ysheaf{\Gamma}(\bar X'),\Oo)$.
这来自 \hyperref[Vbis.cor1.3.11]{(1.3.1.1)} 给出的
$\sisi{\RR^+}{R}\Gamma(\alpha_\ast)$ 的显式计算,
以及假设 (\hyperref[Vbis.par2.4.1.1]{2.4.1.1}).

\subsubsection{}\etiquette[2.4.5]{Vbis.subsubsec2.4.5}
设 $S$ 为 $B$ 的对象. 环化 $D$-拓扑斯
$(\bar{C^D_S},\mathcal O)$ (参见
(\hyperref[Vbis.subsubsec1.2.6]{1.2.6})) 等同于常值环化 $D$-拓扑斯
$(\mathcal E^{\circ}_S\times D,\mathcal O_S)$,
并有下述引理, 它来自 [\cite{sga4.exp.Vbis.bib.7} (I. (7.3). (ii))]:

\begin{enonce*}{引理 2.4.5.1}\etiquette[2.4.5.1]{Vbis.lem2.4.5.1}
函子 $\sisi{\RR^+}{R}\varepsilon_{Z^{\bigdot}\ast}|_{D_G^+(\Ysheaf{\Gamma}(\bar{C^D_S}),\mathcal O)}$,
记为 $\sisi{\RR^+}{R}_G\varepsilon_{Z^{\bigdot}\ast}$,
取值于 $D^+_{G_S}(S)$, 并等同于下述函子的一个右导出函子:
$$
K^+(\varepsilon_{\ast})|_{K_G^+
(\Ysheaf{\Gamma}(\bar{C^D_S}),\mathcal O)}:
K_G^+(\Ysheaf{\Gamma}(\bar{C^D_S}),\mathcal O)
\longrightarrow K^+(S){\quoi}.
$$
\end{enonce*}

\begin{enonce*}{命题 2.4.6}\etiquette[2.4.6]{Vbis.prop2.4.6}
设 $\theta:X\to C^D_S$ 为 $B$ 的增广 $D$-对象.
函子 $K^+(\bar\theta_{\ast})|_{K_G^+(\Ysheaf{\Gamma}(\bar X),\mathcal O)}$
有取值于 $D^+_{G_S}(S)$ 的右导出函子, 记为 $\sisi{\RR^+}{R}_G\bar\theta_{\ast}$.
此外, 存在典范同构
$$
\sisi{\RR^+}{R}_G\bar\theta_{\ast}\xrightarrow{\sim}
\sisi{\RR^+}{R}_G\varepsilon_{Z^{\bigdot}\ast}\circ
\sisi{\RR^+}{R}_G\Gamma(\theta_{\ast}){\quoi}.
$$
\end{enonce*}

由上述讨论得出证明, 留待读者完成.

现在可以给出有关上同调下降的定义.

\setcounter{subsubsection}{6}
\subsubsection{}\etiquette[2.4.7]{Vbis.subsubsec2.4.7}
设\pageoriginale{} $\theta:X\to C^D_S$ 为 $B$ 的增广 $D$-对象.
假定 $\theta$ 平坦, 于是函子 $\bar\theta^{\ast}$ 在 $D^+_{G_S}(S)$ 上的限制
取值于 $D_G^+(\Ysheaf{\Gamma}(\bar X),\mathcal O)$;
将此限制记为 $\bar\theta_G^{\ast}$. 采用此记号:

\begin{enonce*}{引理 2.4.7.1}\etiquette[2.4.7.1]{Vbis.lem2.4.7.1}
存在两个函子性态射
$\alpha:\id_{D^+_{G_S}(S)}\to\sisi{\RR^+}{R}_G\bar\theta_{\ast}\circ\bar\theta_G^{\ast}$
以及
$\beta:\bar\theta_G^{\ast}\circ\sisi{\RR^+}{R}_G\bar\theta_{\ast}\to\id_{D_G^+(\Ysheaf{\Gamma}(\bar X),\mathcal O)}$,
使这两个函子互为伴随.
\end{enonce*}

这立即来自 (\hyperref[Vbis.lem2.2.3.1]{2.2.3.1})
和 (\hyperref[Vbis.prop2.4.6]{2.4.6}).

\begin{enonce*}{定义 2.4.8}\etiquette[2.4.8]{Vbis.def2.4.8}
若
\begin{enumerate}
\renewcommand{\theenumi}{\arabic{enumi}}
\renewcommand{\labelenumi}{\theenumi$^{\circ}$)}
\item $\theta$ 平坦;

\item 函子 $\bar\theta_G^{\ast}$ 全忠实,
\end{enumerate}
就称态射 $\theta$ 为相对于 $G$ 的上同调 $1$-下降增广.
\end{enonce*}

\begin{enonce*}{注 2.4.9}\etiquette[2.4.9]{Vbis.rem2.4.9}
前一定义的条件 $2^{\circ})$ 等价地可表述为:
(\hyperref[Vbis.lem2.4.7.1]{2.4.7.1}) 中的态射 $\alpha$ 为同构.
\end{enonce*}

\begin{enonce*}{定义 2.4.10}\etiquette[2.4.10]{Vbis.def2.4.10}
设 $\theta:X\to C^D_S$ 为 $B$ 的增广 $D$-对象.
若函子
$$
\Gamma(\theta^{\ast})\circ\varepsilon^{\ast}:G_S\longrightarrow
\Ysheaf{\Gamma}^{\Cocart}(D,G)
$$
为范畴等价, 就称态射 $\theta$ 为\emph{相对于 $G$ 的有效下降增广}.
\end{enonce*}

\begin{enonce*}{定义 2.4.11}\etiquette[2.4.11]{Vbis.def2.4.11}
若 $\theta$ 既是上同调 $1$-下降增广,
又是相对于 $G$ 的有效下降增广,
就称态射 $\theta$ 为\emph{相对于 $G$ 的上同调 $2$-下降增广}
(或\emph{有效上同调下降增广}).
\end{enonce*}

(\hyperref[Vbis.prop2.2.7]{2.2.7}) 的证明遂给出下述结果:

\begin{enonce*}{定理 2.4.12}\etiquette[2.4.12]{Vbis.thm2.4.12}
设\pageoriginale{} $\theta:X\to C^D_S$ 为相对于 $G$ 的上同调 $2$-下降增广.
则 $\bar\theta_G^{\ast}$ 的本质像是
$D_G^+(\Ysheaf{\Gamma}(\bar X),\mathcal O)$ 中由满足下述条件的复形 $F^{\bigdot}$ 组成的满子范畴:
对每个 $n$, $H^n(F^{\bigdot})$ 为 $G$ 的余笛卡尔截面.
\end{enonce*}

下述命题将定义 (\hyperref[Vbis.def2.4.10]{2.4.10})
转述为 (\cite{sga4.exp.Vbis.bib.4}, 6) 的语言, 其验证留待读者完成:

\begin{enonce*}{命题 2.4.13}\etiquette[2.4.13]{Vbis.prop2.4.13}
态射 $\theta:X\to C^D_S$ 为相对于 $G$ 的有效下降增广,
当且仅当由 $\theta$ 定义的 $B$ 上函子 $D^{\circ}\to B/S$
在 $G^{\circ}$ 于 $B/S$ 上的笛卡尔截面范畴与 $G^{\circ}$ 于 $D^{\circ}$ 上的笛卡尔截面范畴之间诱导等价.
\end{enonce*}

\begin{enonce*}{推论 2.4.14}\etiquette[2.4.14]{Vbis.cor2.4.14}
假定有限纤维积在 $B$ 中可表, 设 $f:R\to S$ 为态射.
增广半单纯对象 $[R|_fS]$ \emph{(参见 (\hyperref[Vbis.subsubsec1.2.7]{1.2.7}))}
为相对于 $G$ 的有效下降对象,
当且仅当 $f$ 为 $B$ 上纤维范畴 $G^{\circ}$ 的有效下降态射.
\end{enonce*}

这来自 \hyperref[Vbis.prop2.4.13]{(2.4.13)} 及 (\cite{sga4.exp.Vbis.bib.4}, 9).

\begin{enonce*}{定义 2.4.15}\etiquette[2.4.15]{Vbis.def2.4.15}
假定有限纤维积在 $B$ 中可表, 设 $f:R\to S$ 为态射.
若增广半单纯对象 $[R|_fS]$ 为相对于 $G$ 的上同调 $i$-下降对象,
就称 $f$ 为相对于 $G$ 的上同调 $i$-下降态射 $(i=1,2)$.
\end{enonce*}

\begin{enonce*}{定义 2.4.16}\etiquette[2.4.16]{Vbis.def2.4.16}
假定有限纤维积在 $B$ 中可表, 设 $\theta:X\to C^D_S$ 为增广.
若对每个态射 $T\to S$, 经基变换\pageoriginale{}得到的增广
$X_T\xrightarrow{\theta_T}C^D_T$ 均为相对于 $G$ 的上同调 $i$-下降增广,
就称 $\theta$ 为\emph{相对于 $G$ 的普遍上同调 $i$-下降增广 $(i=1,2)$}.
\end{enonce*}

\emph{相对于 $G$ 的普遍上同调 $i$-下降态射}的概念立即来自前两个定义.

我们将常用“$G$-上同调 $i$-下降增广”这一术语,
代替“相对于 $G$ 的上同调 $i$-下降增广”.

\subsection{下降谱序列}\etiquette[2.5]{Vbis.subsec2.5}


% END INCLUDED SGA4_English_current_buildable_source_20260906/source/05bis/drafts/SGA4_VBIS_LINES_1422_1690_English.texfrag


% BEGIN INCLUDED SGA4_English_current_buildable_source_20260906/source/05bis/drafts/SGA4_VBIS_LINES_1691_1929_English.texfrag

\subsubsection{}\etiquette[2.5.1]{Vbis.subsubsec2.5.1}
设 $A$ 为阿贝尔范畴. 用 $F(A)$ 表示如下范畴:
其对象为配备离散且余离散滤过的 $A$ 的对象, 其态射为 $A$ 的滤过态射;
范畴 $F(A)$ 为加性范畴 (但并非阿贝尔范畴).

$A$ 的滤过逐次离散且余离散的滤过复形范畴 $K(F(A))$ 是三角范畴,
且典范函子
\begin{equation*}
\Gr_n : K(F(A)) \longrightarrow K(A)\tag{2.5.1.1}
\end{equation*}\etiquette[2.5.1.1]{Vbis.eq2.5.1.1}
是三角函子. 因而 $K(F(A))$ 中使对每个 $n$, $\Gr_n(f)$ 都为拟同构的态射 $f$
所成的集合 $\Sigma$ 是饱和乘法系
[参见 \cite{sga4.exp.Vbis.bib.9}. \S~2. 第 ${\rm n^\circ}$ 1 节].

将此系中的箭头倒置, 就得到新的三角范畴, 记为 $D F(A)$,
且各函子 $\Gr_n$ 扩张为函子
\begin{equation*}
\Gr_n : D F(A) \longrightarrow D(A){\quoi}.\tag{2.5.1.2}
\end{equation*}\etiquette[2.5.1.2]{Vbis.eq2.5.1.2}

下文只考虑下有界复形: 这自然引出记号 $K^+(F(A))$ 与 $D^+ F(A)$.

\subsubsection{}\etiquette[2.5.2]{Vbis.subsubsec2.5.2}
设\pageoriginale{} $B$ 为 $A$ 的厚子范畴,
用 $K^+_B(F(A))$ 表示 $K^+(F(A))$ 中由满足对每对整数 $(i, j)$,
$H^{i}(\Gr_j(X^{\bigdot})) \in B$ 的复形 $X^{\bigdot}$ 组成的满子范畴:
$K_B^+(F(A))$ 是 $K(F(A))$ 的\emph{局部化三角子范畴}
[参见 \cite{sga4.exp.Vbis.bib.7}. (I. \S\ 5)].
用 $D_B^+ F(A)$ 表示 $D F(A)$ 中由满足对每对 $(i, j)$,
$H^i(\Gr_j(X^{\bigdot})) \in B$ 的下有界复形 $X^{\bigdot}$ 组成的满子范畴:
由 (\loccit), $D^+_B F(A)$ 等同于分式范畴
$K^+_B(F(A))_{\Sigma \cap K^{+}_{B} (F(A))}$.

“遗忘滤过函子”
$$
\iota : K^+(F(A)) \longrightarrow K^+(A)
$$
为三角函子. 此外, 由 [\cite{sga4.exp.Vbis.bib.5}. I. (4.7)] 得出,
若 $f \in \Sigma$, 则 $\iota(f)$ 为拟同构,
且 $\iota(K^+_B(F(A))) \subset K^+_B(A)$.
因此函子 $\iota$ 扩张为三角函子
\begin{equation*}
\iota : D^+ F(A) \longrightarrow D^+ (A) \tag{2.5.2.1}
\end{equation*}\etiquette[2.5.2.1]{Vbis.eq2.5.2.1}
使 $\iota (D^+_B F(A)) \subset D^+_B (A)$.

最后, 注意存在典范\emph{谱函子} $D^+_B F(A) \to B$,
收敛到 $H^n \circ \iota$, 其 $E_1$ 项由 $H^{p+q} \circ \Gr_{p}$ 给出
[参见 \cite{sga4.exp.Vbis.bib.5} I. (4.2)].

\emph{现回到} (2.3) \emph{的记号, 并假定} $D = \Delta$:
$Z^{\bigdot}$ 表示由 \eqhyperref[Vbis.eq2.3.9.1]{2.3.9.1} 定义的 $\Add(\Delta)$ 的复形.

\begin{enonce*}{命题 2.5.3}\etiquette[2.5.3]{Ii.prop2.5.3}
设 $(S, \mathcal{O}_S)$ 为环化拓扑斯. 函子
$$\sisi{\RR^+}{R}
\varepsilon_{Z^{\bigdot} \ast} : D^{+} (\Ysheaf{\Gamma} (S \times D),
\mathcal{O}_S) \to D^{+} (S, \mathcal{O}_S)$$ 有典范分解
$$
D^{+} (\Gamma (S \times D), \mathcal{O}_S) \xrightarrow{F^{\sisi{\RR^+}{R}
\varepsilon} Z^{\bigdot}\ast} D^{+} F(S, \mathcal{O}_S)
\xrightarrow{\iota} D^{+} (S, \mathcal{O}_S){\quoi}.
$$
此外, 对每个整数 $i$, 下图交换:
\[
\xymatrix@=2.5cm{
D^{+}(\Ysheaf{\Gamma}(S\times D),\Oo_S)
\ar[r]^(.55){\sisi{\RR^{+}}{R}e^{\ast}_i} \ar[d]_{F^{\sisi{\RR^{+}}{R}\epsilon_{Z^{\bigdot}\ast}}}& D^{+}(S,\Oo_S)\\
D^{+}F(S,\Oo_S) \ar[ru]_{\Gr_i}& .
}
\]
\end{enonce*}

考虑\pageoriginale{}函子 $C^+ (C^+(S, \mathcal{O}_S))\xrightarrow{(\quad)_s} C^+(S, \mathcal{O}_S)$:
若 $X^{\cdot\cdot}$ 为 $C^+ (C^+(S, \mathcal{O}_S))$ 的对象,
则为 $(X^{\cdot\cdot})_s$ 配备第二典范滤过
[参见 \cite{sga4.exp.Vbis.bib.5}. I.4.8]; 这给出分解
\begin{equation*}
(\quad)_s : C^+ (C^+(S, \mathcal{O}_S)) \longrightarrow C^+ (F(S,
  \mathcal{O}_S)) \xrightarrow{\iota} C^+(S, \mathcal{O}_S)\tag{2.5.3.1}
\end{equation*}\etiquette[2.5.3.1]{Vbis.eq2.5.3.1}
它过渡到范畴 $K^+$:
\begin{equation*}
(\quad)_s : K^+ (C^+(S, \mathcal{O}_S)) \longrightarrow K^+ (F(S,
  \mathcal{O}_S)) \xrightarrow{\iota} K^+(S, \mathcal{O}_S)\tag{2.5.3.2}
\end{equation*}\etiquette[2.5.3.2]{Vbis.eq2.5.3.2}
因为 $C^+ (C^+(S, \mathcal{O}_S))$ 中的同伦诱导滤过同伦.

于是有图
\[\vcenter{
\xymatrix{
D^{+}(\Ysheaf{\Gamma}(S\times\Delta),\Oo_S) \ar[rr]^{R\epsilon_{Z^{\bigdot}}\ast} \ar@{-->}[rd]&                   &D^{+}(S,\Oo_S)\\
                                              & D^{+}F(S,\Oo_S)
\ar[ru]_{\imath}&\\
                                  & K^{+}(F(S,\Oo_S))
\ar[u] \ar[rd]^{\imath}&\\
K^{+}(\Ysheaf{\Gamma}(S\times \Delta),\Oo_S) \ar[uuu]
\ar[rr]_{(\;)_s\circ K^{+}(\epsilon_{Z^{\bigdot}}\ast)} \ar[ru]& & K^{+}(S,\Oo_S)\ar[uuu]
}}\leqno{(2.5.3.3)}
\]\etiquette[2.5.3.3]{Vbis.eq2.5.3.3}
且可验存在唯一函子
$$
{}_F \sisi{\RR^{+}}{R} \epsilon_{Z^{\bigdot} \ast} : D^{+}(\Ysheaf{\Gamma}(S \times
\Delta), \mathcal{O}_S) \longrightarrow D^{+}F(S, \mathcal{O}_S)
$$
使\pageoriginale{} \eqhyperref[Vbis.eq2.5.3.3]{2.5.3.3} 交换.

命题的其余部分显然成立.

\begin{enonce*}{命题 2.5.4}\etiquette[2.5.4]{Vbis.prop2.5.4}
设 $(E, A)$ 为半单纯环化拓扑斯,
$$\theta : (E, A) \to (S \times \Delta,
\mathcal{O}_S)$$ 为增广.
对每个整数 $i$, 用
$$\theta_i : (E_i, A_i) \to (S, \mathcal{O}_S)$$
表示由 $\theta$ 在 $i$ 上诱导的环化拓扑斯态射.

函子 $\sisi{\RR^+}{R} \bar{\theta}_\ast : D^+ (\Ysheaf{\Gamma}(E), A)\to D^+ (S, \mathcal{O}_S)$
有典范分解
$$
D^+ (\Ysheaf{\Gamma}(E), A) \xrightarrow{{}_F\sisi{\RR^{+}}{R} \bar{\theta}_{\ast}}
D^{+} F(S, \mathcal{O}_S) \xrightarrow{\iota} D^+ (S, \mathcal{O}_S)
$$
使对每个整数 $i$, 典范图
\[
\xymatrix@C=2.5cm@R=1.75cm{
D^{+}(\Ysheaf{\Gamma}(E),A)  \ar[r]^{\sisi{\RR^{+}}{R}e_i^{\ast}}
\ar[d]_{{}_F \sisi{\RR^{+}}{R}
\bar{\theta}_\ast}&
D^{+}(E_i,A_i) \ar[d]^{\sisi{\RR^{+}}{R}\theta_{i\ast}}\\
D^{+}F(S,\Oo_S)  \ar[r]_{\Gr_i} & D^{+}(S,\Oo_S)\\
}
\]
本质交换.
\end{enonce*}

这来自上述讨论及 \hyperref[Vbis.cor1.3.12]{(1.3.1.2)}.

\begin{enonce*}{命题 2.5.5}\etiquette[2.5.5]{Vbis.prop2.5.5}
设 $(E, A)$ 为半单纯环化拓扑斯,
$$\theta : (E, A) \to (S \times \Delta, \mathcal{O}_S)$$ 为上同调 $1$-下降增广.
设
$$H : (S, \mathcal{O}_S) \to (R, \mathcal{O}_R)$$ 为环化拓扑斯态射.
则有从 $D^+(S, \mathcal{O}_S)$ 到 $\Mod(R, \mathcal{O}_R)$ 的谱函子:
$$
E^{pq}_1 = R^q(H \circ \theta_p)_{\ast} \circ
\theta^{\ast}_p \longrightarrow R^{p+q}H_{\ast}
$$
称为\pageoriginale\emph{下降谱函子}.
\end{enonce*}

设 $H \circ \theta : (E, A) \to (R \times \Delta, \mathcal{O}_R)$
为由 $\theta$ 与 $H$ 复合而典范得到的增广; 有
\begin{equation*}
\bar{(H \circ \theta)}_{\ast} = H_{\ast} \circ \bar{\theta}_{\ast}\tag{2.5.5.1}
\end{equation*}\etiquette[2.5.5.1]{Vbis.eq2.5.5.1}
从而图
\[\vcenter{
\xymatrix@C=2.25cm@R=2.2cm{
D^{+}(S,\Oo_S) \ar[r]^{\sisi{\LL^{+}(\bar{\theta}^{\ast})}{\bar{\theta}^{\ast}}} \ar[rd]_{\sisi{\RR^+}{R}H_{\ast}} &
D^{+}(\Ysheaf{\Gamma}(E),A)
\ar[d]^{\sisi{\RR^{+}}{R}(\overline{H\circ
\theta})_\ast}\\
                  & D^{+}(R,\Oo_R)
}}\leqno{(2.5.5.2)}
\]\etiquette[2.5.5.2]{Vbis.eq2.5.5.2}
本质交换, 因为
$$
\sisi{\RR^+}{R} \bar{(H \circ \theta)}_{\ast} \xrightarrow{\sim} \sisi{\RR^+}{R}
H_\ast \circ \sisi{\RR^+}{R} \bar{\theta}_{\ast} \mbox{ 和 } \id_{D^{+}(S,
  \mathcal{O}_S)} \xrightarrow{\sim} \sisi{\RR^+}{R} \bar{\theta}_\ast \circ
\bar{\theta}^{\ast}{\quoi}.
$$
由 (\hyperref[Vbis.prop2.5.4]{2.5.4}), 对每个 $i$ 有本质交换图
\[\vcenter{
\xymatrix@C=2.25cm@R=1.75cm{
D^{+}(S,\Oo_S) \ar[r]^{\bar{\theta}^{\ast}} \ar[rdd]_{\sisi{\RR^{+}}{R}H_{\ast}}&
D^{+}(\Ysheaf{\Gamma}(E),A)
\ar[r]^{\sisi{\RR^{+}}{R}e^{\ast}_i} \ar[d]_{{}_F\sisi{\RR^{+}}{R}(\bar{H\circ\theta})_\ast}&
D^{+}(E_i,A_i) \ar[d]^{\sisi{\RR^{+}}{R}(H\circ\theta_{i})_{\ast}}\\
                  & D^{+}F(R,\Oo_R)   \ar[d]_{\imath}\ar[r]_{\Gr_i}      & D^{+}(R,\Oo_R)\\
                  & D^{+}(R,\Oo_R)          &   .
}}\leqno{(2.5.5.3)}
\]\etiquette[2.5.5.3]{Vbis.eq2.5.5.3}
所断言\pageoriginale{}的谱函子, 是将 $D^+ F(R, \mathcal{O}_R)$ 上的典范谱函子
与 ${}_F R(\bar{H \circ \theta})_{\ast}\circ\bar{\theta}^{\ast}$ 复合得到的.

下述命题留待读者验证:

\begin{enonce*}{命题 2.5.6}\etiquette[2.5.6]{Vbis.prop2.5.6}
采用前述记号, 上述谱函子的 $E^{pq}_2$ 项为
$$
\check{H}^p \circ R^q \Gamma ((H \circ \theta)_{\ast}) \circ
\bar{\theta}^\ast
$$
其中 $\check{H}^p : \Mod(\Ysheaf{\Gamma}(R \times \Delta), \mathcal{O}_R) \to\Mod(R, \mathcal{O}_R)$
为将每个函子 $\Delta \xrightarrow{F} \Mod(R, \mathcal{O}_R)$
对应到相伴复形 (在 (\hyperref[Vbis.subsubsec2.3.3]{2.3.3}) 中记为
$\varepsilon_{Z^{\bigdot}\ast} (F)^{\bigdot}$) 的第 $p$ 次同调对象的函子.
\end{enonce*}

现回到 (\hyperref[Vbis.subsec2.4]{2.4}) 中引入的术语:
由连锁命题 (\hyperref[Vbis.subsubsec2.5.2]{2.5.2}),
上述一切讨论均可在所有有意义之处以下标加上字母 $G$ 后逐字推广.
细节留待读者完成; 特别地, 得到:

\begin{enonce*}{命题 2.5.7}\etiquette[2.5.7]{Vbis.prop2.5.7}
设\pageoriginale{} $\theta : X \to C^D_S$ 为相对于 $G$ 的上同调 $1$-下降增广,
$h : S \to R$ 为 $B$ 的态射; 有从 $D^+_{G_S}(S)$ 到 $G_R$ 的谱函子:
$$
E^{pq}_1 = R^q_G(h \circ \theta)_\ast \circ \theta_{p,G}^\ast
\to R^{p + q}_G h_{\ast}
$$
其 $E^{pq}_2 = \check{H}^p \circ R^q_G \Gamma((h \circ \theta)_\ast)\circ \bar\theta_G^{\star}$.
\end{enonce*}

% END INCLUDED SGA4_English_current_buildable_source_20260906/source/05bis/drafts/SGA4_VBIS_LINES_1691_1929_English.texfrag


% BEGIN INCLUDED SGA4_English_current_buildable_source_20260906/source/05bis/drafts/SGA4_VBIS_LINES_1930_2298_English.texfrag
\section{下降判别准则}\etiquette[3]{Vbis.sec3}

\setcounter{subsection}{-1}
\subsection{记号}\etiquette[3.0]{Vbis.subsec3.0}

下文沿用 (\hyperref[Vbis.subsec2.4]{2.4}) 的记号.
始终假定\emph{有限纤维积在} $B$ \emph{中可表}.

还假定 $B$ 满足下述条件:

\setcounter{subsubsection}{-1}
\subsubsection{}\etiquette[3.0.0]{Vbis.subsubsec3.0.0}
\emph{$B$ 中的余积存在, 不交且普遍, 并且是有效 $\mathcal{E}$-下降与有效
$G^{\circ}$-下降的族} [参见 \cite{sga4.exp.Vbis.bib.4} (9.23)、(9.25) 及 (9.27)].

\subsubsection{}\etiquette[3.0.1]{Vbis.subsubsec3.0.1}
现在回顾范畴的单纯对象的一些经典记号
[参见 \cite{sga4.exp.Vbis.bib.3}, 第 II 章, 以及 (V, 附录)].

设 $E$ 为\emph{具有有限投射极限}的范畴.
$\Delta^{\circ}E$ 表示 $E$ 的单纯对象范畴
(参见 (\hyperref[Vbis.subsubsec1.2.6]{1.2.6})).

对整数 $n$, 用 $\mathop{\Delta}_n$
(分别用 ${\mathop{\Delta}_n}^{+}$、${\mathop{\Delta}_n}^{-}$)
表示 $\Delta$ (分别为 $\Delta^+$、$\Delta^-$) 中由满足 $p\leq n$ 的对象 $[p]$ 组成的满子范畴.

限制函子
\begin{equation*}
i_n^{\ast}:\Delta^{\circ}E\longrightarrow
{\mathop{\Delta}_n}^{\circ}E\tag{3.0.1.1}
\end{equation*}\etiquette[3.0.1.1]{Vbis.eq3.0.1.1}
有右伴随 $i_{n\ast}$, 因为 $E$ 中\pageoriginale{}存在有限投射极限
(参见 (\hyperref[Vbis.lem1.2.10]{1.2.10})).
函子 $i_{n\ast}\circ i_n^{\ast}$ 记为 $\cosk_n$, 称为 \emph{$n$ 阶余骨架函子};
滥用记号, 我们也用 $\cosk_n$ 表示函子 $i_{n\ast}$.

对每个整数 $p$, 用 $\mathop{\Delta}_{n[p]}$
(分别用 ${\mathop{\Delta}_{n[p]}}^+$) 表示 $\Delta_{[p]}$
(分别为 $\Delta_{[p]}^+$) 的满子范畴,
其中后者为 $\Delta$ (分别为 $\Delta^+$) 中 $[p]$ 上的对象范畴,
而该满子范畴由 $[p]$ 上满足 $q\leq n$ 的对象 $[q]$ 定义.
请读者验证
\begin{equation*}
(\cosk_n(X))_p=\varprojlim_{\mathop{\Delta}_{n[p]}}X_q
\xrightarrow{\sim}
\varprojlim_{\mathop{\Delta}_{n[p]}^{+}}X_q{\quoi}.
\tag{3.0.1.2}
\end{equation*}\etiquette[3.0.1.2]{Vbis.eq3.0.1.2}

\subsubsection{}\etiquette[3.0.2]{Vbis.subsubsec3.0.2}
设 $X$ 和 $X'$ 为 $\Delta^{\circ}E$ 的两个对象,
$f,g:X\rightrightarrows X'$ 为两个态射.
\emph{从 $f$ 到 $g$ 的同伦}由对每个 $n$ 给出的映射
$h_n:\Hom_{\Delta}([n],[1])\to\Hom_E(X_n,X'_n)$ 组成, 满足以下两个条件:
\begin{equation*}
h_n(\partial^1)=f_n\quad\mbox{和}\quad
h_n(\partial^\circ)=g_n\quad\mbox{对每个 }n{\quoi}.
\tag{3.0.2.1}
\end{equation*}\etiquette[3.0.2.1]{Vbis.eq3.0.2.1}

\setcounter{paragraph}{1}
\paragraph{}\etiquette[3.0.2.2]{Vbis.par3.0.2.2}
对 $\Delta$ 中的每条箭头 $[n]\to[p]$, 下图交换:
\[
\xymatrix@C=2.25cm{
\Hom_{\Delta}\big([p],[1]\big) \ar[r] \ar[dd]&
\Hom_E(X_p,X'_p) \ar[rd]&\\
& & \Hom_E(X_p,X'_p)\\
\Hom_{\Delta}\big([n],[1]\big) \ar[r]&
\Hom_E(X_n,X'_n) \ar[ru]&.
}
\]
附记: 在从 $X$ 到 $X'$ 的态射集上这样引入的关系并非等价关系.
引入\emph{复合同伦}概念可弥补这一不便
[参见 \cite{sga4.exp.Vbis.bib.3}, 第 III 章], 但我们不会使用此概念.

\begin{enonce*}{引理 3.0.2.3}\etiquette[3.0.2.3]{Vbis.lem3.0.2.3}
设\pageoriginale{} $X$ 和 $X'$ 为 $\Delta^{\circ}E$ 的两个对象,
$f,g:X\rightrightarrows X'$ 为两个态射,
$F:E\to C$ 为函子, 其中 $C$ 为阿贝尔范畴.
从 $f$ 到 $g$ 的单纯同伦诱导分别典范对应于 $F(f)$ 和 $F(g)$ 的上链复形态射之间的同伦.
\end{enonce*}

这归结为 $C$ 是某个环上的模范畴的情形, 再应用 [\cite{sga4.exp.Vbis.bib.5}, I, (3.7.1)].

\begin{enonce*}{引理 3.0.2.4}\etiquette[3.0.2.4]{Vbis.lem3.0.2.4}
设 $n$ 为整数, $X$ 和 $X'$ 为 $\Delta^{\circ}E$ 的两个对象.
设 $f$ 和 $g$ 为从 $X$ 到 $X'$ 的两个态射, 使当 $p<n$ 时 $f_p=g_p$.
则从 $\cosk_n(f)$ 到 $\cosk_n(g)$ 存在单纯同伦.
\end{enonce*}

对 $p<n$, 定义
$$h_p:\Hom_{\Delta}([p],[1])\to\Hom(X_p,X'_p)$$
为取值为 $f_p=g_p$ 的常值映射.
继而定义 $h_n:\Hom_{\Delta}([n],[1])\to\Hom(X_n,X'_n)$,
使 $\Hom_{\Delta}([n],[1])$ 中除 $\partial^{\circ}$ 外的每个元素映到 $f_n$,
而 $\partial^{\circ}$ 映到 $g_n$.
最后, 注意对 $k>n$, 有
$$
\Hom_{\Delta}([k],[1])=
\varprojlim_{\substack{[q]\to[k]\\q\leq n}}\Hom([q],[1])
$$
由此可典范定义 $h_k$.

\begingroup\tolerance=9999\emergencystretch=4em
\subsection{从 1-上同调下降的观点比较两个增广}\etiquette[3.1]{Vbis.subsec3.1}

设 $S$ 为 $B$ 的对象.
用 $\Hom_{\plat}(D^{\circ},B/S)$ 表示 $\Hom(D^{\circ},B/S)$ 中由平坦增广组成的满子范畴
[参见 (\hyperref[Vbis.subsubsec1.2.6]{1.2.6}) 及
(\hyperref[Vbis.def2.4.3.1]{2.4.3.1})].
\par\endgroup

\begin{enonce*}{命题 3.1.1}\etiquette[3.1.1]{Vbis.prop3.1.1}
设 $X\xrightarrow{u}C^D_S$ 与 $X'\xrightarrow{u'}C^D_S$
为 $\Hom_{\plat}(D^{\circ},B/S)$ 的两个对象.
设 $f:X\to X'$ 为 $S$ 上的态射 (即 $\Hom(D^{\circ},B/S)$ 中的态射).
可自然地与之联系一个函子性态射
$$
\eta^f_G:\sisi{\RR^+}{R}_G\bar u'_\ast\circ{\bar u'}_G^\ast
\longrightarrow
\sisi{\RR^+}{R}_G\bar u_\ast\circ\bar u_G^\ast
$$
使下图交换:
\[
\xymatrix@=1.5cm{
&R_G\bar u'_\ast\circ{\bar u'}_G^{\ast}
\ar[dd]^{\eta^f_G}\\
\id_{D_{G_S}^{+}(S)}\ar[ru]\ar[rd]&\\
&R_G\bar u_\ast\circ\bar u_G^{\ast}.
}
\]

此外,\pageoriginale{}若 $D=\Delta$, 则 $\eta^f_G$ 仅依赖于
$f$ 在 $\Hom(\Delta^{\circ},B/S)$ 中的同伦类
\emph{(参见 (\hyperref[Vbis.subsubsec3.0.2]{3.0.2}))}.
\end{enonce*}

显然, 为构造 $\eta_G^f$, 可设 $G=\Ysheaf{\Mod}(\mathcal{E}^{\circ},\mathcal{O})$.

令 $I$ 为由图
\begin{equation*}
{0}\longrightarrow{1}{\quoi}.\tag{3.1.1.1}
\end{equation*}\etiquette[3.1.1.1]{Vbis.eq3.1.1.1}
定义的范畴.
用 $r_0:D\to I\times D$ (分别用 $r_1$) 表示由 $r_0(i)=(0,i)$
(分别为 $r_1(i)=(1,i)$) 定义的全忠实函子.

由典范同构
\begin{equation*}
\Hom((I\times D)^{\circ},B/S)\xrightarrow{\sim}
\Hom(I^{\circ},\Hom(D^{\circ},B/S))\tag{3.1.1.2}
\end{equation*}\etiquette[3.1.1.2]{Vbis.eq3.1.1.2}
(\hyperref[Vbis.prop3.1.1]{3.1.1}) 的资料定义平坦增广
\nde{明确地说, 有 $(XfX')_{(0,i)}=X_i$ 与 $(XfX')_{(1,j)}=X'_j$;
与 $\alpha\in\Hom_D(i,j)=\Hom((0,i),(0,j))$ 对应的是
$\alpha\circ f_i:(XfX')_{(0,i)}\to(XfX')_{(1,j)}$,
且 $(\varepsilon,i)$, $\varepsilon=0,1$ 的自态射按显然方式作用.}
\begin{equation*}
XfX'\xrightarrow{ufu'}C^{I\times D}_S{\quoi}.\tag{3.1.1.3}
\end{equation*}\etiquette[3.1.1.3]{Vbis.eq3.1.1.3}

设 $F$ 为 $\Mod(\Ee^\circ_S,\Oo_S)$ 的对象. 典范态射
\begin{equation*}
\Ee_{I\times D}^{\ast}(F)\longrightarrow
\Gamma((ufu')_\ast)\circ\Gamma((ufu')^{\ast})
(\Ee_{I\times D}^{\ast}(F)){\quoi}\tag{3.1.1.4}
\end{equation*}\etiquette[3.1.1.4]{Vbis.eq3.1.1.4}
可解释为 $\Mod(\Ysheaf{\Gamma}(\bar{C^D_S}),\Oo)$ 中的交换三角:
\[\vcenter{
\xymatrix@=1.5cm{
&\Gamma(u'_\ast)\circ\Gamma({u'}^{\ast})(\Ee_D^{\ast}(F))\ar[dd]\\
\Ee_D^{\ast}(F)\ar[ru]\ar[rd]&\\
&\Gamma(u_\ast)\circ\Gamma(u^{\ast})(\Ee_D^{\ast}(F))
}}\leqno{(3.1.1.5)}
\]\etiquette[3.1.1.5]{Vbis.eq3.1.1.5}
从而\pageoriginale{}给出函子性态射
$\bar u'_\ast\circ{\bar u'}^{\ast}\xrightarrow{\alpha_f}\bar u_\ast\circ\bar u^{\ast}$,
使图
\[\vcenter{
\xymatrix@=1.5cm{
&\bar u'_\ast\circ{\bar u'}^{\ast}\ar[dd]^{\alpha_f}\\
\id_{\Mod(\Ee^{\circ}_S,\Oo_S)}\ar[ru]\ar[rd]&\\
&\bar u_\ast\circ\bar u^{\ast}
}}\leqno{(3.1.1.6)}
\]\etiquette[3.1.1.6]{Vbis.eq3.1.1.6}
交换.

因而得到一个函子性态射, 使图
\[\vcenter{
\xymatrix@=1.5cm{
&K^{+}(\bar u'_\ast)\circ K^{+}({\bar u'}^{\ast})
\ar[dd]^{K^{+}(\alpha_f)}\\
\id_{K^{+}(\Ee^{\circ}_S,\Oo_S)}\ar[rd]\ar[ru]&\\
&K^{+}(\bar u_\ast)\circ K^{+}(\bar u^{\ast})
}}\leqno{(3.1.1.7)}
\]\etiquette[3.1.1.7]{Vbis.eq3.1.1.7}
交换.

\noindent
(注意, $K^+(\alpha_f)$ 也可用同构
$K^+(\Gamma(f^{\ast}))\circ K^+({\bar u'}^{\ast})\xrightarrow{\sim}K^+(\bar u^{\ast})$
通过伴随关系形式地得到.)

设 $F^{\bigdot}$ 为 $K^+(S)$ 中的复形,
$\xi:K^+(\overline{ufu'}^{\ast})(F^{\bigdot})\to J^{\bigdot}$
为拟同构, 使对每个整数 $n$, $J^n$ 逐对象完全无上同调
(参见 (\hyperref[Vbis.rem1.3.13]{1.3.13})).
由 (\loccit), 典范箭头
$$
K^+({\bar u'}^\ast)(F^{\bigdot})\to r_1^{\ast}(J^{\bigdot})
\quad\mbox{和}\quad
K^+(\bar u^\ast)(F^{\bigdot})\to r_0^{\ast}(J^{\bigdot})
$$
为\pageoriginale{}拟同构 (记号参见 (\hyperref[Vbis.subsubsec1.2.8]{1.2.8})).
因而得到态射
$$
\sisi{\RR^+}{R}\bar u'_\ast\circ{\bar u'}^\ast(F^{\bigdot})
\xrightarrow{\eta^f(F^{\bigdot})}
\sisi{\RR^+}{R}\bar u_\ast\circ\bar u^\ast(F^{\bigdot})
$$
使下图在 $D^+(S)$ 中交换:
{\def\objectstyle{\scriptstyle}
\def\labelstyle{\scriptstyle}

\[\vcenter{
\xymatrix@C=0cm@R=1cm{
&Q\circ K^{+}(\bar u'_\ast)(r^{\ast}_{0}(J^{\bigdot}))
\ar[rr]^{\sim}\ar@{->}'[d][ddd]&&
R\bar u'_\ast\circ
Q(r^{\ast}_0(J^{\bigdot}))\ar[ddd]\\
Q\circ K^{+}(\bar u'_\ast)\circ K^{+}({\bar u'}^{\ast})(F^{\bigdot})
\ar[rr]\ar[ru]\ar[ddd]_{Q(K^{+}(\alpha_f)(F^{\bigdot}))}&&
R\bar u'_\ast\circ
{\bar u'}^{\ast}(F^{\bigdot})\ar[ru]^{\sim}
\ar@{-->}[ddd]^(.35){\eta^{f}(F^{\bigdot})}&\\
&&&\\
&Q\circ K^{+}(\bar u_\ast)(r^{\ast}_1(J^{\bigdot}))
\ar@{->}'[r][rr]^(.25){\sim}&&
R\bar u_\ast\circ Q(r_1^{\ast}(J^{\bigdot}))\\
Q\circ K^{+}(\bar u_\ast)\circ K^{+}(\bar u^{\ast})(F^{\bigdot})
\ar[ru]\ar[rr]&&
R\bar u_\ast\circ\bar u^{\ast}(F^{\bigdot})
\ar[ru]&
}}\leqno{(3.1.1.8)}
\]\etiquette[3.1.1.8]{Vbis.eq3.1.1.8}
}
\noindent
显然, $\eta^f(F^{\bigdot})$ 不依赖所选分解;
且可验证, 当 $F^{\bigdot}$ 在 $D^{+}(S)$ 中变动时, 它具有函子性.

为完成 \hyperref[Vbis.prop3.1.1]{(3.1.1)} 的证明, 只须证明下述引理:

\begin{enonce*}{引理 3.1.1.9}\etiquette[3.1.1.9]{Vbis.lem3.1.1.9}
设\pageoriginale{} $X\xrightarrow{u}C^{\Delta}_S$ 与 $X'\xrightarrow{u'}C^{\Delta}_S$
为 $\Hom_{\plat}(\Delta^{\circ},B/S)$ 的两个对象.
设 $f,g:X\to X'$ 为 $\Hom(\Delta^{\circ},B/S)$ 中的两个态射.
若从 $f$ 到 $g$ 存在单纯同伦, 则 $\eta^f=\eta^g$.
\end{enonce*}

设 $h_n:\Hom_{\Delta}([n],[1])\to\Hom_{B/S}(X_n,{X'}_n)$
为从 $f$ 到 $g$ 的单纯同伦.
令 $\overline{I\times D}$ 为从 $I\times D$ 添加与
$h_n(\Hom([n],[1]))$ 一一对应的箭头 $(1,n)\to(0,n)$,
以及由 $h_n$ 随 $n$ 变动的相容性所规定的关系后得到的范畴.
引理的资料定义平坦增广
$\overline{XfX'}\xrightarrow{t}C^{\overline{I\times D}}_S$.
设 $F^{\bigdot}$ 为 $C^+(S)$ 中的复形,
$J^{\bigdot}$ 为 $C^+(t^{\ast})(F^{\bigdot})$ 的分解,
使对每个整数 $n$, $J^n$ 逐对象完全无上同调.
在
$$
C^{+}(\Gamma(u'_\ast))(r^{\ast}_0(J^{\bigdot}))
\rightrightarrows
C^{+}(\Gamma(u_\ast))(r^\ast_1(J^{\bigdot})),
$$
之间存在单纯同伦, 因而在过渡到总复形后存在同伦
(参见 (\hyperref[Vbis.lem3.0.2.3]{3.0.2.3})).

\begin{enonce*}{定义 3.1.2}\etiquette[3.1.2]{Vbis.def3.1.2}
采用 \emph{\hyperref[Vbis.prop3.1.1]{(3.1.1)}} 的记号,
若 $\eta^f_G$ 为同构, 就称 $f$ 为 \emph{$G$-1-上同调下降的等价}.
\end{enonce*}

\setcounter{subsubsection}{2}
\subsubsection{}\etiquette[3.1.3]{Vbis.subsubsec3.1.3}
下文用 $L_i:\Delta\to\Delta\times\Delta$
(分别用 $C_i:\Delta\to\Delta\times\Delta$) 表示由
$[n]\longmapsto[n]\times[i]$ (分别为 $[n]\longmapsto[i]\times[n]$) 定义的典范函子.

\begin{enonce*}{命题 3.1.4}\etiquette[3.1.4]{Vbis.prop3.1.4}
设 $X\xrightarrow{u}C^{\Delta\times\Delta}_S$ 与
$X'\xrightarrow{u'}C^{\Delta\times\Delta}_S$
为 $\Hom_{\plat}((\Delta\times\Delta)^{\circ},B/S)$ 的两个对象.
设 $f:X\to X'$ 为 $S$ 上的态射.
假定对每个整数 $i$, $f\ast\id_{L_i}:X\circ L_i\to X'\circ L_i$
(分别为 $f\ast\id_{C_i}:X\circ C_i\to X'\circ C_i$)
为 $G$-1-上同调下降的等价. 则 $f$ 为 $G$-1-上同调下降的等价.
\end{enonce*}

由 $\eta^f$ 的描述 (参见 (\hyperref[Vbis.prop3.1.1]{3.1.1}) 的证明),
只须证明某个三重复形态射在总复形的上同调上诱导同构.
标准的谱序列论证遂给出结论.

\begin{enonce*}{命题 3.1.5}\etiquette[3.1.5]{Vbis.prop3.1.5}
设\pageoriginale{} $X$ 与 $X'$ 为 $\Hom_{\plat}(\Delta^{\circ},B/S)$ 的两个对象.
设 $f:X\to X'$ 为函子性态射, 使 $f_n:X_n\to X'_n$ 为 $G$-1-上同调下降态射.
则典范态射
$$
[[X|_fX']]\to[[X'|_{\id}X']]
$$
\emph{(参见 (\hyperref[Vbis.subsubsec1.2.7]{1.2.7}))} 为 $G$-1-上同调下降的等价.
\end{enonce*}

这来自 (\hyperref[Vbis.prop3.1.4]{3.1.4}) 及下述引理, 其证明留待读者完成:

\begin{enonce*}{引理 3.1.5.1}\etiquette[3.1.5.1]{Vbis.lem3.1.5.1}
设 $f:R\to S$ 为平坦态射, $Y\xrightarrow{u}C^{\Delta}_S$ 与
$Y\xrightarrow{v}C^{\Delta}_R$ 为两个\emph{平坦}增广,
使下图在 $\Hom(\Delta^{\circ},B)$ 中交换:
\[
\xymatrix@=1.5cm{
Y\ar[rr]^{v}\ar[rd]_{u}&&C^{\Delta}_R\ar[ld]^f\\
&C^{\Delta}_S&
}
\]
若 $v$ 为 $G$-1-上同调下降增广, 则它为 $G$-1-上同调下降的等价.
\end{enonce*}

在应用中, 将把 (\hyperref[Vbis.prop3.1.4]{3.1.4})
和 (\hyperref[Vbis.prop3.1.5]{3.1.5}) 与下述结果结合:

\begin{enonce*}{命题 3.1.6}\etiquette[3.1.6]{Vbis.prop3.1.6}
设 $X\xrightarrow{u}C^{\Delta\times\Delta}_S$ 为平坦增广.
要使 $u$ 为 $G$-1-上同调下降增广, 只须对每个整数 $i$,
$$
X\circ L_i\xrightarrow{u\ast\id_{L_i}}C^{\Delta}_S
\quad(\resp\ X\circ C_i\xrightarrow{u\ast\id_{C_i}}C^{\Delta}_S)
$$
均如此.
\end{enonce*}

这来自 $\sisi{\RR^+}{R}_G\bar u_\ast$ 的显式构造
(参见 (\hyperref[Vbis.subsubsec2.3.10]{2.3.10})).

\begin{enonce*}{推论 3.1.7}\etiquette[3.1.7]{Vbis.cor3.1.7}
设 $X\xrightarrow{u}C^{\Delta}_S$ 为增广.
则 $u$ 为 $G$-1-上同调下降增广当且仅当
$[[X|_{\id}X]]\to C^{\Delta\times\Delta}_S$ 如此.
\end{enonce*}

% END INCLUDED SGA4_English_current_buildable_source_20260906/source/05bis/drafts/SGA4_VBIS_LINES_1930_2298_English.texfrag


% BEGIN INCLUDED SGA4_English_current_buildable_source_20260906/source/05bis/drafts/SGA4_VBIS_LINES_2299_2585_English.texfrag
\subsection{局部化判别准则}\etiquette[3.2]{Vbis.subsec3.2}\pageoriginale

\begin{enonce*}{命题 3.2.1}\etiquette[3.2.1]{Vbis.prop3.2.1}
设 $Y\xrightarrow{v}C^{\Delta}_S$ 为 $G$-上同调 $1$-下降增广.
要使平坦增广 $X\xrightarrow{u}C^{\Delta}_S$ 为 $G$-上同调 $1$-下降增广,
只须它在所有基变换 $Y_n\to S$ 后变成这样的增广.
\end{enonce*}

通过基变换 $Y_n\to S$, 构造向 $S$ 增广的双重半单纯对象 $\fprod{X}{Y}{S}$.
由 (\hyperref[Vbis.prop3.1.4]{3.1.4}) 及
(\hyperref[Vbis.lem3.1.5.1]{3.1.5.1}), 典范态射
$$
\fprod{X}{Y}{S}\to[[X|_{\id}X]]
$$
为 $G$-上同调 $1$-下降的等价.
由 (\hyperref[Vbis.cor3.1.7]{3.1.7}), 因而只须证明
$\fprod{X}{Y}{S}$ 向 $S$ 的增广为 $G$-上同调 $1$-下降增广.
这来自典范态射
$$
\fprod{X}{Y}{S}\to[[Y|_{\id}Y]]
$$
为 $G$-上同调 $1$-下降的等价这一事实;
再次使用 (\hyperref[Vbis.prop3.1.4]{3.1.4})、
(\hyperref[Vbis.lem3.1.5.1]{3.1.5.1}) 及 (\hyperref[Vbis.cor3.1.7]{3.1.7}).

\setcounter{subsubsection}{1}
\subsubsection{}\etiquette[3.2.2]{Vbis.subsubsec3.2.2}
设
\[
\vcenter{
\xymatrix@C=3cm@R=1.5cm{
Y'\ar[r]^{\mathsf g'}\ar[d]_{f'}&Y\ar[d]^{f}\\
X'\ar[r]_{\mathsf g}&X
}}\leqno{(3.2.2.1)}
\]\etiquette[3.2.2.1]{Vbis.eq3.2.2.1}
为 $B$ 中的交换图, $F$ 为 $\Mod(\mathcal E^{\circ}_Y,\mathcal O_Y)$ 的对象.
存在唯一态射
$$
\varphi(F):g_{\ast}(f^{\ast}(F))\to
f'^{\ast}(g'_{\ast}(F))
$$
使下图 (在 $\Ysheaf{\Mod}(\mathcal E^{\circ},\mathcal O)$ 中) 交换
\nde{如前所述, 须将交换图视为 $B^\circ$ 中的图,
并使用 $\Mod(\mathcal E^\circ,\mathcal O)\rightarrow B^\circ$ 的纤维与余纤维结构,
这解释了基变换态射的不寻常形式.
下图的各 (伴随) 箭头下写出了它们在 $B^\circ$ 中的像}:
\[
\vcenter{
\xymatrix@R=1.85cm@C=1.75cm{
&f^{\ast}(F)\ar[r]^{f^\circ}\ar[ld]_{g^\circ}&F\ar[d]^{g'^\circ}\\
g_{\ast}(f^{\ast}(F))\ar[rd]_{\Id}&&g'_{\ast}(F)\\
&f'^{\ast}(g'_{\ast}(F))\ar[ru]_{f'^\circ}&
}}\leqno{(3.2.2.2)}
\]\etiquette[3.2.2.2]{Vbis.eq3.2.2.2}

现若\pageoriginale{}假定 \eqhyperref[Vbis.eq3.2.2.1]{3.2.2.1} 的所有箭头均平坦,
则以相同方式定义一个态射, 称为\emph{基变换}态射
\nde{这里同样, 以 $f^\circ$、$g^\circ$、$f'^\circ$、$g'^\circ$
代替 $f$、$g$、$f'$、$g'$ 更为正确.
若考虑非反范畴的纤维与余纤维结构, 则 $\xi_G$ 变为
$$g_G^\ast\circ R_Gf_{\ast}\longrightarrow R_Gf'_{\ast}\circ g_G'^\ast.$$}
\begin{equation*}
\xi_G:g_{G\ast}\circ\sisi{\RR^+}{R}_Gf^{\ast}
\longrightarrow\sisi{\RR^+}{R}_Gf'^{\ast}\circ g'_{G\ast}{\quoi}.
\tag{3.2.2.3}
\end{equation*}\etiquette[3.2.2.3]{Vbis.eq3.2.2.3}

\begin{enonce*}{定义 3.2.2.4}\etiquette[3.2.2.4]{Vbis.def3.2.2.4}
若 $f$、$g$、$f'$、$g'$ 为平坦态射且 $\xi_G$ 为同构,
就称图 \emph{\hyperref[Vbis.eq3.2.2.1]{(3.2.2.1)} 满足相对于 $G$ 的基变换定理}.
\end{enonce*}

\setcounter{subsubsection}{2}
\subsubsection{}\etiquette[3.2.3]{Vbis.subsubsec3.2.3}
设 $h:S\to S'$ 为 $B$ 中的\emph{平坦}态射,
$X\xrightarrow{u}C^{\Delta}_S$ 与 $X'\xrightarrow{u'}C^{\Delta}_{S'}$
为两个\emph{平坦}增广, $f:X\to X'$ 为函子性态射, 使图
(在 $\Hom(\Delta^0,B)$ 中)
\[
\vcenter{
\xymatrix@C=3cm@R=1.5cm{
X\ar[r]^{f}\ar[d]_{u}&X'\ar[d]^{u'}\\
C^{\Delta}_S\ar[r]^{h_c}&C^{\Delta}_{S'}
}}\leqno{(3.2.3.1)}
\]\etiquette[3.2.3.1]{Vbis.eq3.2.3.1}
交换. (这里 $h_c$ 表示由 $h$ 定义的常值函子性态射.)

设 $K^{\bigdot}$ 为 $D_G^+(S')$ 中的复形. 态射
$$
\eta_G^f(K^{\bigdot}):
\sisi{\RR^+}{R}\bar u'_{\ast}\circ\bar u'^{\ast}(K^{\bigdot})
\longrightarrow
\sisi{\RR^+}{R}(\overline{h_c\circ u})_{\ast}\circ
(\overline{h_c\circ u})^{\ast}(K^{\bigdot})
$$
考虑到恒等式
$$
\sisi{\RR^+}{R}(\overline{h_c\circ u})_{\ast}\simeq
\sisi{\RR^+}{R}h_c^{\ast}\circ
\sisi{\RR^+}{R}\bar u_{\ast}
\quad\mbox{和}\quad
(\overline{h_c\circ u})^{\ast}\simeq
\bar u^{\ast}\circ h_{c\ast},
$$
诱导一个对 $K^{\bigdot}$ 函子性的态射 $\ch_G^{f,h}(K^{\bigdot})$,
使下图交换\nde{为便于阅读, 在导出函子中略去了厚范畴 $G$ 的标记.
由于选取余纤维结构, 同样不寻常的导出式 $h_{c\ast}$ 再次出现.}:
\[
\vcenter{
\xymatrix@C=1cm@R=2.25cm{
{}_{}h_{c\ast}\circ
R_{}\bar u'_{\ast}{}_{}\circ
\bar u'^{\ast}(K^{\bigdot})
\ar[rr]^{\ch_G^{f,h}{(K^{\bigdot})}}&&
R_{}\bar u_{\ast}\circ
{}_{}\bar u^{\ast}\circ{}_{}h_{c\ast}(K^{\bigdot})\\
&h_{c\ast}(K^{\bigdot})
\ar[lu]^{{}_{}h_{c\ast}(\alpha(K^{\bigdot}))}\ar[ru]&.
}}\leqno{(3.2.3.2)}
\]\etiquette[3.2.3.2]{Vbis.eq3.2.3.2}

\begin{enonce*}{引理 3.2.3.3}\etiquette[3.2.3.3]{Vbis.lem3.2.3.3}
假定\pageoriginale{}对每个整数 $i$, 图
\[
\xymatrix@C=3cm@R=1.5cm{
X_i\ar[d]_{u_i}\ar[r]^{f_i}&X'_i\ar[d]^{u'_i}\\
S\ar[r]^{h}&S'
}
\]
满足相对于 $G$ 的基变换定理. 则 $\ch_G^{f,h}$ 为同构.
\end{enonce*}

态射 $\ch_G^{f,h}(K^{\bigdot})$ 可如下计算. 考虑增广
$$
X f X'\xrightarrow{(h_c\circ u)f u'}C^{I\times\Delta}_S
$$
(参见 \eqhyperref[Vbis.eq3.1.1.2]{3.1.1.2}),
并选取 $K^+(\overline{(h_c\circ u)f u'}^{\ast})(K^{\bigdot})$ 的分解
$J^{\bigdot}$, 使对每个整数 $n$, $J^n$ 逐对象完全无上同调.

得到拟同构
$$
\theta:K^+(\Gamma(f^{\ast}))(r_0^{\ast}J^{\bigdot})
\longrightarrow r_1^{\ast}(J^{\bigdot})
$$
以及态射
$$
t:K^+(\Gamma(h_c^{\ast}))\circ K^+(\Gamma(u'_{\ast}))
(r_0^{\ast}(J^{\bigdot}))\longrightarrow
K^+(\Gamma(u_{\ast}))(r_1^{\ast}(J^{\bigdot}))
$$
它有分解
\begin{align*}
&K^+(\Gamma(h_c^{\ast}))\circ K^+(\Gamma(u'_{\ast}))
(r_0^{\ast}(J^{\bigdot}))\xrightarrow{\varphi}
K^+(\Gamma(u_{\ast}))\circ K^+(\Gamma(f^{\ast}))
(r_0^{\ast}(J^{\bigdot}))
\xrightarrow{K^+(\Gamma(u_{\ast}))(\theta)}\\
&\longrightarrow K^+(\Gamma(u_{\ast}))(r_1^{\ast}(J^{\bigdot})),
\end{align*}
且\pageoriginale{} $\ch^{f,h}(K^{\bigdot})$ 是取 $t$ 在函子
$\sisi{\RR^+}{R}\epsilon_{Z^{\bigdot\ast}}$ 下的像得到的.

上述分解表明, $t$ “逐对象”等同于与下列各图相联系的基变换态射
\[
\xymatrix@C=3cm@R=1.5cm{
X_i\ar[d]_{u_i}\ar[r]^{f_i}&X'_i\ar[d]^{u'_i}\\
S\ar[r]^{h}&S'
}
\]
由谱序列得出 $\ch_G^{f,h}(K^{\bigdot})$ 为同构, 证明完毕.

由此得到\emph{第二个局部化判别准则}:

\begin{enonce*}{命题 3.2.4}\etiquette[3.2.4]{Vbis.prop3.2.4}
设 $X\xrightarrow{u}C^{\Delta}_S$ 为 $\Hom_{\plat}(\Delta^{\circ},B/S)$ 的对象.
要使 $u$ 为 $G$-上同调 $1$-下降增广,
只须它在“对 $G$ 足够多”的\footnote{“对 $G$ 足够多”意指存在一族满足前述条件的平坦态射
$(S_{\alpha}\xrightarrow{h_{\alpha}}S)_{\alpha\in A}$,
使函子族 $(h_{\alpha\ast}:G_S\to G_{S_\alpha})_{\alpha\in A}$ \emph{保守}.}
平坦基变换 $h:S'\to S$ 后变成这样的增广, 而且笛卡尔图
\[
\xymatrix@C=3cm@R=1.5cm{
X'_i\ar[d]\ar[r]&X_i\ar[d]\\
S'\ar[r]^{h}&S
}
\]
满足相对于 $G$ 的基变换定理.
\end{enonce*}

命题\pageoriginale{} (\hyperref[Vbis.prop3.2.4]{3.2.4})
考虑到图 \eqhyperref[Vbis.eq3.2.3.2]{3.2.3.2}, 立即来自
(\hyperref[Vbis.lem3.2.3.3]{3.2.3.3}).

\subsection{上同调下降态射的性质}
\etiquette[3.3]{Vbis.subsec3.3}

本节假定 \emph{$B$ 中的平坦态射类对基变换稳定}.

\begin{enonce*}{命题 3.3.1}\etiquette[3.3.1]{Vbis.prop3.3.1}
设 $i=1,2$.
\begin{enumerate}
\renewcommand{\theenumi}{\alph{enumi}}
\renewcommand{\labelenumi}{\theenumi)}
\item 每个有截面的平坦态射都是普遍 $G$-上同调 $2$-下降态射.

\item 设 $f:X\to S$ 为平坦态射, $g:S'\to S$ 为 $G$-上同调 $i$-下降态射,
令 $X'=\fprod{X}{S'}{S}$, 并设 $f'=f_{(S')}:X'\to S'$.
则 $f$ 为普遍 $G$-上同调 $i$-下降态射当且仅当 $f'$ 如此.

\item 若两个态射 $f:X\to Y$ 和 $g:Y\to Z$ 的复合为普遍 $G$-上同调 $i$-下降态射,
则 $g$ 为普遍 $G$-上同调 $i$-下降态射.

\item 两个普遍 $G$-上同调 $i$-下降态射的复合仍为普遍 $G$-上同调 $i$-下降态射.

\item 若 $f:X\to X'$ 和 $g:Y\to Y'$ 为两个普遍 $G$-上同调 $i$-下降的 $S$-态射,
则 $\fprod{f}{g}{S}$ 为普遍 $G$-上同调 $i$-下降态射.

\item 设 $(u_{\alpha}:X_{\alpha}\to Y_{\alpha})_{\alpha\in A}$
为一族 $G$-上同调 $i$-下降态射.
则 $\amalg u_{\alpha}:\amalg_{\alpha\in A}X_{\alpha}\to\amalg_{\alpha\in A}Y_{\alpha}$
为 $G$-上同调 $i$-下降态射.
\end{enumerate}
\end{enonce*}

\noindent\emph{证明.}
关于\pageoriginale{}相对于 $G$ 的有效下降, 参阅 \cite{sga4.exp.Vbis.bib.4}.
\begin{enumerate}
\renewcommand{\theenumi}{\alph{enumi}}
\renewcommand{\labelenumi}{\theenumi)}
\item 利用 $f$ 的截面 $s$, 比较向 $S$ 增广的半单纯对象 $[X|_fS]$ 与 $[S|_{\id}S]$.
由 (\hyperref[Vbis.lem3.0.2.4]{3.0.2.4}), 得到 $G$-上同调 $1$-下降的等价
\nde{为此, 注意 $[X|_fS]$ 与 $[S|_{\id}S]$ 分别就是 $f$ 与 $\id$ 的 $0$ 阶余骨架,
故由 \hyperref[Vbis.prop3.1.1]{3.1.1}, $s$ 与 $f$ 在上同调中诱导互逆同构.}.

\item 来自 (\hyperref[Vbis.prop3.2.1]{3.2.1}).

\item 来自 a) 和 b).

\item 设 $f:X\to Y$ 与 $g:Y\to S$ 为两个普遍 $G$-上同调 $1$-下降态射.
令 $h=g\circ f$, 考虑纤维积
\[
\xymatrix@C=3cm@R=1.5cm{
Y\ar[d]_{g}&R\ar[d]^{g'}\ar[l]_{h'}\\
S&X\ar[l]^{h}
}
\]
态射 $g'$ 有截面 $s:X\to R$, 使 $h'\circ s=f$.
由 c), $h'$ 为普遍 $G$-上同调 $1$-下降态射, 再由 b), $h$ 也如此.

\item 由 b) 和 d) 形式地得出.

\item 设 $(Z_{\lambda})_{\lambda\in\Lambda}$ 为 $B$ 的一族对象.
由假设 (\hyperref[Vbis.subsubsec3.0.0]{3.0.0}), 存在典范等价
$$
\Mod\left(\mathcal E^{\circ}_{\coprod_{\lambda}Z_{\lambda}},
\mathcal O_{\coprod_{\lambda}Z_{\lambda}}\right)
\xrightarrow{\sim}\displaystyle\mathop{\prod}_{\lambda}
\Mod(\mathcal E^{\circ}_{Z_\lambda},\mathcal O_{Z_\lambda})
$$
\end{enumerate}
它诱导等价
$$
G_{\coprod_{\lambda}Z_{\lambda}}
\xrightarrow{\sim}\prod_{\lambda}G_{Z_{\lambda}}{\quoi}.
$$
此外, 一个族
$(Q_{\lambda})_{\lambda}\in\prod_{\lambda}\Mod(\mathcal E^{\circ}_{Z_\lambda},\mathcal O_{Z_\lambda})$
对应于完全无上同调对象, 当且仅当对每个 $\lambda$, $Q_{\lambda}$ 完全无上同调.

现设 $(u_{\lambda}:X_{\lambda}\to Y_{\lambda})_{\lambda\in\Lambda}$
为一族 $G$-上同调 $1$-下降态射.
由于 $B$ 中的直和不交且\pageoriginale{}普遍, 对每个 $n$ 有等同
$$
\left[\displaystyle\mathop{\coprod}_{\lambda}X_{\lambda}
\mathrel{|}_{\coprod u_{\lambda}}\coprod Y_{\lambda}\right]_n
\simeq\displaystyle\mathop{\coprod}_{\lambda}
[X_{\lambda}|_{u_{\lambda}}Y_{\lambda}]_n
$$
它对 $n$ 函子性.
于是由上述讨论及 $\sisi{\RR^+}{R}_G\theta_{\ast}$ 的显式计算
(参见 \hyperref[Vbis.subsec2.3]{(2.3)}),
$\coprod_{\lambda}u_{\lambda}$ 为 $G$-上同调 $1$-下降态射; 细节留待读者完成.
读者还可验证, \emph{若每个 $u_{\lambda}$ 在任意基变换后仍为 $G$-上同调
$1$-下降态射, 则 $\coprod_{\lambda}u_{\lambda}$ 也如此}.

\setcounter{subsubsection}{1}

% END INCLUDED SGA4_English_current_buildable_source_20260906/source/05bis/drafts/SGA4_VBIS_LINES_2299_2585_English.texfrag


% BEGIN INCLUDED SGA4_English_current_buildable_source_20260906/source/05bis/drafts/SGA4_VBIS_LINES_2586_2816_English.texfrag
\subsubsection{}\etiquette[3.3.2]{Vbis.subsubsec3.3.2}
对 $B$ 的每个对象 $X$, 对应以满足
$\coprod_\alpha X_\alpha\to X$ 为普遍 $G$-$i$-上同调下降态射的
态射族 $(X_\alpha\to X)_{\alpha\in A}$ 所成的集合.
由 \hyperref[Vbis.prop3.3.1]{3.3.1}, 这定义 $B$ 上的一个预拓扑.

此预拓扑生成的拓扑称为 \emph{$G$-$i$-上同调下降拓扑}.
由 \hyperref[Vbis.prop3.3.1]{3.3.1} c), 此拓扑的覆盖态射恰为普遍
$G$-$i$-上同调下降态射.
最后, 注意 $G$-$2$-上同调下降拓扑比 $G^{\circ}$-下降拓扑粗
[参见 \cite{sga4.exp.Vbis.bib.4} (6.23)].

\begin{enonce*}{定理 3.3.3}\etiquette[3.3.3]{Vbis.thm3.3.3}
假定 \emph{$B$ 的所有箭头均平坦}. 设 $S$ 为 $B$ 的对象.
$S$ 关于 $G$-$1$-上同调下降拓扑的每个对象均可表的超覆盖
\emph{(参见 V, 附录)} 都定义一个普遍 $G$-$1$-上同调下降增广.
\end{enonce*}

证明分两步; 明确规定所有余骨架均在范畴 $B/S$ 中计算.

\begin{enonce*}{引理 3.3.3.1}\etiquette[3.3.3.1]{Vbis.lem3.3.3.1}
设\pageoriginale{} $X$ 为 $\Hom(\Delta^{\circ},B/S)$ 的对象.
要使 $X$ 为 $1$-上同调下降对象 (分别为普遍 $1$-上同调下降对象),
只须对所有充分大的 $n$, $\cosk_n(X)$ 具有此性质.
\end{enonce*}

事实上, 设 $\alpha_n:X\to\cosk_n(X)$ 为典范态射.
可验若 $K^{\bigdot}$ 为 $D^+_{G_S}(S)$ 的复形,
使当 $j<N$ 时 $H^j(K^{\bigdot})=0$,
则当 $i<N+n$ 时 $H^i(\eta_G^{\alpha_n})$ 为同构, 断言由此得出.

\begin{enonce*}{引理 3.3.3.2}\etiquette[3.3.3.2]{Vbis.lem3.3.3.2}
设 $n\geq0$, $X$ 与 $X'$ 为 $B/S$ 的两个单纯对象, $f:X\to X'$ 为态射.
假定:
\begin{enumerate}
\renewcommand{\theenumi}{\roman{enumi}}
\renewcommand{\labelenumi}{(\theenumi)}
\item $X\to\cosk_{n+1}(X)$ 为同构;

\item $X'\to\cosk_{n+1}(X')$ 为同构;

\item 当 $i\leq n$ 时 $f_i:X_i\to X'_i$ 为同构;

\item $f_{n+1}$ 为普遍 $G$-$1$-上同调下降态射.
\end{enumerate}
则若 $X$ 为普遍 $G$-$1$-上同调下降对象, $X'$ 也如此.
\end{enonce*}

引理 \hyperref[Vbis.lem3.3.3.1]{3.3.3.1} 与
\hyperref[Vbis.lem3.3.3.2]{3.3.3.2} 显然通过归纳证明定理 \hyperref[Vbis.thm3.3.3]{3.3.3}.

\begin{enonce*}{引理 3.3.3.3}\etiquette[3.3.3.3]{Vbis.lem3.3.3.3}
在 \hyperref[Vbis.lem3.3.3.2]{3.3.3.2} 的假设下,
所有态射 $f_p$ 均为普遍 $G$-$1$-上同调下降态射.
\end{enonce*}

当 $p\leq n+1$ 时这是显然的.
对 $p>n+1$, $X_p$ (分别为 $X'_p$) 可写为
$\varprojlim_{\substack{\Delta^+\\ n+1[p]}}X_q$
(分别为 $\varprojlim_{\substack{\Delta^+\\ n+1[p]}}X'_q$).

可设 $p>n+1$. 将
$$
\cosq_{n+1}(X)_p=
\varprojlim_{\substack{[q]\rightarrow[p]\\q\leq n+1}}X_q
$$
写为箭头对
$$
\Pi X=\prod_{\substack{[q]\rightarrow[p]\\q\leq n+1}}X_q
\rightrightarrows
\prod_{\substack{[i]\stackrel{\alpha}{\hookrightarrow}[j]\\
\searrow\swarrow\\{}[p]\\j\leq n+1}}X_i.
$$
的核.
此箭头对中以 $\alpha\in\Hom_{[p]}([i],[j])$ 为指标的分量 $\alpha_X$,
一方面由以 $[i]\to[p]$ 为指标的投影 $\Pi X\to X_i$ 组成,
另一方面由以 $[j]\to[p]$ 为指标的投影与 $\alpha\in\Hom(X_j,X_i)$ 的复合
$$
\Pi X\longrightarrow X_j\longrightarrow X_i
$$
组成.

由此得图
$$
\xymatrix{\cosq_{n+1}(X)_p\ar[r]\ar[d]&\Pi X\ar[d]\\
\cosq_{n+1}(X')_p\ar[r]&\Pi X'}
$$
为笛卡尔图. 事实上, 由上述描述,
只须证明对每个\textit{递增单射} $\alpha\in\Hom_{[p]}([i],[j])$, 图
$$
\xymatrix{\Ker(\alpha_X:\Pi X\rightrightarrows X_i)\ar[r]\ar[d]&
\Pi X\ar[d]\\
\Ker(\alpha_{X'}:\Pi X'\rightrightarrows X'_i)\ar[r]&\Pi X'}
$$
为笛卡尔图. 若 $i=n+1$, 则 $\alpha:[i]\hookrightarrow[j]$
为 $[n+1]$ 的恒等态射, 因而该箭头对不施加条件.
相应的核为 $\Pi X$ 与 $\Pi X'$, 此情形得证.
否则 $i<n+1$, 且 $X_i\to X'_i$ 为同构,
故此时该图也为笛卡尔图 (在 $X_i$ 或 $X'_i$ 中检验等式是同一回事)
\nde{引理 \hyperref[Vbis.lem3.3.3.3]{3.3.3.3} 的证明略作了改写.}.

因此由 \hyperref[Vbis.prop3.3.1]{3.3.1},
相关态射均为普遍 $G$-$1$-上同调下降态射.
令 $[X'/X]^p$ 为 $X'$ 在 $X$ 上的 $(p+1)$ 重纤维积.
由 \hyperref[Vbis.lem3.3.3.3]{3.3.3.3},
只须在基变换 $[X'/X]^p\to X$ 后验证 \hyperref[Vbis.lem3.3.3.2]{3.3.3.2}.
作此基变换后, \hyperref[Vbis.lem3.3.3.2]{3.3.3.2} 的假设仍成立,
且 $f_{n+1}$ 还容许截面. 引理 \hyperref[Vbis.lem3.0.2.4]{3.0.2.4} 遂完成证明.

对于有效下降, 有下述结果.

\begin{enonce*}{命题 3.3.4}\etiquette[3.3.4]{Vbis.prop3.3.4}
设 $X:\Delta^{\circ}\to B/S$ 为函子.
假定 $X_0\to S$ 为普遍 $G^{\circ}$-$2$-下降态射,
且对 $n=0,1$, 态射 $X_{n+1}\to(\cosk_n(X))_{n+1}$ 也如此.
则函子 $X$ 是 $G^{\circ}$-$2$-忠实的, 并在任意基变换后仍然如此
(换言之, $X$ 定义 \hyperref[Vbis.def2.4.11]{2.4.11} 意义下的普遍有效下降增广).
\end{enonce*}

由 \cite{sga4.exp.Vbis.bib.4} (7.12), 只须证明 $i_2^{\ast}(X)$ 为 $G^{\circ}$-$2$-忠实;
为此使用下述引理.

\begin{enonce*}{引理 3.3.4.1}\etiquette[3.3.4.1]{Vbis.lem3.3.4.1}
设\pageoriginale{} $X$ 与 $X'$ 为两个函子 $\mathop{\Delta}_2^{\circ}\to B/S$,
并考虑交换图
\[
\xymatrix@R=1.5cm@C=2cm{
X'_2 \ar[r]\ar@<8pt>[r]\ar@<-8pt>[r]\ar[d]_{k}&X'_1\ar@<4pt>[l]
\ar@<-4pt>[l]\ar@<4pt>[r]\ar@<-4pt>[r]\ar@{=}[d]&X'_0\ar[l]
\ar[r]\ar@{=}[d]&S\ar@{=}[d]\\
X_2 \ar[r]\ar@<8pt>[r]\ar@<-8pt>[r]&X_1\ar@<4pt>[l]
\ar@<-4pt>[l]\ar@<4pt>[r]\ar@<-4pt>[r]&X_0\ar[l]\ar[r]&S
}
\]
使 $k$ 为 $G^{\circ}$-$0$-下降态射.
若 $X$ 为 $G^{\circ}$-$2$-忠实, 则 $X'$ 也如此.
\end{enonce*}

这是显然的.

\begin{enonce*}{引理 3.3.4.2}\etiquette[3.3.4.2]{Vbis.lem3.3.4.2}
设 $X$ 与 $X'$ 为两个函子 $\mathop{\Delta}_2^{\circ}\to B/S$,
并考虑交换图
\def\labelstyle{\scriptstyle}
\[
\xymatrix@R=2cm@C=2cm{
X'_2 \ar[r]^{\partial^{1}} \ar@<17pt>[r]^{\partial^{0}}
\ar@<-14pt>[r]_{\partial^{2}} \ar[d]_{f_2}&X'_1\ar@<8pt>[l]
\ar@<-11pt>[l]\ar@<11pt>[r]^{\partial^{0}}\ar@<-8pt>[r]_{\partial^{1}}
\ar[d]^{f_1}&X'_0\ar[l]_{\sigma^{0}}\ar[r]\ar@{=}[d]&S\ar@{=}[d]\\
X_2 \ar[r]\ar@<12pt>[r]\ar@<-12pt>[r]&X_1\ar@<6pt>[l]
\ar@<-6pt>[l]\ar@<6pt>[r]\ar@<-6pt>[r]&X_0\ar[l]\ar[r]&S
}
\]
使 $f_1$ 为 $G^{\circ}$-$1$-下降态射, $f_2$ 为 $G^{\circ}$-$0$-下降态射.
再假定 $X'_2\xrightarrow{\sim}(\cosk_1(X'))_2$.
若 $X$ 为 $G^{\circ}$-$2$-忠实, 则 $X'$ 也如此.
\end{enonce*}

令 $X''_1=\fprod{X'_1}{X'_1}{X_1}$,
并设 $\Ysheaf{\partial}_0:X''_1\to X'_1$ 和
$\Ysheaf{\partial}_1:X''_1\to X'_1$ 为两个投影.

余骨架的定义给出箭头 $\varphi:X''_1\to X'_2$, 使
\begin{equation*}
\begin{cases}
\partial_0\circ\varphi=\Ysheaf{\partial}_0,\\
\partial_1\circ\varphi=\Ysheaf{\partial}_1,\\
\partial_2\circ\varphi=\sigma_0\partial_1\Ysheaf{\partial}_0
=\sigma_0\partial_1\Ysheaf{\partial}_1{\quoi}.
\end{cases}
\end{equation*}

现取 $X'$ 上的一个下降资料:
它给出 $X_0$ 上的一个对象、$X_1$ 上的两个对象以及它们在 $X'_1$ 上的一个同构.
箭头 $\varphi$ 表明, 这些同构在 $X''_1$ 上的两个逆像相等.
由于 $f_1$ 为 $1$-下降态射, 得到 $X_1$ 上两个对象之间的同构;
此同构为下降资料 (因为 $X'_2\to X_2$ 为 $0$-下降态射), 结论随之得出.

\begin{enonce*}{推论 3.3.5}\etiquette[3.3.5]{Vbis.cor3.3.5}
假定 \emph{$B$ 的所有箭头均平坦}.
$S$ 关于 $G$-$2$-上同调下降拓扑的每个对象可表的超覆盖,
都定义一个普遍 $G$-$2$-上同调下降增广.
\end{enonce*}

% END INCLUDED SGA4_English_current_buildable_source_20260906/source/05bis/drafts/SGA4_VBIS_LINES_2586_2816_English.texfrag


% BEGIN INCLUDED SGA4_English_current_buildable_source_20260906/source/05bis/drafts/SGA4_VBIS_LINES_2817_2992_English.texfrag
\section{例子}\etiquette[4]{Vbis.sec4}

\subsection{拓扑空间上的阿贝尔群层}\etiquette[4.1]{Vbis.subsec4.1}

\setcounter{subsubsection}{-1}
\subsubsection{}\etiquette[4.1.0]{Vbis.subsubsec4.1.0}
本节用 Top 表示拓扑空间范畴 [固定宇宙 $\Uu$ 的元素].
定义一个在 Top 上为拓扑斯对偶双纤维化的范畴 $\mathcal{E}$:
对象为二元组 $(X, F)$, 其中 $X$ 为拓扑空间, $F$ 为 $X$ 上的集合层;
态射为二元组 $(f, \varphi) : (X, F) \to (Y, H)$,
其中 $f : X \to Y$ 为连续映射, $\varphi : H \to f_{\ast} (F)$ 为层态射.
取 $\mathcal{O}$ 为 $\mathcal{E}^0$ 在 $\Top^0$ 上将每个拓扑空间对应到
常值层 $\ZZ_X$ 的截面: 令 $G = \Ysheaf{\Mod}(\mathcal{E}^0, \mathcal{O})$,
于是对每个拓扑空间 $X$, $G_X$ 等同于 $X$ 上的阿贝尔群层范畴.

回顾, 若 $\fprod{X}{X}{Y}$ 中的对角线闭,
就称态射 $f : X \to Y$ \emph{分离}.
\emph{固有}态射是分离且普遍闭的态射 (注意此定义比 Bourbaki 的定义更严格).

下述“基变换定理”的\pageoriginale{}证明受到 (\cite{sga4.exp.Vbis.bib.5} II(4.11.1)) 的启发.

\begin{enonce*}{定理 4.1.1}\etiquette[4.1.1]{Vbis.thm4.1.1}
设 $f : X \to Y$ 为固有态射, $F$ 为 $X$ 上的阿贝尔层, 并考虑笛卡尔图
\[
\xymatrix@C=2.5cm@R=1.5cm{
X' \ar[r]^{g'} \ar[d]^{f'}& X\ar[d]^{f}\\
Y' \ar[r]^{g}& Y.
}
\]
则典范映射 \emph{(XII \hyperref[XII.eq4.2]{4.2})}
或 \emph{\eqhyperref[Vbis.eq3.2.2.3]{3.2.2.3}}
$$
g^{\ast} R^i f_{\ast} F \xrightarrow{\sim} R^i f'_{\ast} {g'}^{\ast}F
$$
为同构\nde{与前文不同, 这里是通常函子的导出函子,
而不是由纤维范畴结构 $G^\circ\rightarrow\Top^\circ$ 产生的导出函子.}.
\end{enonce*}

此定理等价于下述推论 (取 $Y'= (\Point)$ 得到推论; 对 $f$ 和 $f'$ 应用推论得到定理).

\begin{enonce*}{推论 4.1.2}\etiquette[4.1.2]{Vbis.cor4.1.2}
设 $f : X \to Y$ 为固有态射, $F$ 为 $X$ 上的阿贝尔层, $y \in Y$.
典范映射
$$
(R^i f_{\ast} F)_y \longrightarrow H^i(f^{-1} (y), F)
$$
为双射.
\end{enonce*}

按定义, 当 $U$ 遍历 $y$ 的邻域时, 有
$$
(R^if_{\ast} F)_y = \varinjlim H^i(f^{-1} (U), F){\quoi}.
$$
由于 $f$ 闭, 各 $f^{-1}(U)$ 构成 $f^{-1}(y)$ 的一个基本邻域系,
结论来自下述引理 (\hyperref[Vbis.lem4.1.3]{4.1.3}).

\begin{enonce*}{引理 4.1.3}\etiquette[4.1.3]{Vbis.lem4.1.3}
设\pageoriginale{} $X$ 为拓扑空间, $K \subset X$, $F$ 为 $X$ 上的阿贝尔层.
假定 $K$ 紧, 且 $K$ 的任意两个不同点在 $X$ 中有不交邻域.
则当 $U$ 遍历 $K$ 的邻域时, 有
\begin{equation*}
\varinjlim H^i(U, F) \xrightarrow{\sim} H^i(K, F){\quoi}.\tag{4.1.3.1}
\end{equation*}\etiquette[4.1.3.1]{Vbis.eq4.1.3.1}
\end{enonce*}

先处理 $i = 0$ 的情形.
此时 \eqhyperref[Vbis.eq4.1.3.1]{4.1.3.1} 显然单射.
为证满射性, 使用下述结果.

\begin{enonce*}{引理 4.1.4}\etiquette[4.1.4]{Vbis.lem4.1.4}
在 \emph{(\hyperref[Vbis.lem4.1.3]{4.1.3})} 的假设下,
若 $A$ 和 $B$ 为 $K$ 的两个闭子集, $W$ 为 $A \cap B$ 在 $X$ 中的邻域,
则 $A$ 与 $B$ 在 $X$ 中存在邻域 $U$ 和 $V$, 使 $U \cap V \subset W$.
\end{enonce*}

先处理 $A$ 只含单点 $a$ 的情形. 若 $a \in B$, 断言显然 (取 $U = W$).
若 $a \not \in B$, 则对每个 $b \in B$, $a$ 与 $b$ 在 $X$ 中有不交开邻域 $U_b$ 和 $V_b$.
取 $V$ 为各 $V_b$ 中覆盖 $B$ 的有限多个的并,
取 $U$ 为相应各 $U_b$ 的交.

在一般情形中, 对每个 $a \in A$, $a$ 与 $B$ 在 $X$ 中存在开邻域 $U_a$ 和 $V_a$,
使 $U_a \cap V_a \subset W$.
取 $U$ 为各 $U_a$ 中覆盖 $A$ 的有限多个的并,
取 $V$ 为相应各 $V_a$ 的交.

回到 (\hyperref[Vbis.lem4.1.3]{4.1.3}).
若 $s \in H^{\circ}(K, F)$, 则在 $X$ 中存在 $K$ 的开覆盖 $U_i$,
以及截面 $s_i \in H^{\circ}(U_i, F)$, 使在 $U_i \cap K$ 上 $s_i = s$.
可假定只有有限多个 $U_i$.
设 $(K_i)$ 为 $K$ 的闭覆盖, 满足 $K_i \subset K \cap U_i$.
设 $W_y$ 为 $U_i \cap U_j$ 中使 $s_i = s_j$ 的开子集;
有 $K_i \cap K_j \subset W_{ij}$.
对所有二元组 $(K_i, K_j)$ 及各 $W_{ij}$ 应用 (\hyperref[Vbis.lem4.1.4]{4.1.4}),
得到开子集 $V_{ij}$, 使 $K_i \subset V_{ij} \subset U_i$
且 $V_{ij} \cap V_{ji} \subset W_{ij}$.
令 $U'_i = \displaystyle\mathop{\cap}_j V_{ij}$:
则 $K_i \subset U'_i \subset U_i$, 且在 $U'_i \cap U'_j$ 上 $s_i = s_j$.
因此各 $s_i$ 可在由各 $U'_i$ 的并给出的 $K$ 的邻域上粘合,\pageoriginale
且当 $i = 0$ 时, \eqhyperref[Vbis.eq4.1.3.1]{4.1.3.1} 满射 (从而双射)
\nde{不要混淆符号 $i$ 的两种含义.}.

\begin{enonce*}{引理 4.1.5}\etiquette[4.1.5]{Vbis.lem4.1.5}
若 $F$ 松弛, 则 $K$ 上的层 $F | K$ 为软层.
\end{enonce*}

设 $A \subset K$ 在 $K$ 中闭. 由上述讨论, $F$ 在 $A$ 上的每个截面可延拓到某邻域.
由于 $F$ 松弛, 此截面可延拓到 $X$, 因而特别可延拓到 $K$.

现证明 (\hyperref[Vbis.lem4.1.3]{4.1.3}). 设 $F^{\ast}$ 为 $F$ 的松弛分解. 有
\begin{eqnarray*}
\varinjlim H^i (U, F) &\stackrel{\ndemark}{=}& \varinjlim H^i (\Gamma(U, F^{\ast})) = H^i\varinjlim \Gamma(U,F^{\ast})
\stackrel{\ndemark}{=} H^i \Gamma (K, F^{\ast})\\
& \stackrel{\text{\hyperref[Vbis.lem4.1.5]{(4.1.5)}}}{=} & H^i (K, F){\quoi}.
\end{eqnarray*}
\addtocounter{cnote}{-1}\ndetext{$H^i (\Gamma(U, F^{\ast}))$ 在此表示复形 $\Gamma(U,F^\ast)$ 的第 $i$ 个上同调群.}\addtocounter{cnote}{1}
\ndetext{由 \hyperref[Vbis.lem4.1.3]{(4.1.3)}, 取 $\Gamma=H^0$}
\begin{enonce*}{推论 4.1.6}\etiquette[4.1.6]{Vbis.cor4.1.6}
设 $f : X \to Y$ 为拓扑空间之间的固有满射态射.
则 $f$ 为 $G$-$1$-上同调下降态射.
\end{enonce*}

这来自 (\hyperref[Vbis.cor4.1.2]{4.1.2}) 和 (\hyperref[Vbis.prop3.2.4]{3.2.4}).

\begin{enonce*}{推论 4.1.7}\etiquette[4.1.7]{Vbis.cor4.1.7}
设 $\Uu = (U_i)_{i \in I}$ 为拓扑空间 $X$ 的局部有限闭覆盖.
则典范态射 $f : \coprod_{i\in I} U_i \to X$ 为 $G$-$1$-上同调下降态射.
\end{enonce*}

事实上, $f$ 固有且满射.
应用 (\hyperref[Vbis.prop2.5.6]{2.5.6}), 便重新得到
[\cite{sga4.exp.Vbis.bib.5}. II (5.2.4)] 的谱序列 (参见引言).

\begin{enonce*}{命题 4.1.8}\etiquette[4.1.8]{Vbis.prop4.1.8}
设 $f : X \to Y$ 为拓扑空间之间的满射局部同胚.
则 $f$ 为普遍 $G$-$1$-上同调下降态射.
\end{enonce*}

在 $Y$ 上局部化 (参见 (\hyperref[Vbis.prop3.2.4]{3.2.4})), 可归结为 $f$ 有截面的情形.

\begin{enonce*}{推论 4.1.9}\etiquette[4.1.9]{Vbis.cor4.1.9}
设\pageoriginale{} $\Uu = (U_i)_{i \in I}$ 为拓扑空间 $X$ 的开覆盖.
则典范态射 $f : \coprod_{i \in I} U_i \to X$ 为普遍 $G$-$1$-上同调下降态射.
\end{enonce*}

下降谱序列遂给出 [\cite{sga4.exp.Vbis.bib.5} II (5.4.1)] (参见引言).

\begin{enonce*}[remark]{注 4.1.10}\etiquette[4.1.10]{Vbis.rem4.1.10}
%%%%%%%%%%%%%%%%%%%
%%%%%%%%%%%%Check that this compiles
%%%%%%%%%%%%%%%%%%%%
(\hyperrefext[VIII][VIII.sec9]{9}) 中的证明表明,
上述所有命题中的“1-上同调下降”均可换为“2-上同调下降”.
\end{enonce*}

\subsection{概形上的拟凝聚模}\etiquette[4.2]{Vbis.subsec4.2}

\setcounter{subsubsection}{-1}

% END INCLUDED SGA4_English_current_buildable_source_20260906/source/05bis/drafts/SGA4_VBIS_LINES_2817_2992_English.texfrag


% BEGIN INCLUDED SGA4_English_current_buildable_source_20260906/source/05bis/drafts/SGA4_VBIS_LINES_2993_3150_English.texfrag
\subsubsection{}\etiquette[4.2.0]{Vbis.subsubsec4.2.0}
设 $(\Sch)_s$ 为对象是属于 $\Uu$ 的概形、箭头是拟紧且拟分离态射的范畴.
定义一个在 $(\Sch)_s$ 上为拓扑斯对偶双纤维化的范畴 $\mathcal{E}$:
对象为二元组 $(X,F)$, 其中 $X$ 为概形, $F$ 为 $X$ 上的层 (关于 Zariski 拓扑);
态射为二元组 $(f,\varphi):(X,F)\to(Y,G)$,
其中 $f:X\to Y$ 为连续映射, $\varphi:G\to f_{\ast}(F)$ 为层态射.
取 $\mathcal{O}$ 为 $\mathcal{E}^0$ 在 $(\Sch)^0_s$ 上
将每个概形对应到其结构层的截面;
取 $G$ 为 $\Ysheaf{\Mod}(\mathcal{E}^0,\mathcal{O})$ 的子范畴,
使对每个概形 $X$, $G_X$ 为 $X$ 上的拟凝聚模范畴.

\begin{enonce*}{命题 4.2.1}\etiquette[4.2.1]{Vbis.prop4.2.1}
每个忠实平坦、拟紧且拟分离的态射 $f:X\to Y$ 都是普遍 $G$-$2$-上同调下降态射.
\end{enonce*}

由 [\SGA I (VIII.5.2)], $f$ 为有效 $G^0$-下降态射.
由 [\EGA IV (1.7.21)], 可对笛卡尔方块
\[
\xymatrix@C=3cm@R=1.5cm{
  \fprod{X}{X}{Y} \ar[r] \ar[d]& X \ar[d]^{f}\\
X   \ar[r]_{f}            & Y
}
\]
应用 (\hyperref[Vbis.prop3.2.4]{3.2.4}).

\subsection{概形上的平展层}\etiquette[4.3]{Vbis.subsec4.3}
\pageoriginale

\setcounter{subsubsection}{-1}
\subsubsection{}\etiquette[4.3.0]{Vbis.subsubsec4.3.0}
设 $(\Sch)$ 为所有属于 $\Uu$ 的概形的范畴.
于是定义一个在 $(\Sch)$ 上为拓扑斯对偶双纤维化的范畴 $\mathcal{E}$:
对象为二元组 $(X,\mathcal{F})$, 其中 $X$ 为概形, $\mathcal{F}$ 为 $X$ 上的平展层;
态射为二元组 $(f,\varphi):(X,\mathcal{F})\to(Y,\mathcal{G})$,
其中 $f:X\to Y$ 为概形态射, $\varphi:\mathcal{G}\to f_{\ast}(\mathcal{F})$ 为层态射.

设 $n$ 为 $\geq0$ 的整数.
用 $\mathcal{O}_n$ 表示 $\mathcal{E}^0$ 在 $(\Sch)^0$ 上
将每个概形 $X$ 对应到常值层 $(\ZZ/n\ZZ)_X$ 的截面.
令 $F_n=\Ysheaf{\Mod}(\mathcal{E}^0,\mathcal{O}_n)$,
于是对每个概形 $X$, $F_{n_X}$ 等同于 $X$ 上的 $\ZZ/n\ZZ$-模层范畴.

\subsubsection{}\etiquette[4.3.1]{Vbis.subsubsec4.3.1}
现设 $(\Sch)_s$ 为 $(\Sch)$ 的子范畴, 对象是属于 $\Uu$ 的概形,
箭头为拟紧且拟分离的态射.
用 $\mathcal{E}_s$ 表示经基变换 $(\Sch)_s\to(\Sch)$ 从 $\mathcal{E}$ 得到的
在 $(\Sch)_s$ 上为拓扑斯对偶双纤维化的范畴;
仍用 $\mathcal{O}_0$ 表示 $\mathcal{O}_0$ 在 $(\Sch)_s$ 上的限制.
设 $G\subset\Ysheaf{\Mod}(\mathcal{E}^0_s,\mathcal{O}_0)$
为 $\Ysheaf{\Mod}(\mathcal{E}^0_s,\mathcal{O}_0)$ 的纤维且余纤维子范畴,
使对每个概形 $X$, $G_X$ 为 $X$ 上的挠阿贝尔层范畴.
读者会注意到条件 (\hyperref[Vbis.par2.4.1.1]{2.4.1.1}) 满足.

\begin{enonce*}{命题 4.3.2}\etiquette[4.3.2]{Vbis.prop4.3.2}
设 $f:X\to Y$ 为\emph{固有满射}态射.
则对 $n\geq1$, $f$ 为普遍 $F_n$-$2$-上同调下降态射.
将其视为 $(\Sch)_s$ 的态射时, $f$ 为普遍 $G$-$2$-上同调下降态射.
\end{enonce*}

由\pageoriginale{} (\hyperrefext[VIII][VIII.thm9.4]{9.4}), $f$ 为有效 $F^0$-下降态射,
且由 (XII.\hyperref[XII.thm5.1]{(5.1)}), 可对笛卡尔图
\[
\xymatrix@C=3cm@R=1.5cm{
  \fprod{X}{X}{Y} \ar[r] \ar[d]& X \ar[d]^{f}\\
X   \ar[r]_{f}            & Y.
}
\]
应用 (\hyperref[Vbis.prop3.2.4]{3.2.4}).

\begin{enonce*}{命题 4.3.3}\etiquette[4.3.3]{Vbis.prop4.3.3}
设 $f:X\to Y$ 为概形的\emph{满射}态射. 假定下述假设之一满足:
\begin{enumerate}
\renewcommand{\theenumi}{\alph {enumi}}
\renewcommand{\labelenumi}{\theenumi)}
\item $f$ 整.

\item $Y$ 局部 Noether, 且 $f$ 普遍开且局部有限表现.

\item $f$ 平坦且局部有限表现.
\end{enumerate}
则对每个 $n\geq0$, $f$ 为 $F_n$-$2$-上同调下降态射.
\end{enonce*}

\begin{enumerate}
\renewcommand{\theenumi}{\alph {enumi}}
\renewcommand{\labelenumi}{\theenumi)}
\item 重复 (\hyperref[Vbis.prop4.3.2]{4.3.2}) 的证明,
用 ((\hyperrefext[VIII][VIII.cor5.6]{5.6})) 代替 (XII.\hyperref[XII.thm5.1]{(5.1)}).

\item 由于问题在 $Y$ 上是局部的, 可设 $Y$ 为 Noether 概形.
应用 (\EGA IV (14.5.10)). 由 (\hyperref[Vbis.prop3.3.1]{3.3.1}) 和 a),
归结为 $Y$ 的每点 $y$ 有邻域 $U$ 使 $f^{-1}(U)\to U$ 有截面的情形,
再由 (\hyperref[Vbis.prop3.2.4]{3.2.4}) 得到结论.

\item 须证明 $f$ 为 \emph{$F$-$1$-上同调下降态射}
(由 ((\hyperrefext[VIII][VIII.thm9.4]{9.4})), 已知 $f$ 为有效 $F^0$-下降态射).
\end{enumerate}

在 $Y$ 上局部化 (参见 (\hyperref[Vbis.prop3.2.4]{3.2.4})),
可设 $Y$ 为环为 $A$ 的仿射概形.
于是可写 $\Spec A=\varprojlim_{i\in I}\Spec A_i$,
其中 $I$ 为滤过有序集, 每个 $i$ 对应的 $A_i$ 为 Noether 环,
并有 $X=\varprojlim X_i$, 其中对每个 $i$,
$X_i$ 为 $\Spec A_i$ 上\emph{平坦}且\emph{局部有限表现}的概形,
使 $f_i:X_i\to\Spec A_i$ \emph{满射}; 此外, 图
\[
\xymatrix@C=3cm@R=1.5cm{
X \ar[r]\ar[d]_{f}& X_i \ar[d]^{f_i}\\
Y \ar[r]^{u_i}& \Spec A_i
}
\]
为\pageoriginale{}笛卡尔图.

设 $\mathfrak{F}$ 为 $Y$ 上的阿贝尔层.
有 $\mathfrak{F}={\varinjlim}_i u^{\ast}_i u_{i\ast}(\mathfrak{F})$,
故由标准的取极限过程, 归结为 $\mathfrak{F}$ 来自某个 $i$ 的
$\Spec A_i$ 上的层的情形, 再由 b) 得到结论.

\begin{enonce*}[remark]{注 4.3.4}\etiquette[4.3.4]{Vbis.rem4.3.4}
若 $f:X\to Y$ 整, 则当 $q>0$ 时 $R^qf_{\ast}=0$.
容易得出, 通过 $f$ 作下降时不必限于下有界复形.
\end{enonce*}

\begin{enonce*}{推论 4.3.5}\etiquette[4.3.5]{Vbis.cor4.3.5}
设 $(X_{\alpha}\xrightarrow{u_\alpha}X)_{\alpha\in A}$ 为平展拓扑的覆盖族.
则对每个 $n\geq0$, $\coprod_{\alpha\in A}X_{\alpha}\to X$
为普遍 $F_n$-$2$-上同调下降态射.
\end{enonce*}

事实上, 存在分解
\[
\xymatrix@=1.5cm{
  &\displaystyle{\mathop{\coprod}_{a\in A} X_{\alpha}} \ar[rd]&\\
X \ar[ru]         \ar[rr]^{h} &                              & X
}
\]
其中 $h$ 为满射平展态射.


% END INCLUDED SGA4_English_current_buildable_source_20260906/source/05bis/drafts/SGA4_VBIS_LINES_2993_3150_English.texfrag


% BEGIN INCLUDED SGA4_English_current_buildable_source_20260906/source/05bis/drafts/SGA4_VBIS_LINES_3151_3452_English.texfrag
\section{应用}\etiquette[5]{Vbis.sec5}\pageoriginale

\subsection{单纯对象的构造}\etiquette[5.1]{Vbis.subsec5.1}

\setcounter{subsubsection}{-1}
\subsubsection{}\etiquette[5.1.0]{Vbis.subsubsec5.1.0}

直到并包括 \hyperref[Vbis.prop5.1.3]{5.1.3},
设 $H$ 为非空有限投射极限存在、有限余积存在且不交并普遍的范畴.
若 $X:\Delta^0\to H$ 为 $H$ 的单纯对象,
对 $i\in[k]$ 和 $0\leq k$, 用 $d^k_i:X_k\to X_{k-1}$
(分别用 $s^k_i:X_k\to X_{k+1}$) 表示对应于面算子
$\partial^i_k:[k-1]\to[k]$ (分别为退化算子 $\sigma^i_k:[k]\to[k+1]$) 的箭头
(参见 \cite{sga4.exp.Vbis.bib.3} II 2).

\begin{enonce*}{定义 5.1.1}\etiquette[5.1.1]{Vbis.def5.1.1}
若存在 $X_j$ 的一族子对象 $NX_j$ $(j\geq0)$, 使态射
$$
h_i:\coprod_{\substack{\Hom_{\Delta^{-}}([i], [j])\\
j\leq i}} NX_j \longrightarrow X_i
$$
为同构, 就称 $H$ 的单纯对象 $X:\Delta^0\to H$ 是 $\sigma$-分裂的.
\end{enonce*}

同样地, 若存在 $NX_j$ $(0\leq j\leq k)$,
使对 $i\leq k$ 满足上述条件,
就称 $k$-截断单纯对象 $X:\Delta^0_k\to H$ 是 $\sigma$-分裂的.

各 $NX_j$ 由 $X$ 唯一确定.
事实上, 若 $X$ 为 $\sigma$-分裂,
则退化算子 $s^{k-1}_\ell:X_{k-1}\to X_k$
是从 $X_{k-1}$ 到 $X_k$ 的直和项的同构, 并且
$$
NX_k = \mathop{\cap}_\ell\quad (s^{k-1}_\ell(X_{k-1}) \text{ 的补对象}) {\quoi}.
$$
当 $k=0$ 时, $NX_0=X_0$.

\setcounter{subsubsection}{1}
\subsubsection{}\etiquette[5.1.2]{Vbis.subsubsec5.1.2}
对 $H$ 的 $\sigma$-分裂 $(n+1)$-截断单纯对象 $X$,
用 $\alpha_n(X)$ 表示由下列资料组成的三元组:
\begin{enumerate}
\renewcommand{\theenumi}{\alph {enumi}}
\renewcommand{\labelenumi}{\theenumi)}
\item $X$ 在 $\Delta^0_n$ 上的限制 $i^{\ast}_n(X)$;

\item $NX_{n+1}$;\pageoriginale

\item 从 $NX_{n+1}$ 到 $(\cosq_n(X))_{n+1}$ 的显然映射.
\end{enumerate}

此三元组 $(Y,N,\beta)$ 满足下述条件:

\noindent
$(\ast)$ $Y$ 为 $H$ 的 $\sigma$-分裂 $n$-截断单纯对象,
$\beta$ 为从 $N$ 到 $(\cosq Y)_{n+1}$ 的映射.

\begin{enonce*}{命题 5.1.3}\etiquette[5.1.3]{Vbis.prop5.1.3}
设 $(Y,N,\beta)$ 为满足上述 $(\ast)$ 的三元组.
\begin{enumerate}
\renewcommand{\theenumi}{\roman {enumi}}
\renewcommand{\labelenumi}{(\theenumi)}
\item 精确到唯一同构, 存在唯一的 $\sigma$-分裂
$X:\Delta^0_{n+1}\to H$, 使 $\alpha_n(X)\simeq(Y,N,\beta)$.

\item 给出从 $X$ 到截断单纯对象 $Z:\Delta^0_{n+1}\to H$ 的态射 $f$,
等价于给出
\begin{enumerate}
\renewcommand{\labelenumii}{\theenumii)}
\item 态射 $f':Y\to i^{\ast}_n(Z)$;

\item 态射 $f'':N\to Z_{n+1}$, 使图
\[
\xymatrix@C=4cm@R=2cm{
N  \ar[r] \ar[d]^{f''}& (\cosq_n Y)_{n+1} \ar[d]^{f'}\\
Z_{n+1} \ar[r]& (\cosq_n Z)_{n+1}
}
\]
交换.
\end{enumerate}
\end{enumerate}
\end{enonce*}

现在构造 $X$.

令 $Y'_{n+1}=(\cosk_nY)_{n+1}$, 并用
$$
{d'}^{n+1}_i:Y'_{n+1}\to Y_n
\quad(\textup{分别为 }{s'}^n_i:Y_n\to Y'_{n+1})
$$
表示由函子性得到的面算子 (分别为退化算子).

\noindent \sisi{$+$}{} 令
$Y_{n+1}=N\coprod(\coprod_{\Hom_{\Delta^{-}_{\ell\leq n}}([n+1], [\ell])}NY_{\ell})$,
并设 $\alpha:Y_{n+1}\to Y'_{n+1}$ 为其分量是 $\beta$
以及对每个满足 $\ell\leq n$ 的满态射 $[n+1]\to[\ell]$
相应的箭头 $NY_{\ell}\to Y_{\ell}\to Y'_{n+1}$ 的箭头.

\noindent \sisi{$+$}{} 设\pageoriginale{}
$s^n_i:Y_n\simeq\coprod_{\Hom_{\Delta^{-}_{\ell'\leq n}}([n],[{\ell}'])}NY_{{\ell}'}\to Y_{n+1}$
为 $\sigma^i_n$ 诱导的箭头, 于是
\begin{enumerate}
\renewcommand{\theenumi}{\Roman{enumi}}
\renewcommand{\labelenumi}{(\theenumi)}
\item $\alpha\circ s^n_i={{s}'}^n_i$.

令

\item $d^{n+1}_i={d'}^{n+1}_i\circ\alpha$.
\end{enumerate}

\noindent \sisi{$+$}{} 为证这定义 $\Hom(\Delta^0_{n+1},H)$ 的对象,
由 [\cite{sga4.exp.Vbis.bib.3}. II (2.4)], 只须验证公式
\begin{enumerate}
\renewcommand{\theenumi}{\alph{enumi}}
\renewcommand{\labelenumi}{\theenumi)}
\item $d^n_i\cdot d^{n+1}_j=d^n_{j-1}\cdot d^{n+1}_i$ \qquad $i<j$;

\item $s^n_i\cdot s^{n-1}_j=s^n_{j+1}\cdot s^{n-1}_i$ \qquad $i\leq j$;

\item $d^{n+1}_i\cdot s^n_j=
\begin{cases}
s^{n-1}_{j-1}\cdot d^n_i & i<j,\\
\id & i=j \mbox{ 或 } i=j+1,\\
s^{n-1}_{j}\cdot d^n_{i-1} & i>j+1.
\end{cases}$
\end{enumerate}

公式 a) 和 c) 由 (I) 与 (II) 显然成立,
因为已知将 $d^{n+1}_i$ 换为 ${d'}^{n+1}_i$、$s^n_i$ 换为 ${s'}^n_i$ 时它们成立.

公式 b) 按构造成立.

所得对象 $X$ 满足 \hyperref[Vbis.prop5.1.3]{5.1.3} 的验证留待读者完成.

\noindent \sisi{$+$}{} 注意, \hyperref[Vbis.prop5.1.3]{5.1.3} 的证明中构造的,
从满足 (\hyperref[Vbis.subsubsec5.1.2]{5.1.2}) $(\ast)$ 的三元组范畴
到 $(n+1)$-截断单纯对象范畴的函子忠实但并不满.

命题 \hyperref[Vbis.prop5.1.3]{5.1.3} 提供了一种方便的归纳构造 $H$ 的单纯对象的方法.
本节余下部分, 将此观察形式化为 \hyperref[Vbis.subsec5.2]{5.2}
和 \hyperref[Vbis.subsec5.3]{5.3} 中将使用的形式.

\setcounter{subsubsection}{3}
\subsubsection{}\etiquette[5.1.4]{Vbis.subsubsec5.1.4}

设\pageoriginale{} $E\xrightarrow{\pi}B$ 为函子,
使对 $B$ 的每个对象 $b$, 纤维范畴 $E_b$ 非空且满足 \hyperref[Vbis.subsubsec5.1.0]{5.1.0} 的条件.

设 $Q$ 为 $E$ 的对象的一个在同构下稳定的性质. 假定下述条件成立:
\begin{enumerate}
\renewcommand{\theenumi}{\alph{enumi}}
\renewcommand{\labelenumi}{\theenumi)}
\item \emph{对 $B$ 的每个对象 $b$, 存在 $E_b$ 的对象 $x$ 满足 $Q$.}

\item \emph{对 $B$ 的每个对象 $b$ 及 $E_b$ 中满足 $Q$ 的每对对象 $(x,y)$,
存在 $E_b$ 的对象 $z$ 满足 $Q$, 并在 $E_b$ 中有态射 $z\to x$、$z\to y$.}

\item \emph{对 $B$ 的每个态射 $b'\xrightarrow{t}b$ 及 $E_b$ 中满足 $Q$ 的每个对象 $x$,
存在 $E_{b'}$ 的对象 $x'$ 满足 $Q$, 并在 $t$ 上有态射 $x'\to x$.}

另一方面, 设 $P$ 为 $E$ 中恒等态射上方的箭头的一个在同构下稳定的性质.
假定下述条件成立:

\item \emph{对 $B$ 的每个对象 $b$ 及 $E_b$ 的每个对象 $x$,
存在 $E_b$ 中的箭头 $y\to x$ 满足 $P$.}

\item \emph{对 $B$ 的每个对象 $b$, 以及 $E_b$ 中形如}
\[
\xymatrix{
u \ar[rrrr]^{f} & & &                 &x\\
                & & &z\ar[ru]\ar[rd] &\\
v \ar[rrrr]_{g} & & &               & y
}
\]
\emph{且 $f$ 和 $g$ 满足 $P$ 的每个图, 可以在 $E_b$ 中补成交换图}
\[
\xymatrix@R=1.25cm{
  & u\ar[rrrr]^{f} & & &  & x\\
w\ar[rrrr]^{h}\ar[ru]\ar[rd] &   & & &z\ar[ru]\ar[rd] & \\
  & v\ar[rrrr]_{g} & & &  & y
}
\]
\emph{其中}\pageoriginale{} $h$ \emph{满足 $P$.}

\item \emph{对 $B$ 的每个态射 $b'\xrightarrow{t}b$
及 $E_b$ 中满足 $P$ 的每个态射 $x\xrightarrow{f}y$,
存在 $E_{b'}$ 中满足 $P$ 的态射 $x'\xrightarrow{f'}y'$ 以及交换图}
\[
\xymatrix@C=3cm@R=1.5cm{
x' \ar[r]^{\alpha} \ar[d]_{f'}& x\ar[d]^{f}\\
y' \ar[r]_{\beta}& y
}
\]
\emph{其中 $\alpha$ 和 $\beta$ 为 $t$ 上的态射.}
\end{enumerate}

\begin{enonce*}{定义 5.1.5}\etiquette[5.1.5]{Vbis.def5.1.5}
采用 \hyperref[Vbis.subsubsec5.1.4]{5.1.4} 的记号,
若下述条件成立, 就称 $\Hom(\Delta^0,E)$ 的对象 $X$ 满足条件 \emph{(APQ)}:
\begin{enumerate}
\item[(i)] $\pi\circ X$ 的像是由 $B$ 的对象 $b$ 定义的常值单纯对象.

\item[(ii)] 作为 $\Hom(\Delta^0,E_b)$ 的对象, $X$ 是 $\sigma$-分裂的.

\item[(iii)] $X_0$ 满足条件 $Q$.

\item[(iv)] 对每个整数 $n$, 箭头 $NX_{n+1}\to(\cosq_nX)_{n+1}$ 满足条件 $P$.
\end{enumerate}
\end{enonce*}

\setcounter{subsubsection}{5}
\subsubsection{}\etiquette[5.1.6]{Vbis.subsubsec5.1.6}
用 $E_{(APQ)}$ 表示 $\Hom(\Delta^0,E)$ 的如下子范畴:
其对象为满足 $(APQ)$ 的单纯对象,
其态射为 $\Hom(\Delta^0,E)$ 中的态射 $X\xrightarrow{\alpha}Y$,
使 $\pi(\alpha_n):\pi(X_n)\to\pi(Y_n)$ 对每个 $n$ 都是 $B$ 的同一态射.
于是有典范函子 $\bar{\pi}:E_{(APQ)}\to B$.

\begin{enonce*}{命题 5.1.7}\etiquette[5.1.7]{Vbis.prop5.1.7}
函子\pageoriginale{} $\bar{\pi}$ 满射.
此外, 对 $B$ 的每个对象 $b$ 及 $E_{(APQ),b}$ 的每对对象 $(X,Y)$,
在 $E_{(APQ),b}$ 中存在图
\[
\xymatrix@=1.5cm{
  & Z \ar[ld]\ar[rd]&\\
X &   & Y
}
\]
\end{enonce*}

由 \hyperref[Vbis.subsubsec5.1.4]{5.1.4} 中的假设,
证明立即来自 \hyperref[Vbis.prop5.1.3]{5.1.3}.

\begin{enonce*}{命题 5.1.8}\etiquette[5.1.8]{Vbis.prop5.1.8}
设 $C$ 为范畴, $F$ 为函子 $E_{(APQ)}\to C$.
对 $F$ 引入下述条件:
\begin{enumerate}
\renewcommand{\theenumi}{\roman{enumi}}
\renewcommand{\labelenumi}{(\theenumi)}
\item 对 $B$ 的每个对象 $b$, $F$ 在 $E_{(APQ),b}$ 上的限制
经由 $E_{(APQ),b}$ 的对象所定义的\emph{粗}范畴分解
\footnote{若对每对对象 $(x,y)$, $\Hom(x,y)$ 恰有一个元素, 就称范畴为\emph{粗}范畴.}.

\item 对 $B$ 的每个态射 $t:b'\to b$
及 $t$ 上的每对态射 $(X'\to X,Y'\to Y)$, 下图交换:
\[
\xymatrix@C=3cm@R=2cm{
F(X')         \ar[r]  \ar[d]_{\alpha'}& F(X)\ar[d]^{\alpha}\\
F(Y')          \ar[r] & F(Y)
}
\]
其中 $\alpha$ 和 $\alpha'$ 为条件 \emph{(i)} 所确定的箭头.
\end{enumerate}

则函子 $G\mapsto G\circ\bar{\pi}$
在范畴 $\Hom(B,C)$ 与 $\Hom(E_{(APQ)},C)$ 中
由满足条件 \emph{(i) 和 (ii)} 的函子定义的满子范畴之间诱导等价.
\end{enonce*}

证明的\pageoriginale{}细节留待读者完成.
这里只说明如何证明上述函子本质满射.

设 $F:E_{(APQ)}\to C$ 为满足条件 (i) 和 (ii) 的函子.
对 $B$ 的每个对象 $b$, 选取 $E_{(APQ),b}$ 的对象 $X$ 并令 $\bar{F}(b)=F(X)$;
继而可验 $\bar{F}(b)$ 对 $b$ 函子性地变动.


% END INCLUDED SGA4_English_current_buildable_source_20260906/source/05bis/drafts/SGA4_VBIS_LINES_3151_3452_English.texfrag


% BEGIN INCLUDED SGA4_English_current_buildable_source_20260906/source/05bis/drafts/SGA4_VBIS_LINES_3453_3612_English.texfrag
\subsection{特征零域上有限型概形的奇异上同调}\etiquette[5.2]{Vbis.subsec5.2}

本节用 $k$ 表示特征零域.

\setcounter{subsubsection}{-1}
\subsubsection{}\etiquette[5.2.0]{Vbis.subsubsec5.2.0}
用 $\Sch_f /k$ (分别用 Ann) 表示 $k$ 上\emph{局部有限型}概形的范畴
(分别为复解析空间范畴).
借助典范函子 $\Sch_f /k \to  \Top$ (分别为 $\Ann \to \Top$),
将 $\Sch_f /k$ (分别为 Ann) 的所有单纯对象视为 Top 的单纯对象.
保留 (\hyperref[Vbis.subsubsec4.1.0]{4.1.0}) 中有关上同调下降的资料.

按通常记号, 将 G.A.G.A. 中定义的典范函子 $\Sch_f / \CC \to$ Ann 写为 $S \to S^{\an}$.

\subsubsection{}\etiquette[5.2.1]{Vbis.subsubsec5.2.1}
设 $X$ 为 $\Sch / k$ (分别为 $\Ann$) 的单纯对象:
概形 (分别为解析空间) 态射的 de Rham 复形构造中通常的相容性,
使我们能在 $D^+(\Ysheaf{\Gamma} (\bar{X},\ZZ))$ 中构造复形
$\Omega_{X/k}^{\Zar.}$ (分别为 $\Omega_{X / \CC}^{\bigdot}$),
使对每个整数 $i$, $\sisi{\RR^+}{R} e^{\ast}_i(\Omega_{X/k}^{\Zar.})$
(分别为 $\sisi{\RR^+}{R} e^{\ast}_i$ $(\Omega^{\bigdot}_{X /\CC})$)
为 de Rham 复形 $\Omega_{X/k}^{\Zar.}$ (分别为 $\Omega^{\bigdot}_{X_i/ \CC}$).

设 $X \xrightarrow{u} C^{\Delta}_S$ 为 $\Sch_f / k$ (分别为 Ann) 的增广单纯对象.
定义 $H^{\Zar.}_{DR}(u)$ (分别为 $H^{\bigdot}_{DR}(u)$ 和 $H^{\bigdot}(u, \CC)$)
为复形 $\sisi{\RR^+}{R} \bar{u}_+(\Omega_{X/k}^{\Zar.})$
(分别为 $\sisi{\RR^+}{R} \bar{u}_{\ast}(\Omega_{X/\CC}^{\bigdot})$
及 $\sisi{\RR^+}{R} \bar{u} (\CC)$) 的超上同调群.

\subsubsection{}\etiquette[5.2.2]{Vbis.subsubsec5.2.2}
若\pageoriginale{} $X \xrightarrow{u} C^{\Delta}_S$ 为 $\Sch_f / \CC$ 的半单纯对象,
则有三个典范态射:
\begin{align*}
H^{\Zar.}_{DR} (u) &\xrightarrow{\alpha} H^{\bigdot}_{DR}(u^{\an})\\
H^{\bigdot} (u^{\an}, \CC) &\xrightarrow{\beta}
H^{\bigdot}_{DR}(u^{\an})\\
H^{\bigdot} (S^{\an}, \CC) &\xrightarrow{\gamma}
H^{\bigdot}(u^{\an}, \CC){\quoi}.
\end{align*}

\begin{enonce*}{命题 5.2.3}\etiquette[5.2.3]{Vbis.prop5.2.3}
采用 (\hyperref[Vbis.subsubsec5.2.2]{5.2.2}) 的记号,
假定每个 $n$ 对应的 $X_n$ 光滑: 则 $\alpha$ 和 $\beta$ 为同构.
若进一步 $X^{\an} \xrightarrow{u^{\an}}C^{\Delta}_{S^{\an}}$
为 $S^{\an}$ 关于普遍 $1$-上同调下降拓扑的超覆盖, 则 $\gamma$ 为同构.
\end{enonce*}

$\alpha$ 为同构这一断言来自 [\cite{sga4.exp.Vbis.bib.6}, 定理 $1'$].
Poincar\'e 引理表明 $\beta$ 为同构.
命题最后一部分来自 (\hyperref[Vbis.thm3.3.3]{3.3.3}).

\setcounter{subsubsection}{3}
\subsubsection{}\etiquette[5.2.4]{Vbis.subsubsec5.2.4}
现在将 (\hyperref[Vbis.subsubsec5.1.4]{5.1.4})
应用于如下情形: $B = \Sch_f /k$, 而 $E$ 为 $B$ 上的范畴,
使对每个概形 $S$, 纤维范畴 $E_S$ 为 $\Sch_f /k$ 中 $S$ 上概形的范畴,
$B$ 的箭头上方的箭头按显然方式定义.
若 $X$ 在 $k$ 上光滑, 且 $h$ 为有限多个满足下述条件之一的态射的复合,
就称 $S$ 上的态射 $h : X \to Y$ 满足性质 $P$:
\begin{enumerate}
\renewcommand{\theenumi}{\roman{enumi}}
\renewcommand{\labelenumi}{(\theenumi)}
\item 它是一族平展满射态射的余积;
\item 它是一族固有满射态射的余积.
\end{enumerate}
若结构态射 $X \to S$ 满足性质 $P$, 就称 $E_S$ 的概形 $X$ 满足性质 $Q$.

请读者利用奇点消解验证 (\hyperref[Vbis.subsubsec5.1.4]{5.1.4}) 的全部条件满足.

各群\pageoriginale{} $H^{\Zar.}_{DR}(u)$
定义从范畴 $E_{(APQ)}$ 到分次阿贝尔群范畴的反变函子, 记为 $H^{\Zar.}_{DR}$.

\begin{enonce*}{命题 5.2.5}\etiquette[5.2.5]{Vbis.prop5.2.5}
函子 $H^{\Zar.}_{DR}$ 满足 (\hyperref[Vbis.prop5.1.8]{5.1.8}) 的条件.
\end{enonce*}

由 Lefschetz 原理, 归结为 $k = \CC$ 的情形.
此时对每个态射 $S'\xrightarrow{t} S$ 及 $t$ 上的每对态射 $(u'\to u$, $v'\to v)$,
有交换图:
\[
\xymatrix@C=3.5cm@R=2cm{
H^{\Zar.}_{DR}(u') \ar[r] \ar[d]_{\gamma^{-1}.\beta^{-1}.\alpha}&  H_{DR}^{\Zar.}(u)\ar[d]^{\gamma^{-1}.\beta^{-1}.\alpha}\\
H^{\bigdot}(S'^{\an},\CC)
\ar[r]^{H^{\bigdot}(t^{\an},\CC)}
& H^{\bigdot}(S^{\an},
\CC) \\
H^{\Zar.}_{DR}(v') \ar[r] \ar[u]^{\gamma^{-1}.\beta^{-1}.\alpha}& H_{DR}^{\Zar.}(v) \ar[u]_{\gamma^{-1}.\beta^{-1}.\alpha}
}
\]
证明完毕.

\begin{enonce*}{定义 5.2.6}\etiquette[5.2.6]{Vbis.def5.2.6}
由 (\hyperref[Vbis.prop5.1.8]{5.1.8}) 从 $H^{\Zar.}_{DR}$ 定义的反变函子
记为 $H^{\bigdot} (\ast, k)$.
对 $k$ 上局部有限型的每个概形 $S$,
分次阿贝尔群 $H^{\bigdot}(S, k)$ 称为 \emph{$S$ 的奇异上同调}.
\end{enonce*}

\setcounter{subsubsection}{6}
\subsubsection{}\etiquette[5.2.7]{Vbis.subsubsec5.2.7}
以一个例子结束本节, 说明本节引入的一般原理.
先回顾, 若 $S$ 为复解析空间, 则 $S$ 的 \emph{de Rham 上同调},
记为 $\HH^{\bigdot}_{DR}(S)$,
按定义为 de Rham 复形 $\Omega^{\bigdot}_{S / \CC}$ 的超上同调.

\begin{enonce*}{命题 5.2.8}\etiquette[5.2.8]{Vbis.prop5.2.8}
\emph{(Blum--Herrera)}.\pageoriginale{} 设 $S$ 为 $\CC$ 上局部有限型的概形.
典范箭头
$$
H^{\bigdot} (S^{\an}, \CC) \longrightarrow
\HH^{\bigdot}_{DR}(S^{\an})
$$
容许一个对 $S$ 函子性的回缩 (在分次阿贝尔群范畴中).
\end{enonce*}

由 (\hyperref[Vbis.prop5.1.8]{5.1.8}),
只须考虑由前述态射与 $\bar{\pi}$ 复合得到的态射
$H^{\bigdot}(+, \CC) \circ \bar{\pi} \to\HH^{\bigdot}_{DR} \circ \bar{\pi}$.

对每个增广 $X \xrightarrow{u} C^{\Delta}_S$, 利用典范态射
$\LL \bar{u}^{+}(\Omega^{\bigdot}_{S/C}) \to\Omega^{\bigdot}_{X/C}$,
定义函子性态射 $\HH^{\bigdot}_{DR} \circ \bar{\pi} \to \HH^{\bigdot}_{DR}$,
使下图交换:
\[
\xymatrix@=2cm{
H^{\bigdot}_{DR}  & H^{\bigdot}_{DR} \circ \bar{\pi} \ar[l]\\
H^{\bigdot}(\ast,\CC)_{\circ}\bar{\pi}
\ar[u]_{\text{\rotatebox{90}{$\backsim$}}}^{\beta \circ \gamma}
\ar[ru]&
}
\]
由此得到所需的回缩.

\begin{enonce*}[remark]{注 5.2.9}\etiquette[5.2.9]{Vbis.rem5.2.9}
一旦在解析情形中也能作奇点消解,
命题 (\hyperref[Vbis.prop5.2.8]{5.2.8}) 对每个解析空间仍成立.
\end{enonce*}


% END INCLUDED SGA4_English_current_buildable_source_20260906/source/05bis/drafts/SGA4_VBIS_LINES_3453_3612_English.texfrag


% BEGIN INCLUDED SGA4_English_current_buildable_source_20260906/source/05bis/drafts/SGA4_VBIS_LINES_3613_3730_English.texfrag
\subsection{混合 Hodge 理论}\etiquette[5.3]{Vbis.subsec5.3}

这里给出一个在 \cite{sga4.exp.Vbis.bib.2} 中起关键作用的技术结果.

$k$ 为特征零域.

\subsubsection{}\etiquette[5.3.1]{Vbis.subsubsec5.3.1}
设 $\Sch_{\fs} /k$ 为 $k$ 上的有限型分离概形的范畴;
用 $E \xrightarrow{\pi} \Sch_{\fs}/k$ 表示如下定义的 $\Sch_{\fs}/k$ 上的范畴:

\noindent \sisi{$+$}{} 若\pageoriginale{} $S$ 为 $\Sch_{\fs} /k$ 的对象,
则 $E_S$ 的对象为三元组 $(\bar{X}, X, i)$,
其中 $\bar{X}$ 为 $k$ 上的约化射影概形,
$X$ 为 $S$ 上的固有概形, $i : X \to \bar{X}$ 为 $k$ 上的开浸入,
使 $i(X)$ 在 $\bar{X}$ 中稠密.

\noindent \sisi{$+$}{} 若 $t : S' \to S$ 为概形态射,
$(\bar{X}, X, i)$ 为 $S$ 上的对象,
$(\bar{X}', X', i')$ 为 $S'$ 上的对象,
则 $t$ 上从 $(\bar{X}, X, i)$ 到 $(\bar{X}', X', i')$ 的态射为二元组 $(\bar{h}, h)$,
其中 $\bar{h} : \bar{X}'\to \bar{X}$ 为 $k$ 上的态射,
$h : X'\to X$ 为 $S$ 上的态射, 使下图交换.
\[
\xymatrix@C=3.5cm@R=2cm{
X' \ar[r]^{i'} \ar[d]_{h}& \bar{X}'\ar[d]^{\bar{h}}\\
X  \ar[r]^{i}& \bar{X}.
}
\]
附记: 固有性假设蕴含 $\bar{h}^{-1}(i(X)) = i(X')$.

\begin{enonce*}{引理 5.3.2}\etiquette[5.3.2]{Vbis.lem5.3.2}
对 $\Sch_{\fs} / k$ 的每个对象 $S$, 范畴 $E_S$ 非空,
具有非空投射极限, 且具有不交并普遍的有限余积.
\end{enonce*}

范畴 $E_S$ 非空这一事实来自 Chow 引理.

这里只描述 $E_S$ 的两个对象 $(\bar{X}, X, i)$ 与 $(\bar{X}', X', i')$ 的直积构造.

典范箭头 $j : \fprod{X}{X'}{S} \to\fprod{\bar{X}}{\bar{X}'}{k}$ 为浸入.
令 $T$ 为 $\fprod{\bar{X}}{\bar{X}'}{k}$ 的约化闭子概形,
其底层空间为 $j$ 的像的闭包.
则 $(T, {\fprod{X}{X'_{\red}}{S}}, j_{\red})$ 具有所需的泛性质.

\setcounter{subsubsection}{2}
\subsubsection{}\etiquette[5.3.3]{Vbis.subsubsec5.3.3}
设\pageoriginale{} $S$ 为 $\Sch_{\fs} /k$ 的对象.
若 $\bar{X}$ 在 $k$ 上光滑且 $i(X)$ 为正规交叉除子的补集,
就称 $E_S$ 的对象 $(\bar{X}, X, i)$ 满足性质 $Q$.
若 $\bar{h}$ 满射且 $(\bar{X}', X', i')$ 满足性质 $Q$,
就称态射 $(\bar{h}, h) : (\bar{X}', X', i') \to (\bar{X}, X, i)$ 满足性质 $P$.

\begin{enonce*}{命题 5.3.4}\etiquette[5.3.4]{Vbis.prop5.3.4}
采用 (\hyperref[Vbis.subsubsec5.3.3]{5.3.3}) 的记号,
(\hyperref[Vbis.subsubsec5.1.4]{5.1.4}) 的全部条件均满足.
\end{enonce*}

利用奇点消解容易验证条件 a)、b)、d) 和 e).
此外, c) 来自 f), 故只须验证 f).

须证明每个交换图
\[
\xymatrix@C=3cm@R=1cm{
              &Z         \ar[r]^{i_Z} \ar[d]_{h}& \bar{Z}\ar[d]^{\bar{h}}\\
              &X \ar[r]^{i_X} \ar[d]& \bar{X}\ar[d]\\
 S'\ar[r]^{t} &S  \ar[r]            & k
}
\]
其中 $(\bar{X}, X, i_X)$ 与 $(\bar{Z}, Z, i_Z)$ 为 $E_S$ 的对象,
且 $\bar{h}$ 满射, 都可补成交换图
\[
\xymatrix@C=1.7cm@R=.45cm{
Z' \ar[rr] \ar[rd]_{i_{Z'}} \ar[dd]_{h'}&    & Z \ar[rd]^{i_Z} \ar@{->} '[d][dd]_{h} &\\
         & \bar{Z}' \ar[rr] \ar@{->} '[d][dd]&    & \bar{Z}\ar[dd]^{\bar{h}}\\
X'  \ar[rr]^(.65){h'} \ar[dd]\ar[rd]_{i_{X'}}   &    & X \ar[rd]_{i_X} \ar@{->} '[d][dd]&\\
         & \bar{X}'  \ar@{->} '[d][dd] \ar[rr]&    & \bar{X} \ar[dd]\\
S'\ar[rr]^(.65){t}\ar[rd]        &    & S \ar[rd] &\\
         & k  \ar@{=}[rr]&    & k
}
\]
其中\pageoriginale{} $(\bar{X}', X', i_{X'})$ 与
$(\bar{Z}', Z', i_{Z'})$ 为 $E_{S'}$ 的对象, 且 $\bar{h'}$ 满射.

考虑图
\[
\xymatrix@C=2cm@R=1.5cm@M=.2cm{
T \ar[r]^{\varphi} \ar@{_{(}->}[d]_{i_T}& S'\ar[r]^{t} & S\ar[d]\\
\bar{T} \ar[rr]_{\psi}&  & k\\
}
\]
其中 $\varphi$ 固有, $\psi$ 射影. 经基变换得到图
\[
\xymatrix@C=1cm@R=.6cm{
\fprod{Z}{T_{\red}}{S} \ar[rrr] \ar[rd]_{p} \ar[dd]&&    & Z \ar[rd] \ar@{->} '[d][dd] &\\
         & \fprod{\bar{Z}}{\bar{T}_{\red}}{k} \ar[rrr] \ar@{->} '[d][dd]& &   & \bar{Z}\ar[dd]\\
\fprod{X}{T_{\red}}{S}  \ar[rrr] \ar[dd]\ar[rd]_{q} &  &    & X \ar[rd] \ar@{->} '[d][dd]&\\
         &\fprod{\bar{X}}{\bar{T}_{\red}}{k}  \ar@{->} '[d][dd] \ar[rrr]& &   & \bar{X} \ar[dd]\\
T \ar[rd]\ar[r]& S'\ar[rr]        &    & S \ar[rd] &\\
         & \bar{T}  \ar[rrr]&  &  & k
}
\]
其中\pageoriginale{} $p$ 和 $q$ 为浸入.
于是只须将 $\fprod{\bar{Z}}{\bar{T}_{\red}}{k}$
和 $\fprod{\bar{X}}{\bar{T}_{\red}}{k}$ 分别换为 $p$ 与 $q$ 的闭像,
便得到所需的图. 这就完成了 (\hyperref[Vbis.prop5.3.4]{5.3.4}) 的证明.


% END INCLUDED SGA4_English_current_buildable_source_20260906/source/05bis/drafts/SGA4_VBIS_LINES_3613_3730_English.texfrag


% BEGIN INCLUDED SGA4_English_current_buildable_source_20260906/source/05bis/drafts/SGA4_VBIS_LINES_3731_3755_English.texfrag
\begin{thebibliography}{99}\pageoriginale
\bibitem{sga4.exp.Vbis.bib.1} T. Blum -- M. Herrera : De Rham Cohomology of an
  analytic Space. Inv. Math. vol 7. Fas 4-1969. p.275-296.
\bibitem{sga4.exp.Vbis.bib.2} P. Deligne : Théorie de Hodge III (I.H.E.S).
\bibitem{sga4.exp.Vbis.bib.3} P. Gabriel -- M. Zisman : Homotopy theory and calculus of
  fractions Ergebnisse der Mathematik -- Band 35 -- Springer 1967.
\bibitem{sga4.exp.Vbis.bib.4} J. Giraud : Méthode de la descente : Mémoires de la
  S.M.F. 2. 1964.
\bibitem{sga4.exp.Vbis.bib.5} R. Godement : Théorie des faisceaux.
                                Publications de l'Institut de
                                Mathématique de l'Université de
                                Strasbourg. Hermann 1964.
\bibitem{sga4.exp.Vbis.bib.6} A. Grothendieck : On the De Rham Cohomology of Algebraic
  Varieties, Publications de 1' I.H.E.S ${\rm n^{\circ}}$ 29.
\bibitem{sga4.exp.Vbis.bib.7} R. Hartshorne : Residues and Duality. Springer
  Lecture Notes ${\rm n^{\circ}}$ 20.
\bibitem{sga4.exp.Vbis.bib.8} G. Segal : Classifying spaces and Spectral sequences
  Publications de 1'I.H.E.S. ${\rm n^{\circ}}$ 34.
\bibitem{sga4.exp.Vbis.bib.9} J.L.Verdier : Catégories dérivées - Etat 0 -
  (I.H.E.S)
\end{thebibliography}

\ifx\danslelivre\undefined
\end{document}
\fi

% END INCLUDED SGA4_English_current_buildable_source_20260906/source/05bis/drafts/SGA4_VBIS_LINES_3731_3755_English.texfrag


